% Öffentliche Aussagen bleiben in Band 37; ihre Beweise stehen hier.
\section*{Multiplikationstafeln und Isomorphie endlicher Halbgruppen}
\hypertarget{table-isomorphisms.proofs}{}
Diese Ergänzung enthält die Beweistabellen zu den neun Sätzen des
Tafelabschnitts in Band~37, die Wohldefiniertheit der nummerierten Tafel
und die ergänzenden Nachweise für Automorphismen, Nummerierungswechsel
und kanonische Vertreter. Die \TableIsomorphismsReadingLink{} erklärt
Zusammenhang und Beispiele; zur Einordnung führt die \TableIsomorphismsMainLink{}.
Die ursprünglichen sieben Hilfssätze behalten ihre Nummern und Kennungen.
Neue Hilfssätze tragen die Bandkennung \(37\mathrm T\).
Alle Aussagen gelten auch für den leeren Träger. Eine Nummerierung der
leeren Menge ist die leere Bijektion; ihre Tafel und ihr Tafeltupel sind leer.

\section{Nummerierte Tafeln}
\subsection{Wohldefiniertheit}
\hypertarget{table-isomorphisms.proof.FiniteSemigroupMultiplicationTableDef}{}
\noindent Nachweis zur \FormulaRefAuto{FiniteSemigroupMultiplicationTableDef}[def].
\noindent\textbf{Zu zeigen: Werte der nummerierten Tafel.}
Die folgende Aussage ist der Typnachweis für die eben wiedergegebene
Funktionsvorschrift; alle Voraussetzungen stehen vor dem Folgerungszeichen.
\[
\begin{aligned}
 &\FiniteSemigroupAlg{A}{\star}\dsep n\in\mathbb N\dsep
 e\colon\NatSeg n\bij A\vdash{}\\
 &\forall(i,j)\in\NatSeg n\times\NatSeg n\,
 e^{-1}\!\bigl(e(i)\star e(j)\bigr)\in\NatSeg n.
\end{aligned}
\]
\begin{tabproofwide}
  \proofstepwidestar[1]{\FiniteSemigroupAlg{A}{\star}}{\rA}
  \proofstepwidestar[2]{n\in\mathbb N}{\rA}
  \proofstepwidestar[3]{e\colon\NatSeg{n}\bij A}{\rA}
  \proofstepwidestar[4]{%
    (i,j)\in\NatSeg{n}\times\NatSeg{n}}{\rA}
  \proofstepwidestar[1]{\SemigroupAlg{A}{\star}}{%
    \FormulaRefAuto{FiniteSemigroupBaseAxiom}[ax]{1}}
  \proofstepwidestar[4]{i\in\NatSeg{n}}{%
    \FormulaRefAuto{(a,b)\in A\times B\vdash a\in A}{4}}
  \proofstepwidestar[4]{j\in\NatSeg{n}}{%
    \FormulaRefAuto{(a,b)\in A\times B\vdash b\in B}{4}}
  \proofstepwidestar[3,4]{e(i)\in A}{%
    \FormulaRefAuto{F\colon A\bij B\dsep x\in A\vdash F(x)\in B}%
      [thm]{3,6}}
  \proofstepwidestar[3,4]{e(j)\in A}{%
    \FormulaRefAuto{F\colon A\bij B\dsep x\in A\vdash F(x)\in B}%
      [thm]{3,7}}
  \proofstepwidestar[1,3,4]{e(i)\star e(j)\in A}{%
    \FormulaRefAuto{SemigroupClosureAxiom}[ax]{5,8,9}}
  \proofstepwidestar[3]{e^{-1}\colon A\to\NatSeg{n}}{%
      \FormulaRefAuto{InverseFunctionType}%
        [thm]{3}}
  \proofstepwidestar[1,3,4]{%
    e^{-1}\!\bigl(e(i)\star e(j)\bigr)\in\NatSeg{n}}{%
    \FormulaRefAuto{F\colon A\to B\dsep x\in A\vdash F(x)\in B}[thm]{%
      11,10}}
  \proofstepwidestar[1,3]{%
    \forall (i,j)\in\NatSeg{n}\times\NatSeg{n}\,
      e^{-1}\!\bigl(e(i)\star e(j)\bigr)\in\NatSeg{n}}{%
    \rUI{\rRI{4,12}}}
\end{tabproofwide}


\subsection{Transport der Halbgruppenstruktur}
\hypertarget{table-isomorphisms.proof.FiniteSemigroupMultiplicationTableTransport}{}
\noindent Beweis zu \FormulaRefAuto{FiniteSemigroupMultiplicationTableTransport}[thm].
\begingroup
\renewcommand{\theformulaThm}{37.4.1.1}
% Beweis zu FiniteSemigroupMultiplicationTableTransport.
In den folgenden Tabellen stehen \(S=\NatSeg n\) und
\(m=m_{\star,e}\) nur als Abkürzungen.
\begin{tabproofsplit}
\proofpartwideindDelta[Die Transportgleichung]{%
  \begin{aligned}[t]
  &\FiniteSemigroupAlg{A}{\star}\dsep n\in\mathbb N\dsep
    e\colon S\bij A\dsep i,j\in S\vdash{}\\[-2pt]
  &e(m(i,j))=e(i)\star e(j)
  \end{aligned}
}{%
  \DeltaRow{Mengen}{A\dsep n\dsep i\dsep j}
  \DeltaRow{Funktionen}{e\dsep\star\dsep m}
}[][B37TableTransportEvaluation]
  \proofstepwidestar[1]{\FiniteSemigroupAlg{A}{\star}}{\rA}
  \proofstepwidestar[2]{n\in\mathbb N}{\rA}
  \proofstepwidestar[3]{e\colon S\bij A}{\rA}
  \proofstepwidestar[4]{i\in S}{\rA}
  \proofstepwidestar[5]{j\in S}{\rA}
  \proofstepwidestar[1]{\SemigroupAlg{A}{\star}}{\FormulaRefAuto{FiniteSemigroupBaseAxiom}[ax]{1}}
  \proofstepwidestar[3,4]{e(i)\in A}{\FormulaRefAuto{F\colon A\bij B\dsep x\in A\vdash F(x)\in B}[thm]{3,4}}
  \proofstepwidestar[3,5]{e(j)\in A}{\FormulaRefAuto{F\colon A\bij B\dsep x\in A\vdash F(x)\in B}[thm]{3,5}}
  \proofstepwidestar[1,3,4,5]{e(i)\star e(j)\in A}{\FormulaRefAuto{SemigroupClosureAxiom}[ax]{6,7,8}}
  \proofstepwidestarbreak[1,2,3,4,5]{m(i,j)=e^{-1}(e(i)\star e(j))}{\FormulaRefAuto{FiniteSemigroupMultiplicationTableDef}[def]{1,2,3,4,5}}
  \proofstepwidestarbreak[1,3,4,5]{e(e^{-1}(e(i)\star e(j)))=e(i)\star e(j)}{\FormulaRefAuto{InverseFunctionRightInverse}[thm]{3,9}}
  \proofstepwidestar[1,2,3,4,5]{e(m(i,j))=e(i)\star e(j)}{\rIE{10,11}}
\closeproofpartwideindDelta
\proofpartwideindDelta[Assoziativität der nummerierten Operation]{%
  \begin{aligned}[t]
  &\FiniteSemigroupAlg{A}{\star}\dsep n\in\mathbb N\dsep
    e\colon S\bij A\vdash{}\\[-2pt]
  &\SemigroupAlg{S}{m}
  \end{aligned}
}{%
  \DeltaRow{Mengen}{A\dsep n\dsep i\dsep j\dsep k}
  \DeltaRow{Funktionen}{e\dsep\star\dsep m}
}[][B37TableTransportSemigroup]
  \proofstepwidestar[1]{\FiniteSemigroupAlg{A}{\star}}{\rA}
  \proofstepwidestar[2]{n\in\mathbb N}{\rA}
  \proofstepwidestar[3]{e\colon S\bij A}{\rA}
  \proofstepwidestar[1,2,3]{m\colon S\times S\to S}{\FormulaRefAuto{FiniteSemigroupMultiplicationTableDef}[def]{1,2,3}}
  \proofstepwidestar[1]{\SemigroupAlg{A}{\star}}{\FormulaRefAuto{FiniteSemigroupBaseAxiom}[ax]{1}}
  \proofstepwidestar[3]{e\colon S\inj A}{\FormulaRefAuto{F\colon A\bij B\vdash F\colon A\inj B}[thm]{3}}
  \proofstepwidestar[7]{i\in S}{\rA}
  \proofstepwidestar[8]{j\in S}{\rA}
  \proofstepwidestar[9]{k\in S}{\rA}
  \proofstepwidestar[1,2,3,7,8]{m(i,j)\in S}{\FormulaRefAuto{F\colon A\to B\dsep x\in A\vdash F(x)\in B}[thm]{4,\FormulaRefAuto{a\in A\dsep b\in B\vdash(a,b)\in A\times B}[thm]{7,8}}}
  \proofstepwidestar[1,2,3,8,9]{m(j,k)\in S}{\FormulaRefAuto{F\colon A\to B\dsep x\in A\vdash F(x)\in B}[thm]{4,\FormulaRefAuto{a\in A\dsep b\in B\vdash(a,b)\in A\times B}[thm]{8,9}}}
  \proofstepwidestarbreak[1,2,3,7,8,9]{e(m(m(i,j),k))=e(m(i,j))\star e(k)}{\FormulaRefAuto{B37TableTransportEvaluation}[thm]{1,2,3,10,9}}
  \proofstepwidestarbreak[1,2,3,7,8,9]{e(m(m(i,j),k))=(e(i)\star e(j))\star e(k)}{\rIE{\FormulaRefAuto{B37TableTransportEvaluation}[thm]{1,2,3,7,8},12}}
  \proofstepwidestarbreak[1,2,3,7,8,9]{e(m(i,m(j,k)))=e(i)\star e(m(j,k))}{\FormulaRefAuto{B37TableTransportEvaluation}[thm]{1,2,3,7,11}}
  \proofstepwidestarbreak[1,2,3,7,8,9]{e(m(i,m(j,k)))=e(i)\star(e(j)\star e(k))}{\rIE{\FormulaRefAuto{B37TableTransportEvaluation}[thm]{1,2,3,8,9},14}}
  \proofstepwidestarbreak[1,3,7,8,9]{(e(i)\star e(j))\star e(k)=e(i)\star(e(j)\star e(k))}{\FormulaRefAuto{SemigroupAssociativityAxiom}[ax]{5,\FormulaRefAuto{F\colon A\bij B\dsep x\in A\vdash F(x)\in B}[thm]{3,7},\FormulaRefAuto{F\colon A\bij B\dsep x\in A\vdash F(x)\in B}[thm]{3,8},\FormulaRefAuto{F\colon A\bij B\dsep x\in A\vdash F(x)\in B}[thm]{3,9}}}
  \proofstepwidestarbreak[1,2,3,7,8,9]{e(m(m(i,j),k))=e(m(i,m(j,k)))}{\rIE{13,15,16}}
  \proofstepwidestarbreak[1,2,3,7,8,9]{m(m(i,j),k)=m(i,m(j,k))}{\FormulaRefAuto{F\colon A\inj B\dsep x\in A\dsep y\in A\dsep F(x)=F(y)\vdash x=y}[ax]{6,\FormulaRefAuto{FiniteSemigroupMultiplicationTableDef}[def]{1,2,3,10,9},\FormulaRefAuto{FiniteSemigroupMultiplicationTableDef}[def]{1,2,3,7,11},17}}
  \proofstepwidestarbreak[1,2,3]{\forall i,j,k\in S\,(m(m(i,j),k)=m(i,m(j,k)))}{\rUI{\rRI{7,\rUI{\rRI{8,\rUI{\rRI{9,18}}}}}}}
  \proofstepwidestar[1,2,3]{\SemigroupAlg{S}{m}}{\FormulaRefAuto{SemigroupDef}[def]{4,19}}
\closeproofpartwideindDelta
\proofpartwide{Endlichkeit und Isomorphismus}
  \proofstepwidestar[1]{\FiniteSemigroupAlg{A}{\star}}{\rA}
  \proofstepwidestar[2]{n\in\mathbb N}{\rA}
  \proofstepwidestar[3]{e\colon S\bij A}{\rA}
  \proofstepwidestar[1]{\SemigroupAlg{A}{\star}}{\FormulaRefAuto{FiniteSemigroupBaseAxiom}[ax]{1}}
  \proofstepwidestar[1,2,3]{\SemigroupAlg{S}{m}}{\FormulaRefAuto{B37TableTransportSemigroup}[thm]{1,2,3}}
  \proofstepwidestar[2]{\Finite S}{\FormulaRefAuto{PeanoNatSegFinite}[thm]{2}}
  \proofstepwidestar[1,2,3]{\FiniteSemigroupAlg{S}{m}}{\FormulaRefAuto{FiniteSemigroupDef}[def]{5,6}}
  \proofstepwidestar[3]{e\colon S\to A}{\FormulaRefAuto{Bijektive Funktion}[def]{3}}
  \proofstepwidestar[9]{i\in S}{\rA}
  \proofstepwidestar[10]{j\in S}{\rA}
  \proofstepwidestar[1,2,3,9,10]{e(m(i,j))=e(i)\star e(j)}{\FormulaRefAuto{B37TableTransportEvaluation}[thm]{1,2,3,9,10}}
  \proofstepwidestarbreak[1,2,3]{\forall i,j\in S\,(e(m(i,j))=e(i)\star e(j))}{\rUI{\rRI{9,\rUI{\rRI{10,11}}}}}
  \proofstepwidestar[1,2,3]{\SemigroupHom{e}{S,m}{A,\star}}{\FormulaRefAuto{SemigroupHomDef}[def]{5,4,8,12}}
  \proofstepwidestar[1,2,3]{\FiniteSemigroupHom{e}{S,m}{A,\star}}{\FormulaRefAuto{FiniteSemigroupHomDef}[def]{7,1,13}}
  \proofstepwidestar[1,2,3]{\FiniteSemigroupIso{e}{S,m}{A,\star}}{\FormulaRefAuto{FiniteSemigroupIsoDef}[def]{14,3}}
  \proofstepwidestarbreak[1,2,3]{\FiniteSemigroupAlg{S}{m}\land\FiniteSemigroupIso{e}{S,m}{A,\star}}{\rAI{7,15}}
\closeproofpartwide
\end{tabproofsplit}

\endgroup

\subsection{Das Umbenennungskriterium}
\hypertarget{table-isomorphisms.proof.FiniteSemigroupTableRelabelCriterion}{}
\noindent Beweis zu \FormulaRefAuto{FiniteSemigroupTableRelabelCriterion}[thm].
% Beweis zu FiniteSemigroupTableRelabelCriterion.
Wir kürzen \(S=\NatSeg n\) und
\(Q=\forall i,j\in S\,(\pi(i\star j)=\pi(i)\diamond\pi(j))\) ab.
\begin{tabproofwide}
  \proofstepwidestar[1]{n\in\mathbb N}{\rA}
  \proofstepwidestar[2]{\FiniteSemigroupAlg{S}{\star}}{\rA}
  \proofstepwidestar[3]{\FiniteSemigroupAlg{S}{\diamond}}{\rA}
  \proofstepwidestar[4]{\pi\colon S\bij S}{\rA}
  \proofstepwidestar[5]{\FiniteSemigroupIso{\pi}{S,\star}{S,\diamond}}{\rA}
  \proofstepwidestar[5]{\FiniteSemigroupHom{\pi}{S,\star}{S,\diamond}}{\FormulaRefAuto{FiniteSemigroupIsoHomAxiom}[ax]{5}}
  \proofstepwidestar[5]{\SemigroupHom{\pi}{S,\star}{S,\diamond}}{\FormulaRefAuto{FiniteSemigroupHomBaseHomAxiom}[ax]{6}}
  \proofstepwidestar[8]{i\in S}{\rA}
  \proofstepwidestar[9]{j\in S}{\rA}
  \proofstepwidestar[5,8,9]{\pi(i\star j)=\pi(i)\diamond\pi(j)}{\FormulaRefAuto{SemigroupHomOperationAxiom}[ax]{7,8,9}}
  \proofstepwidestar[5]{Q}{\rUI{\rRI{8,\rUI{\rRI{9,10}}}}}
  \proofstepwidestar[]{\FiniteSemigroupIso{\pi}{S,\star}{S,\diamond}\to Q}{\rRI{5,11}}
  \proofstepwidestar[13]{Q}{\rA}
  \proofstepwidestar[2]{\SemigroupAlg{S}{\star}}{\FormulaRefAuto{FiniteSemigroupBaseAxiom}[ax]{2}}
  \proofstepwidestar[3]{\SemigroupAlg{S}{\diamond}}{\FormulaRefAuto{FiniteSemigroupBaseAxiom}[ax]{3}}
  \proofstepwidestar[4]{\pi\colon S\to S}{\FormulaRefAuto{Bijektive Funktion}[def]{4}}
  \proofstepwidestar[2,3,4,13]{\SemigroupHom{\pi}{S,\star}{S,\diamond}}{\FormulaRefAuto{SemigroupHomDef}[def]{14,15,16,13}}
  \proofstepwidestar[2,3,4,13]{\FiniteSemigroupHom{\pi}{S,\star}{S,\diamond}}{\FormulaRefAuto{FiniteSemigroupHomDef}[def]{2,3,17}}
  \proofstepwidestar[2,3,4,13]{\FiniteSemigroupIso{\pi}{S,\star}{S,\diamond}}{\FormulaRefAuto{FiniteSemigroupIsoDef}[def]{18,4}}
  \proofstepwidestar[2,3,4]{Q\to\FiniteSemigroupIso{\pi}{S,\star}{S,\diamond}}{\rRI{13,19}}
  \proofstepwidestarbreak[2,3,4]{\FiniteSemigroupIso{\pi}{S,\star}{S,\diamond}\leftrightarrow Q}{\rLRI{12,20}}
\end{tabproofwide}


\section{Umbenennungen und Automorphismen}
\subsection{Assoziative Operationen}
\hypertarget{table-isomorphisms.proof.FiniteSemigroupAssociativeOperationSemigroup}{}
\noindent Beweis zu \FormulaRefAuto{FiniteSemigroupAssociativeOperationSemigroup}[thm].
\begin{tabproofwide}
  \proofstepwidestar[1]{m\in\operatorname{Assoc}(S)}{\rA}
  \proofstepwidestar[1]{\SemigroupAlg{S}{m}}{%
    \FormulaRefAuto{FiniteSemigroupAssociativeOperationsDef}[def]{1}}
\end{tabproofwide}
\subsection{Selbstbijektionen}
\hypertarget{table-isomorphisms.proof.FiniteSemigroupSymmetricSetBijection}{}
\TableIsomorphismsStatement{FiniteSemigroupSymmetricSetBijection}
\begin{tabproofwide}
  \proofstepwidestar[1]{\pi\in\operatorname{Sym}(S)}{\rA}
  \proofstepwidestar[1]{\pi\colon S\bij S}{%
    \FormulaRefAuto{FiniteSemigroupSymmetricSetDef}[def]{1}}
\end{tabproofwide}


\subsection{Das Wirkungsgesetz}
\hypertarget{table-isomorphisms.proof.FiniteSemigroupRelabellingAction}{}
\noindent Beweis zu \FormulaRefAuto{FiniteSemigroupRelabellingAction}[thm].
\begingroup
\renewcommand{\theformulaThm}{37.4.1.5}
% Beweis zu FiniteSemigroupRelabellingAction.
\begin{tabproofsplit}
\proofpartwideindDelta[Die umbenannte Operation ist assoziativ]{%
  m\in\operatorname{Assoc}(S)\dsep\pi\in\operatorname{Sym}(S)
  \vdash m^\pi\in\operatorname{Assoc}(S)
}{%
  \DeltaRow{Mengen}{S\dsep u\dsep v\dsep w}
  \DeltaRow{Funktionen}{m\dsep\pi\dsep m^\pi}
}[][B37RelabelledOperationAssociative]
  \proofstepwidestar[1]{m\in\operatorname{Assoc}(S)}{\rA}
  \proofstepwidestar[2]{\pi\in\operatorname{Sym}(S)}{\rA}
  \proofstepwidestar[1]{\SemigroupAlg{S}{m}}{\FormulaRefAuto{FiniteSemigroupAssociativeOperationSemigroup}[thm]{1}}
  \proofstepwidestar[2]{\pi\colon S\bij S}{\FormulaRefAuto{FiniteSemigroupSymmetricSetBijection}[thm]{2}}
  \proofstepwidestar[1,2]{m^\pi\colon S\times S\to S}{\FormulaRefAuto{FiniteSemigroupRelabelledOperationDef}[def]{1,2}}
  \proofstepwidestar[6]{u\in S}{\rA}
  \proofstepwidestar[7]{v\in S}{\rA}
  \proofstepwidestar[8]{w\in S}{\rA}
  \proofstepwidestar[2,6]{\pi^{-1}(u)\in S}{\FormulaRefAuto{InverseFunctionValueType}[thm]{4,6}}
  \proofstepwidestar[2,7]{\pi^{-1}(v)\in S}{\FormulaRefAuto{InverseFunctionValueType}[thm]{4,7}}
  \proofstepwidestar[2,8]{\pi^{-1}(w)\in S}{\FormulaRefAuto{InverseFunctionValueType}[thm]{4,8}}
  \proofstepwidestarbreak[1,2,6,7]{m(\pi^{-1}(u),\pi^{-1}(v))\in S}{\FormulaRefAuto{SemigroupClosureAxiom}[ax]{3,9,10}}
  \proofstepwidestarbreak[1,2,7,8]{m(\pi^{-1}(v),\pi^{-1}(w))\in S}{\FormulaRefAuto{SemigroupClosureAxiom}[ax]{3,10,11}}
  \proofstepwidestarbreak[1,2,6,7]{m^\pi(u,v)=\pi(m(\pi^{-1}(u),\pi^{-1}(v)))}{\FormulaRefAuto{FiniteSemigroupRelabelledOperationDef}[def]{1,2,6,7}}
  \proofstepwidestarbreak[1,2,7,8]{m^\pi(v,w)=\pi(m(\pi^{-1}(v),\pi^{-1}(w)))}{\FormulaRefAuto{FiniteSemigroupRelabelledOperationDef}[def]{1,2,7,8}}
  \proofstepwidestarbreak[1,2,6,7]{\pi^{-1}(m^\pi(u,v))=m(\pi^{-1}(u),\pi^{-1}(v))}{\rIE{14,\FormulaRefAuto{InverseFunctionLeftInverse}[thm]{4,12}}}
  \proofstepwidestarbreak[1,2,7,8]{\pi^{-1}(m^\pi(v,w))=m(\pi^{-1}(v),\pi^{-1}(w))}{\rIE{15,\FormulaRefAuto{InverseFunctionLeftInverse}[thm]{4,13}}}
  \proofstepwidestar[1,2,6,7]{m^\pi(u,v)\in S}{\FormulaRefAuto{F\colon A\to B\dsep x\in A\vdash F(x)\in B}[thm]{5,\FormulaRefAuto{a\in A\dsep b\in B\vdash(a,b)\in A\times B}[thm]{6,7}}}
  \proofstepwidestar[1,2,7,8]{m^\pi(v,w)\in S}{\FormulaRefAuto{F\colon A\to B\dsep x\in A\vdash F(x)\in B}[thm]{5,\FormulaRefAuto{a\in A\dsep b\in B\vdash(a,b)\in A\times B}[thm]{7,8}}}
  \proofstepwidestarbreak[1,2,6,7,8]{\begin{aligned}[t]
    &m^\pi(m^\pi(u,v),w)\\[-2pt]
    &\quad=\pi(m(m(\pi^{-1}(u),\pi^{-1}(v)),\pi^{-1}(w)))
    \end{aligned}}{\FormulaRefAuto{FiniteSemigroupRelabelledOperationDef}[def]{1,2,18,8},\rIE{16}}
  \proofstepwidestarbreak[1,2,6,7,8]{\begin{aligned}[t]
    &m(m(\pi^{-1}(u),\pi^{-1}(v)),\pi^{-1}(w))\\[-2pt]
    &\quad=m(\pi^{-1}(u),m(\pi^{-1}(v),\pi^{-1}(w)))
    \end{aligned}}{\FormulaRefAuto{SemigroupAssociativityAxiom}[ax]{3,9,10,11}}
  \proofstepwidestarbreak[1,2,6,7,8]{\begin{aligned}[t]
    &m^\pi(u,m^\pi(v,w))\\[-2pt]
    &\quad=\pi(m(\pi^{-1}(u),m(\pi^{-1}(v),\pi^{-1}(w))))
    \end{aligned}}{\FormulaRefAuto{FiniteSemigroupRelabelledOperationDef}[def]{1,2,6,19},\rIE{17}}
  \proofstepwidestarbreak[1,2,6,7,8]{m^\pi(m^\pi(u,v),w)=m^\pi(u,m^\pi(v,w))}{\rIE{20,21,22}}
  \proofstepwidestarbreak[1,2]{\forall u,v,w\in S\,(m^\pi(m^\pi(u,v),w)=m^\pi(u,m^\pi(v,w)))}{\rUI{\rRI{6,\rUI{\rRI{7,\rUI{\rRI{8,23}}}}}}}
  \proofstepwidestar[1,2]{\SemigroupAlg{S}{m^\pi}}{\FormulaRefAuto{SemigroupDef}[def]{5,24}}
  \proofstepwidestar[1,2]{m^\pi\in\operatorname{Assoc}(S)}{\FormulaRefAuto{FiniteSemigroupAssociativeOperationsDef}[def]{5,25}}
\closeproofpartwideindDelta
\proofpartwide{Identität und Komposition}
  \proofstepwidestar[1]{m\in\operatorname{Assoc}(S)}{\rA}
  \proofstepwidestar[2]{\pi\in\operatorname{Sym}(S)}{\rA}
  \proofstepwidestar[3]{\sigma\in\operatorname{Sym}(S)}{\rA}
  \proofstepwidestar[]{\Id_S\colon S\bij S}{\FormulaRefAuto{\Id_A\colon A\bij A}[thm]}
  \proofstepwidestar[]{\Id_S\in\operatorname{Sym}(S)}{\FormulaRefAuto{FiniteSemigroupSymmetricSetDef}[def]{4}}
  \proofstepwidestar[2]{\pi\colon S\bij S}{\FormulaRefAuto{FiniteSemigroupSymmetricSetBijection}[thm]{2}}
  \proofstepwidestar[3]{\sigma\colon S\bij S}{\FormulaRefAuto{FiniteSemigroupSymmetricSetBijection}[thm]{3}}
  \proofstepwidestar[2,3]{\sigma\circ\pi\colon S\bij S}{\FormulaRefAuto{BijectiveFunctionComposition}[thm]{6,7}}
  \proofstepwidestar[2,3]{\sigma\circ\pi\in\operatorname{Sym}(S)}{\FormulaRefAuto{FiniteSemigroupSymmetricSetDef}[def]{8}}
  \proofstepwidestar[1,2]{m^\pi\in\operatorname{Assoc}(S)}{\FormulaRefAuto{B37RelabelledOperationAssociative}[thm]{1,2}}
  \proofstepwidestar[1]{m^{\Id_S}\colon S\times S\to S}{\FormulaRefAuto{FiniteSemigroupRelabelledOperationDef}[def]{1,5}}
  \proofstepwidestar[1]{m\colon S\times S\to S}{\FormulaRefAuto{FiniteSemigroupAssociativeOperationsDef}[def]{1}}
  \proofstepwidestar[13]{u\in S}{\rA}
  \proofstepwidestar[14]{v\in S}{\rA}
  \proofstepwidestarbreak[1,13,14]{m^{\Id_S}(u,v)=m(u,v)}{\FormulaRefAuto{FiniteSemigroupRelabelledOperationDef}[def]{1,5,13,14},\FormulaRefAuto{\Id_A^{-1}=\Id_A}[thm],\FormulaRefAuto{IdentityFunctionDef}[def]}
  \proofstepwidestar[1]{\forall u,v\in S\,(m^{\Id_S}(u,v)=m(u,v))}{\rUI{\rRI{13,\rUI{\rRI{14,15}}}}}
  \proofstepwidestar[1]{m^{\Id_S}=m}{\FormulaRefAuto{FunctionExtensionality}[thm]{11,12,16}}
  \proofstepwidestarbreak[2,3]{\pi^{-1}\circ\sigma^{-1}=(\sigma\circ\pi)^{-1}}{\FormulaRefAuto{InverseFunctionComposition}[thm]{6,7}}
  \proofstepwidestar[1,2,3]{(m^\pi)^\sigma\colon S\times S\to S}{\FormulaRefAuto{FiniteSemigroupRelabelledOperationDef}[def]{10,3}}
  \proofstepwidestar[1,2,3]{m^{\sigma\circ\pi}\colon S\times S\to S}{\FormulaRefAuto{FiniteSemigroupRelabelledOperationDef}[def]{1,9}}
  \proofstepwidestar[21]{u\in S}{\rA}
  \proofstepwidestar[22]{v\in S}{\rA}
  \proofstepwidestarbreak[1,2,3,21,22]{\begin{aligned}[t]
    &(m^\pi)^\sigma(u,v)\\[-2pt]
    &\quad=\sigma\bigl(\pi(m(\pi^{-1}(\sigma^{-1}(u)),
                             \pi^{-1}(\sigma^{-1}(v))))\bigr)
    \end{aligned}}{\FormulaRefAuto{FiniteSemigroupRelabelledOperationDef}[def]{10,3,21,22;1,2},\FormulaRefAuto{InverseFunctionValueType}[thm]{7,21;7,22}}
  \proofstepwidestarbreak[1,2,3,21,22]{\begin{aligned}[t]
    &m^{\sigma\circ\pi}(u,v)\\[-2pt]
    &\quad=\sigma\bigl(\pi(m(\pi^{-1}(\sigma^{-1}(u)),
                             \pi^{-1}(\sigma^{-1}(v))))\bigr)
    \end{aligned}}{\FormulaRefAuto{FiniteSemigroupRelabelledOperationDef}[def]{1,9,21,22},\FormulaRefAuto{CompositionDef}[def],\rIE{18}}
  \proofstepwidestar[1,2,3,21,22]{(m^\pi)^\sigma(u,v)=m^{\sigma\circ\pi}(u,v)}{\rIE{23,24}}
  \proofstepwidestarbreak[1,2,3]{\forall u,v\in S\,((m^\pi)^\sigma(u,v)=m^{\sigma\circ\pi}(u,v))}{\rUI{\rRI{21,\rUI{\rRI{22,25}}}}}
  \proofstepwidestar[1,2,3]{(m^\pi)^\sigma=m^{\sigma\circ\pi}}{\FormulaRefAuto{FunctionExtensionality}[thm]{19,20,26}}
  \proofstepwidestarbreak[1,2,3]{\begin{aligned}[t]
    &\Id_S\in\operatorname{Sym}(S)\land\sigma\circ\pi\in\operatorname{Sym}(S)\\[-2pt]
    &{}\land m^\pi\in\operatorname{Assoc}(S)\land m^{\Id_S}=m
      \land(m^\pi)^\sigma=m^{\sigma\circ\pi}
    \end{aligned}}{\rAI{5,\rAI{9,\rAI{10,\rAI{17,27}}}}}
\closeproofpartwide
\end{tabproofsplit}

\endgroup

\subsection{Graph der Wirkung}
\hypertarget{table-isomorphisms.proof.FiniteSemigroupRelabelGraphCharacterization}{}
\noindent Beweis zu \FormulaRefAuto{FiniteSemigroupRelabelGraphCharacterization}[thm].
% Beweis zu FiniteSemigroupRelabelGraphCharacterization.
\begin{tabproofwide}
  \proofstepwidestar[1]{m\in\operatorname{Assoc}(S)}{\rA}
  \proofstepwidestar[2]{d\in\operatorname{Assoc}(S)}{\rA}
  \proofstepwidestar[3]{\pi\in\operatorname{Sym}(S)}{\rA}
  \proofstepwidestar[3]{\pi\colon S\bij S}{\FormulaRefAuto{FiniteSemigroupSymmetricSetBijection}[thm]{3}}
  \proofstepwidestar[2]{d\colon S\times S\to S}{\FormulaRefAuto{FiniteSemigroupAssociativeOperationsDef}[def]{2}}
  \proofstepwidestar[1,3]{m^\pi\colon S\times S\to S}{\FormulaRefAuto{FiniteSemigroupRelabelledOperationDef}[def]{1,3}}
  \proofstepwidestar[7]{\operatorname{Relab}_S(\pi;m,d)}{\rA}
  \proofstepwidestarbreak[1,7]{\forall x,y\in S\,(\pi(m(x,y))=d(\pi(x),\pi(y)))}{\FormulaRefAuto{FiniteSemigroupRelabelGraphDef}[def]{1,7}}
  \proofstepwidestar[9]{u\in S}{\rA}
  \proofstepwidestar[10]{v\in S}{\rA}
  \proofstepwidestar[3,9]{\pi^{-1}(u)\in S}{\FormulaRefAuto{InverseFunctionValueType}[thm]{4,9}}
  \proofstepwidestar[3,10]{\pi^{-1}(v)\in S}{\FormulaRefAuto{InverseFunctionValueType}[thm]{4,10}}
  \proofstepwidestarbreak[1,3,7,9,10]{\begin{aligned}[t]
    &\pi(m(\pi^{-1}(u),\pi^{-1}(v)))\\[-2pt]
    &\quad=d(\pi(\pi^{-1}(u)),\pi(\pi^{-1}(v)))
    \end{aligned}}{\rRE{\rUE{\rRE{\rUE{8},11}},12}}
  \proofstepwidestar[3,9]{\pi(\pi^{-1}(u))=u}{\FormulaRefAuto{InverseFunctionRightInverse}[thm]{4,9}}
  \proofstepwidestar[3,10]{\pi(\pi^{-1}(v))=v}{\FormulaRefAuto{InverseFunctionRightInverse}[thm]{4,10}}
  \proofstepwidestar[1,3,7,9,10]{m^\pi(u,v)=d(u,v)}{\FormulaRefAuto{FiniteSemigroupRelabelledOperationDef}[def]{1,3,9,10},\rIE{13,14,15}}
  \proofstepwidestarbreak[1,3,7]{\forall u,v\in S\,(m^\pi(u,v)=d(u,v))}{\rUI{\rRI{9,\rUI{\rRI{10,16}}}}}
  \proofstepwidestar[1,2,3,7]{m^\pi=d}{\FormulaRefAuto{FunctionExtensionality}[thm]{6,5,17}}
  \proofstepwidestar[1,2,3,7]{d=m^\pi}{\FormulaRefAuto{a=b\vdash b=a}[thm]{18}}
  \proofstepwidestar[1,2,3]{\operatorname{Relab}_S(\pi;m,d)\to d=m^\pi}{\rRI{7,19}}
  \proofstepwidestar[21]{d=m^\pi}{\rA}
  \proofstepwidestar[22]{x\in S}{\rA}
  \proofstepwidestar[23]{y\in S}{\rA}
  \proofstepwidestar[3,22]{\pi(x)\in S}{\FormulaRefAuto{F\colon A\bij B\dsep x\in A\vdash F(x)\in B}[thm]{4,22}}
  \proofstepwidestar[3,23]{\pi(y)\in S}{\FormulaRefAuto{F\colon A\bij B\dsep x\in A\vdash F(x)\in B}[thm]{4,23}}
  \proofstepwidestarbreak[1,3,22,23]{\begin{aligned}[t]
    &m^\pi(\pi(x),\pi(y))\\[-2pt]
    &\quad=\pi(m(\pi^{-1}(\pi(x)),\pi^{-1}(\pi(y))))
    \end{aligned}}{\FormulaRefAuto{FiniteSemigroupRelabelledOperationDef}[def]{1,3,24,25}}
  \proofstepwidestar[3,22]{\pi^{-1}(\pi(x))=x}{\FormulaRefAuto{InverseFunctionLeftInverse}[thm]{4,22}}
  \proofstepwidestar[3,23]{\pi^{-1}(\pi(y))=y}{\FormulaRefAuto{InverseFunctionLeftInverse}[thm]{4,23}}
  \proofstepwidestar[1,3,21,22,23]{\pi(m(x,y))=d(\pi(x),\pi(y))}{\rIE{21,27,28,26},\FormulaRefAuto{a=b\vdash b=a}[thm]}
  \proofstepwidestarbreak[1,3,21]{\forall x,y\in S\,(\pi(m(x,y))=d(\pi(x),\pi(y)))}{\rUI{\rRI{22,\rUI{\rRI{23,29}}}}}
  \proofstepwidestar[1,2,3,21]{\operatorname{Relab}_S(\pi;m,d)}{\FormulaRefAuto{FiniteSemigroupRelabelGraphDef}[def]{1,3,2,30}}
  \proofstepwidestar[1,2,3]{d=m^\pi\to\operatorname{Relab}_S(\pi;m,d)}{\rRI{21,31}}
  \proofstepwidestarbreak[1,2,3]{\operatorname{Relab}_S(\pi;m,d)\leftrightarrow d=m^\pi}{\rLRI{20,32}}
\end{tabproofwide}


\subsection{Identität, Komposition und Inverse von Automorphismen}
\hypertarget{table-isomorphisms.proof.automorphisms}{}
\FormulaThmDeltaK[Die Identität ist ein Automorphismus]{%
 m\in\operatorname{Assoc}(S)\vdash\Id_S\in\operatorname{Aut}(S,m)
}{B37TAutomorphismIdentity}{\DeltaRow{Mengen}{S}\DeltaRow{Funktionen}{m}}
\begin{tabproofwide}
 \proofstepwidestar[1]{m\in\operatorname{Assoc}(S)}{\rA}
 \proofstepwidestar[]{\Id_S\colon S\bij S}{\FormulaRefAuto{\Id_A\colon A\bij A}[thm]}
 \proofstepwidestar[]{\Id_S\in\operatorname{Sym}(S)}{\FormulaRefAuto{FiniteSemigroupSymmetricSetDef}[def]{2}}
 \proofstepwidestar[1]{m^{\Id_S}=m}{\rAEa{\rAEb{\rAEb{\rAEb{\FormulaRefAuto{FiniteSemigroupRelabellingAction}[thm]{1,\rAI{3,3}}}}}}}
 \proofstepwidestar[1]{\Id_S\in\operatorname{Aut}(S,m)}{\FormulaRefAuto{FiniteSemigroupAutomorphismGroupDef}[def]{1,3,4}}
\end{tabproofwide}

\FormulaThmDeltaK[Automorphismen sind unter Komposition abgeschlossen]{%
 \begin{aligned}[t]
 &m\in\operatorname{Assoc}(S)\dsep\pi,\sigma\in\operatorname{Aut}(S,m)\vdash{}\\[-2pt]
 &\sigma\circ\pi\in\operatorname{Aut}(S,m)
 \end{aligned}
}{B37TAutomorphismComposition}{\DeltaRow{Mengen}{S}\DeltaRow{Funktionen}{m\dsep\pi\dsep\sigma}}
\begin{tabproofwide}
 \proofstepwidestar[1]{m\in\operatorname{Assoc}(S)}{\rA}
 \proofstepwidestar[2]{\pi\in\operatorname{Aut}(S,m)}{\rA}
 \proofstepwidestar[3]{\sigma\in\operatorname{Aut}(S,m)}{\rA}
 \proofstepwidestar[1,2]{\pi\in\operatorname{Sym}(S)\land m^\pi=m}{\FormulaRefAuto{FiniteSemigroupAutomorphismGroupDef}[def]{1,2}}
 \proofstepwidestar[1,3]{\sigma\in\operatorname{Sym}(S)\land m^\sigma=m}{\FormulaRefAuto{FiniteSemigroupAutomorphismGroupDef}[def]{1,3}}
 \proofstepwidestar[1,2,3]{\begin{aligned}[t]&\Id_S\in\operatorname{Sym}(S)\land\sigma\circ\pi\in\operatorname{Sym}(S)\\[-2pt]&{}\land m^\pi\in\operatorname{Assoc}(S)\land m^{\Id_S}=m\land(m^\pi)^\sigma=m^{\sigma\circ\pi}\end{aligned}}{\FormulaRefAuto{FiniteSemigroupRelabellingAction}[thm]{1,\rAI{\rAEa{4},\rAEa{5}}}}
 \proofstepwidestar[1,2,3]{\sigma\circ\pi\in\operatorname{Sym}(S)}{\rAEa{\rAEb{6}}}
 \proofstepwidestar[1,2,3]{(m^\pi)^\sigma=m^{\sigma\circ\pi}}{\rAEb{\rAEb{\rAEb{\rAEb{6}}}}}
 \proofstepwidestar[1,2,3]{m^\sigma=m^{\sigma\circ\pi}}{\rIE{\rAEb{4},8}}
 \proofstepwidestar[1,2,3]{m^{\sigma\circ\pi}=m^\sigma}{\FormulaRefAuto{a=b\vdash b=a}[thm]{9}}
 \proofstepwidestar[1,2,3]{m^{\sigma\circ\pi}=m}{\rChain{10,\rAEb{5}}}
 \proofstepwidestar[1,2,3]{\sigma\circ\pi\in\operatorname{Aut}(S,m)}{\FormulaRefAuto{FiniteSemigroupAutomorphismGroupDef}[def]{1,7,11}}
\end{tabproofwide}

\FormulaThmDeltaK[Das Inverse eines Automorphismus ist ein Automorphismus]{%
 m\in\operatorname{Assoc}(S)\dsep\pi\in\operatorname{Aut}(S,m)
 \vdash\pi^{-1}\in\operatorname{Aut}(S,m)
}{B37TAutomorphismInverse}{\DeltaRow{Mengen}{S}\DeltaRow{Funktionen}{m\dsep\pi}}
\begin{tabproofwide}
 \proofstepwidestar[1]{m\in\operatorname{Assoc}(S)}{\rA}
 \proofstepwidestar[2]{\pi\in\operatorname{Aut}(S,m)}{\rA}
 \proofstepwidestar[1,2]{\pi\in\operatorname{Sym}(S)\land m^\pi=m}{\FormulaRefAuto{FiniteSemigroupAutomorphismGroupDef}[def]{1,2}}
 \proofstepwidestar[1,2]{\pi\colon S\bij S}{\FormulaRefAuto{FiniteSemigroupSymmetricSetBijection}[thm]{\rAEa{3}}}
 \proofstepwidestar[1,2]{\pi^{-1}\colon S\bij S}{\FormulaRefAuto{InverseFunctionBijection}[thm]{4}}
 \proofstepwidestar[1,2]{\pi^{-1}\in\operatorname{Sym}(S)}{\FormulaRefAuto{FiniteSemigroupSymmetricSetDef}[def]{5}}
 \proofstepwidestar[1,2]{\pi^{-1}\circ\pi=\Id_S}{\FormulaRefAuto{InverseFunctionLeftComposition}[thm]{4}}
 \proofstepwidestar[1,2]{(m^\pi)^{\pi^{-1}}=m^{\pi^{-1}\circ\pi}}{\rAEb{\rAEb{\rAEb{\rAEb{\FormulaRefAuto{FiniteSemigroupRelabellingAction}[thm]{1,\rAI{\rAEa{3},6}}}}}}}
 \proofstepwidestar[1,2]{m^{\pi^{-1}}=m^{\Id_S}}{\rIE{\rAEb{3},7,8}}
 \proofstepwidestar[1,2]{m^{\Id_S}=m}{\rAEa{\rAEb{\rAEb{\rAEb{\FormulaRefAuto{FiniteSemigroupRelabellingAction}[thm]{1,\rAI{\rAEa{3},6}}}}}}}
 \proofstepwidestar[1,2]{m^{\pi^{-1}}=m}{\rChain{9,10}}
 \proofstepwidestar[1,2]{\pi^{-1}\in\operatorname{Aut}(S,m)}{\FormulaRefAuto{FiniteSemigroupAutomorphismGroupDef}[def]{1,6,11}}
\end{tabproofwide}
Die drei Tabellen beweisen die behauptete Stabilität unter Identität,
Komposition und Inversen. Assoziativität der Komposition und die
Identitäts- und Inversengesetze werden von den Selbstbijektionen geerbt;
eine Gruppendefinition aus einem späteren Band wird hier nicht vorausgesetzt.

\subsection{Orbits und Isomorphieklassen}
\hypertarget{table-isomorphisms.proof.FiniteSemigroupRelabelOrbitsAreIsomorphismClasses}{}
\noindent Beweis zu \FormulaRefAuto{FiniteSemigroupRelabelOrbitsAreIsomorphismClasses}[thm].
% Beweis zu FiniteSemigroupRelabelOrbitsAreIsomorphismClasses.
\begin{tabproofwide}
  \proofstepwidestar[1]{\Finite S}{\rA}
  \proofstepwidestar[2]{m\in\operatorname{Assoc}(S)}{\rA}
  \proofstepwidestar[3]{d\in\operatorname{Assoc}(S)}{\rA}
  \proofstepwidestar[2]{\SemigroupAlg{S}{m}}{\FormulaRefAuto{FiniteSemigroupAssociativeOperationSemigroup}[thm]{2}}
  \proofstepwidestar[3]{\SemigroupAlg{S}{d}}{\FormulaRefAuto{FiniteSemigroupAssociativeOperationSemigroup}[thm]{3}}
  \proofstepwidestar[1,2]{\FiniteSemigroupAlg{S}{m}}{\FormulaRefAuto{FiniteSemigroupDef}[def]{4,1}}
  \proofstepwidestar[1,3]{\FiniteSemigroupAlg{S}{d}}{\FormulaRefAuto{FiniteSemigroupDef}[def]{5,1}}
  \proofstepwidestar[8]{d\in\operatorname{Orb}_S(m)}{\rA}
  \proofstepwidestarbreak[2,8]{\exists\pi\,(\pi\in\operatorname{Sym}(S)\land\operatorname{Relab}_S(\pi;m,d))}{\FormulaRefAuto{FiniteSemigroupRelabelOrbitDef}[def]{2,8}}
  \proofstepwidestar[10]{\pi\in\operatorname{Sym}(S)\land\operatorname{Relab}_S(\pi;m,d)}{\rA}
  \proofstepwidestar[10]{\pi\in\operatorname{Sym}(S)}{\rAEa{10}}
  \proofstepwidestar[10]{\pi\colon S\bij S}{\FormulaRefAuto{FiniteSemigroupSymmetricSetBijection}[thm]{11}}
  \proofstepwidestar[10]{\pi\colon S\to S}{\FormulaRefAuto{Bijektive Funktion}[def]{12}}
  \proofstepwidestarbreak[2,10]{\forall x,y\in S\,(\pi(m(x,y))=d(\pi(x),\pi(y)))}{\FormulaRefAuto{FiniteSemigroupRelabelGraphDef}[def]{2,\rAEb{10}}}
  \proofstepwidestar[2,3,10]{\SemigroupHom{\pi}{S,m}{S,d}}{\FormulaRefAuto{SemigroupHomDef}[def]{4,5,13,14}}
  \proofstepwidestar[1,2,3,10]{\FiniteSemigroupHom{\pi}{S,m}{S,d}}{\FormulaRefAuto{FiniteSemigroupHomDef}[def]{6,7,15}}
  \proofstepwidestar[1,2,3,10]{\FiniteSemigroupIso{\pi}{S,m}{S,d}}{\FormulaRefAuto{FiniteSemigroupIsoDef}[def]{16,12}}
  \proofstepwidestar[1,2,3,10]{\exists f\,\FiniteSemigroupIso{f}{S,m}{S,d}}{\rEI{17}}
  \proofstepwidestar[1,2,3,8]{\exists f\,\FiniteSemigroupIso{f}{S,m}{S,d}}{\rEE{9,10,18}}
  \proofstepwidestarbreak[1,2,3]{d\in\operatorname{Orb}_S(m)\to\exists f\,\FiniteSemigroupIso{f}{S,m}{S,d}}{\rRI{8,19}}
  \proofstepwidestar[21]{\exists f\,\FiniteSemigroupIso{f}{S,m}{S,d}}{\rA}
  \proofstepwidestar[22]{\FiniteSemigroupIso{f}{S,m}{S,d}}{\rA}
  \proofstepwidestar[22]{f\colon S\bij S}{\FormulaRefAuto{FiniteSemigroupIsoBijectiveAxiom}[ax]{22}}
  \proofstepwidestar[22]{f\in\operatorname{Sym}(S)}{\FormulaRefAuto{FiniteSemigroupSymmetricSetDef}[def]{23}}
  \proofstepwidestar[22]{\FiniteSemigroupHom{f}{S,m}{S,d}}{\FormulaRefAuto{FiniteSemigroupIsoHomAxiom}[ax]{22}}
  \proofstepwidestar[22]{\SemigroupHom{f}{S,m}{S,d}}{\FormulaRefAuto{FiniteSemigroupHomBaseHomAxiom}[ax]{25}}
  \proofstepwidestar[27]{x\in S}{\rA}
  \proofstepwidestar[28]{y\in S}{\rA}
  \proofstepwidestar[22,27,28]{f(m(x,y))=d(f(x),f(y))}{\FormulaRefAuto{SemigroupHomOperationAxiom}[ax]{26,27,28}}
  \proofstepwidestar[22]{\forall x,y\in S\,(f(m(x,y))=d(f(x),f(y)))}{\rUI{\rRI{27,\rUI{\rRI{28,29}}}}}
  \proofstepwidestar[2,3,22]{\operatorname{Relab}_S(f;m,d)}{\FormulaRefAuto{FiniteSemigroupRelabelGraphDef}[def]{2,24,3,30}}
  \proofstepwidestarbreak[2,3,22]{\exists\pi\,(\pi\in\operatorname{Sym}(S)\land\operatorname{Relab}_S(\pi;m,d))}{\rEI{\rAI{24,31}}}
  \proofstepwidestar[2,3,22]{d\in\operatorname{Orb}_S(m)}{\FormulaRefAuto{FiniteSemigroupRelabelOrbitDef}[def]{2,3,32}}
  \proofstepwidestar[2,3,21]{d\in\operatorname{Orb}_S(m)}{\rEE{21,22,33}}
  \proofstepwidestarbreak[2,3]{(\exists f\,\FiniteSemigroupIso{f}{S,m}{S,d})\to d\in\operatorname{Orb}_S(m)}{\rRI{21,34}}
  \proofstepwidestarbreak[1,2,3]{d\in\operatorname{Orb}_S(m)\leftrightarrow\exists f\,\FiniteSemigroupIso{f}{S,m}{S,d}}{\rLRI{20,35}}
\end{tabproofwide}


\section{Tafeltupel und vollständiger Code}
\subsection{Die Tafel aus ihrem Tupel zurückgewinnen}
\hypertarget{table-isomorphisms.proof.FiniteSemigroupTableTupleInjective}{}
\noindent Beweis zu \FormulaRefAuto{FiniteSemigroupTableTupleInjective}[thm].
% Beweis zu FiniteSemigroupTableTupleInjective.
Setze \(S=\NatSeg n\). Die beiden Tupel werden als verschachtelte
endliche Folgen gelesen; die äußere Folge nummeriert die Zeilen.
\begin{tabproofwide}
  \proofstepwidestar[1]{n\in\mathbb N}{\rA}
  \proofstepwidestar[2]{m\colon S\times S\to S}{\rA}
  \proofstepwidestar[3]{d\colon S\times S\to S}{\rA}
  \proofstepwidestar[4]{w_n(m)=w_n(d)}{\rA}
  \proofstepwidestar[5]{i\in S}{\rA}
  \proofstepwidestar[6]{j\in S}{\rA}
  \proofstepwidestarbreak[1,2,3,4,5]{(m(i,j))_{j\in S}=(d(i,j))_{j\in S}}{\FormulaRefAuto{FiniteSemigroupTableTupleDef}[def]{1,2,3},\rIE{4}}
  \proofstepwidestar[]{\bigl((m(i,k))_{k\in S}\bigr)(j)=\bigl((m(i,k))_{k\in S}\bigr)(j)}{\rII}
  \proofstepwidestar[1,2,3,4,5]{\bigl((m(i,k))_{k\in S}\bigr)(j)=\bigl((d(i,k))_{k\in S}\bigr)(j)}{\rIE{7,8}}
  \proofstepwidestar[1,2,3,4,5,6]{m(i,j)=d(i,j)}{\FormulaRefAuto{FiniteSemigroupTableTupleDef}[def]{1,2,3,5,6,9}}
  \proofstepwidestar[1,2,3,4]{\forall i,j\in S\,(m(i,j)=d(i,j))}{\rUI{\rRI{5,\rUI{\rRI{6,10}}}}}
  \proofstepwidestar[1,2,3,4]{m=d}{\FormulaRefAuto{FunctionExtensionality}[thm]{2,3,11}}
\end{tabproofwide}


\subsection{Vollständigkeit des kanonischen Codes}
\hypertarget{table-isomorphisms.proof.FiniteSemigroupCanonicalTableCodeComplete}{}
\noindent Beweis zu \FormulaRefAuto{FiniteSemigroupCanonicalTableCodeComplete}[thm].
\begingroup
\renewcommand{\theformulaThm}{37.4.1.9}
% Beweis zu FiniteSemigroupCanonicalTableCodeComplete.
Wir kürzen \(H_A=\FiniteSemigroupAlg{A}{\star}\),
\(H_B=\FiniteSemigroupAlg{B}{\diamond}\),
\(n_A=\FiniteCard(A)\), \(n_B=\FiniteCard(B)\) und
\[
 \begin{aligned}
 W_A&=\{w_{n_A}(m_{\star,e})\mid e\colon\NatSeg{n_A}\bij A\},\\
 W_B&=\{w_{n_B}(m_{\diamond,e})\mid e\colon\NatSeg{n_B}\bij B\}
 \end{aligned}
\]
ab. Somit ist \(\operatorname{can}(A,\star)=(n_A,W_A)\),
und entsprechend für \(B\).
\begin{tabproofsplit}
\proofpartwideindDelta[Ein Isomorphismus erhält jede nummerierte Tafel]{%
  \begin{aligned}[t]
  &H_A\dsep H_B\dsep
    \FiniteSemigroupIso{f}{A,\star}{B,\diamond}\dsep
    n\in\mathbb N\dsep e\colon\NatSeg n\bij A\vdash{}\\[-2pt]
  &f\circ e\colon\NatSeg n\bij B\land
    m_{\diamond,f\circ e}=m_{\star,e}
  \end{aligned}
}{%
  \DeltaRow{Mengen}{A\dsep B\dsep n\dsep i\dsep j}
  \DeltaRow{Funktionen}{f\dsep e\dsep\star\dsep\diamond}
}[][B37TableCodeIsomorphismTransport]
  \proofstepwidestar[1]{H_A}{\rA}
  \proofstepwidestar[2]{H_B}{\rA}
  \proofstepwidestar[3]{\FiniteSemigroupIso{f}{A,\star}{B,\diamond}}{\rA}
  \proofstepwidestar[4]{n\in\mathbb N}{\rA}
  \proofstepwidestar[5]{e\colon\NatSeg n\bij A}{\rA}
  \proofstepwidestar[3]{f\colon A\bij B}{\FormulaRefAuto{FiniteSemigroupIsoBijectiveAxiom}[ax]{3}}
  \proofstepwidestar[3,5]{f\circ e\colon\NatSeg n\bij B}{\FormulaRefAuto{BijectiveFunctionComposition}[thm]{5,6}}
  \proofstepwidestar[3]{\FiniteSemigroupHom{f}{A,\star}{B,\diamond}}{\FormulaRefAuto{FiniteSemigroupIsoHomAxiom}[ax]{3}}
  \proofstepwidestar[3]{\SemigroupHom{f}{A,\star}{B,\diamond}}{\FormulaRefAuto{FiniteSemigroupHomBaseHomAxiom}[ax]{8}}
  \proofstepwidestar[10]{i\in\NatSeg n}{\rA}
  \proofstepwidestar[11]{j\in\NatSeg n}{\rA}
  \proofstepwidestar[5,10]{e(i)\in A}{\FormulaRefAuto{F\colon A\bij B\dsep x\in A\vdash F(x)\in B}[thm]{5,10}}
  \proofstepwidestar[5,11]{e(j)\in A}{\FormulaRefAuto{F\colon A\bij B\dsep x\in A\vdash F(x)\in B}[thm]{5,11}}
  \proofstepwidestarbreak[1,4,5,10,11]{m_{\star,e}(i,j)\in\NatSeg n}{\FormulaRefAuto{FiniteSemigroupMultiplicationTableDef}[def]{1,4,5,10,11}}
  \proofstepwidestarbreak[1,4,5,10,11]{e(m_{\star,e}(i,j))=e(i)\star e(j)}{\FormulaRefAuto{B37TableTransportEvaluation}[thm]{1,4,5,10,11}}
  \proofstepwidestarbreak[3,5,10,11]{f(e(i)\star e(j))=f(e(i))\diamond f(e(j))}{\FormulaRefAuto{SemigroupHomOperationAxiom}[ax]{9,12,13}}
  \proofstepwidestarbreak[1,3,4,5,10,11]{\begin{aligned}[t]
    &(f\circ e)(m_{\star,e}(i,j))\\[-2pt]
    &\quad=(f\circ e)(i)\diamond(f\circ e)(j)
    \end{aligned}}{\FormulaRefAuto{CompositionDef}[def]{5,6},\rIE{15,16}}
  \proofstepwidestarbreak[2,3,4,5,10,11]{\begin{aligned}[t]
    &m_{\diamond,f\circ e}(i,j)\\[-2pt]
    &\quad=(f\circ e)^{-1}((f\circ e)(i)\diamond(f\circ e)(j))
    \end{aligned}}{\FormulaRefAuto{FiniteSemigroupMultiplicationTableDef}[def]{2,4,7,10,11}}
  \proofstepwidestarbreak[1,2,3,4,5,10,11]{m_{\diamond,f\circ e}(i,j)=m_{\star,e}(i,j)}{\rIE{17,18},\FormulaRefAuto{InverseFunctionLeftInverse}[thm]{7,14}}
  \proofstepwidestarbreak[1,2,3,4,5]{\forall i,j\in\NatSeg n\,(m_{\diamond,f\circ e}(i,j)=m_{\star,e}(i,j))}{\rUI{\rRI{10,\rUI{\rRI{11,19}}}}}
  \proofstepwidestar[1,4,5]{m_{\star,e}\colon\NatSeg n\times\NatSeg n\to\NatSeg n}{\FormulaRefAuto{FiniteSemigroupMultiplicationTableDef}[def]{1,4,5}}
  \proofstepwidestar[2,3,4,5]{m_{\diamond,f\circ e}\colon\NatSeg n\times\NatSeg n\to\NatSeg n}{\FormulaRefAuto{FiniteSemigroupMultiplicationTableDef}[def]{2,4,7}}
  \proofstepwidestar[1,2,3,4,5]{m_{\diamond,f\circ e}=m_{\star,e}}{\FormulaRefAuto{FunctionExtensionality}[thm]{22,21,20}}
  \proofstepwidestarbreak[1,2,3,4,5]{f\circ e\colon\NatSeg n\bij B\land m_{\diamond,f\circ e}=m_{\star,e}}{\rAI{7,23}}
\closeproofpartwideindDelta
\proofpartwideindDelta[Gleiche nummerierte Tafeln liefern einen Isomorphismus]{%
  \begin{aligned}[t]
  &H_A\dsep H_B\dsep n\in\mathbb N\dsep
    e_A\colon\NatSeg n\bij A\dsep e_B\colon\NatSeg n\bij B\dsep{}\\[-2pt]
  &m_{\star,e_A}=m_{\diamond,e_B}\vdash
    \FiniteSemigroupIso{e_B\circ e_A^{-1}}{A,\star}{B,\diamond}
  \end{aligned}
}{%
  \DeltaRow{Mengen}{A\dsep B\dsep n\dsep x\dsep y}
  \DeltaRow{Funktionen}{e_A\dsep e_B\dsep\star\dsep\diamond}
}[][B37TableCodeCommonTableIsomorphism]
  \proofstepwidestar[1]{H_A}{\rA}
  \proofstepwidestar[2]{H_B}{\rA}
  \proofstepwidestar[3]{n\in\mathbb N}{\rA}
  \proofstepwidestar[4]{e_A\colon\NatSeg n\bij A}{\rA}
  \proofstepwidestar[5]{e_B\colon\NatSeg n\bij B}{\rA}
  \proofstepwidestar[6]{m_{\star,e_A}=m_{\diamond,e_B}}{\rA}
  \proofstepwidestar[4]{e_A^{-1}\colon A\bij\NatSeg n}{\FormulaRefAuto{InverseFunctionBijection}[thm]{4}}
  \proofstepwidestar[4,5]{e_B\circ e_A^{-1}\colon A\bij B}{\FormulaRefAuto{BijectiveFunctionComposition}[thm]{7,5}}
  \proofstepwidestar[9]{x\in A}{\rA}
  \proofstepwidestar[10]{y\in A}{\rA}
  \proofstepwidestar[4,9]{e_A^{-1}(x)\in\NatSeg n}{\FormulaRefAuto{InverseFunctionValueType}[thm]{4,9}}
  \proofstepwidestar[4,10]{e_A^{-1}(y)\in\NatSeg n}{\FormulaRefAuto{InverseFunctionValueType}[thm]{4,10}}
  \proofstepwidestar[4,9]{e_A(e_A^{-1}(x))=x}{\FormulaRefAuto{InverseFunctionRightInverse}[thm]{4,9}}
  \proofstepwidestar[4,10]{e_A(e_A^{-1}(y))=y}{\FormulaRefAuto{InverseFunctionRightInverse}[thm]{4,10}}
  \proofstepwidestarbreak[1,3,4,9,10]{m_{\star,e_A}(e_A^{-1}(x),e_A^{-1}(y))\in\NatSeg n}{\FormulaRefAuto{FiniteSemigroupMultiplicationTableDef}[def]{1,3,4,11,12}}
  \proofstepwidestarbreak[1,3,4,9,10]{\begin{aligned}[t]
    &e_A(m_{\star,e_A}(e_A^{-1}(x),e_A^{-1}(y)))\\[-2pt]
    &\quad=x\star y
    \end{aligned}}{\FormulaRefAuto{B37TableTransportEvaluation}[thm]{1,3,4,11,12},\rIE{13,14}}
  \proofstepwidestarbreak[1,3,4,9,10]{e_A^{-1}(x\star y)=m_{\star,e_A}(e_A^{-1}(x),e_A^{-1}(y))}{\rIE{16,\FormulaRefAuto{InverseFunctionLeftInverse}[thm]{4,15}}}
  \proofstepwidestarbreak[2,3,4,5,9,10]{\begin{aligned}[t]
    &e_B(m_{\diamond,e_B}(e_A^{-1}(x),e_A^{-1}(y)))\\[-2pt]
    &\quad=e_B(e_A^{-1}(x))\diamond e_B(e_A^{-1}(y))
    \end{aligned}}{\FormulaRefAuto{B37TableTransportEvaluation}[thm]{2,3,5,11,12}}
  \proofstepwidestar[1]{\SemigroupAlg{A}{\star}}{\FormulaRefAuto{FiniteSemigroupBaseAxiom}[ax]{1}}
  \proofstepwidestar[1,9,10]{x\star y\in A}{\FormulaRefAuto{SemigroupClosureAxiom}[ax]{19,9,10}}
  \proofstepwidestarbreak[1,2,3,4,5,6,9,10]{\begin{aligned}[t]
    &(e_B\circ e_A^{-1})(x\star y)\\[-2pt]
    &\quad=(e_B\circ e_A^{-1})(x)\diamond(e_B\circ e_A^{-1})(y)
    \end{aligned}}{\FormulaRefAuto{CompositionDef}[def]{7,5},\rIE{6,17,18},20}
  \proofstepwidestarbreak[1,2,3,4,5,6]{\begin{aligned}[t]
    &\forall x,y\in A\,\bigl((e_B\circ e_A^{-1})(x\star y)\\[-2pt]
    &\qquad=(e_B\circ e_A^{-1})(x)\diamond(e_B\circ e_A^{-1})(y)\bigr)
    \end{aligned}}{\rUI{\rRI{9,\rUI{\rRI{10,21}}}}}
  \proofstepwidestar[2]{\SemigroupAlg{B}{\diamond}}{\FormulaRefAuto{FiniteSemigroupBaseAxiom}[ax]{2}}
  \proofstepwidestar[4,5]{e_B\circ e_A^{-1}\colon A\to B}{\FormulaRefAuto{Bijektive Funktion}[def]{8}}
  \proofstepwidestar[1,2,3,4,5,6]{\SemigroupHom{e_B\circ e_A^{-1}}{A,\star}{B,\diamond}}{\FormulaRefAuto{SemigroupHomDef}[def]{19,23,24,22}}
  \proofstepwidestar[1,2,3,4,5,6]{\FiniteSemigroupHom{e_B\circ e_A^{-1}}{A,\star}{B,\diamond}}{\FormulaRefAuto{FiniteSemigroupHomDef}[def]{1,2,25}}
  \proofstepwidestar[1,2,3,4,5,6]{\FiniteSemigroupIso{e_B\circ e_A^{-1}}{A,\star}{B,\diamond}}{\FormulaRefAuto{FiniteSemigroupIsoDef}[def]{26,8}}
\closeproofpartwideindDelta
\proofpartwideindDelta[Isomorphie erhält die beiden Codekomponenten]{%
  H_A\dsep H_B\dsep\FiniteSemigroupIso{f}{A,\star}{B,\diamond}
  \vdash\operatorname{can}(A,\star)=\operatorname{can}(B,\diamond)
}{%
  \DeltaRow{Mengen}{A\dsep B}
  \DeltaRow{Funktionen}{f\dsep\star\dsep\diamond}
}[][B37TableCodeIsomorphismInvariant]
  \proofstepwidestar[1]{H_A}{\rA}
  \proofstepwidestar[2]{H_B}{\rA}
  \proofstepwidestar[3]{\FiniteSemigroupIso{f}{A,\star}{B,\diamond}}{\rA}
  \proofstepwidestar[1]{\Finite A}{\FormulaRefAuto{FiniteSemigroupFinitenessAxiom}[ax]{1}}
  \proofstepwidestar[2]{\Finite B}{\FormulaRefAuto{FiniteSemigroupFinitenessAxiom}[ax]{2}}
  \proofstepwidestar[1]{n_A\in\mathbb N}{\FormulaRefAuto{FiniteCardNatural}[thm]{4}}
  \proofstepwidestar[2]{n_B\in\mathbb N}{\FormulaRefAuto{FiniteCardNatural}[thm]{5}}
  \proofstepwidestar[1]{\NatSeg{n_A}\EqCard A}{\FormulaRefAuto{FiniteCardCounts}[thm]{4}}
  \proofstepwidestar[2]{\NatSeg{n_B}\EqCard B}{\FormulaRefAuto{FiniteCardCounts}[thm]{5}}
  \proofstepwidestar[3]{f\colon A\bij B}{\FormulaRefAuto{FiniteSemigroupIsoBijectiveAxiom}[ax]{3}}
  \proofstepwidestar[3]{A\EqCard B}{\FormulaRefAuto{Gleichmächtigkeit}[def]{\rEI{10}}}
  \proofstepwidestar[1,3]{\NatSeg{n_A}\EqCard B}{\FormulaRefAuto{EqCardTransitive}[thm]{8,11}}
  \proofstepwidestar[1,2,3]{n_A=n_B}{\FormulaRefAuto{FiniteCardWitnessUnique}[thm]{6,7,12,9}}
  \proofstepwidestar[14]{w\in W_A}{\rA}
  \proofstepwidestarbreak[14]{\exists e\,(e\colon\NatSeg{n_A}\bij A\land w=w_{n_A}(m_{\star,e}))}{\FormulaRefAuto{FiniteSemigroupCanonicalTableCodeDef}[def]{1,14}}
  \proofstepwidestar[16]{e\colon\NatSeg{n_A}\bij A\land w=w_{n_A}(m_{\star,e})}{\rA}
  \proofstepwidestarbreak[1,2,3,16]{f\circ e\colon\NatSeg{n_A}\bij B\land m_{\diamond,f\circ e}=m_{\star,e}}{\FormulaRefAuto{B37TableCodeIsomorphismTransport}[thm]{1,2,3,6,\rAEa{16}}}
  \proofstepwidestar[1,2,3,16]{w\in W_B}{\FormulaRefAuto{FiniteSemigroupCanonicalTableCodeDef}[def]{2,\rAEa{17}},\rIE{13,\rAEb{16},\rAEb{17}}}
  \proofstepwidestar[1,2,3,14]{w\in W_B}{\rEE{15,16,18}}
  \proofstepwidestar[1,2,3]{\forall w\,(w\in W_A\to w\in W_B)}{\rUI{\rRI{14,19}}}
  \proofstepwidestar[1,2,3]{W_A\subseteq W_B}{\FormulaRefAuto{A\subseteq B\coloneq\forall x\,(x\in A\to x\in B)}[def]{20}}
  \proofstepwidestar[22]{w\in W_B}{\rA}
  \proofstepwidestarbreak[22]{\exists g\,(g\colon\NatSeg{n_B}\bij B\land w=w_{n_B}(m_{\diamond,g}))}{\FormulaRefAuto{FiniteSemigroupCanonicalTableCodeDef}[def]{2,22}}
  \proofstepwidestar[24]{g\colon\NatSeg{n_B}\bij B\land w=w_{n_B}(m_{\diamond,g})}{\rA}
  \proofstepwidestar[3]{f^{-1}\colon B\bij A}{\FormulaRefAuto{InverseFunctionBijection}[thm]{10}}
  \proofstepwidestar[3,24]{f^{-1}\circ g\colon\NatSeg{n_B}\bij A}{\FormulaRefAuto{BijectiveFunctionComposition}[thm]{\rAEa{24},25}}
  \proofstepwidestar[27]{i\in\NatSeg{n_B}}{\rA}
  \proofstepwidestar[24,27]{g(i)\in B}{\FormulaRefAuto{F\colon A\bij B\dsep x\in A\vdash F(x)\in B}[thm]{\rAEa{24},27}}
  \proofstepwidestar[3,24,27]{(f\circ(f^{-1}\circ g))(i)=g(i)}{\FormulaRefAuto{CompositionDef}[def],\FormulaRefAuto{InverseFunctionRightInverse}[thm]{10,28}}
  \proofstepwidestar[3,24]{\forall i\in\NatSeg{n_B}\,((f\circ(f^{-1}\circ g))(i)=g(i))}{\rUI{\rRI{27,29}}}
  \proofstepwidestar[3,24]{f\circ(f^{-1}\circ g)=g}{\FormulaRefAuto{FunctionExtensionality}[thm]{\FormulaRefAuto{BijectiveFunctionComposition}[thm]{26,10},\rAEa{24},30},\FormulaRefAuto{Bijektive Funktion}[def]}
  \proofstepwidestarbreak[1,2,3,24]{\begin{aligned}[t]
    &f\circ(f^{-1}\circ g)\colon\NatSeg{n_B}\bij B\\[-2pt]
    &{}\land m_{\diamond,f\circ(f^{-1}\circ g)}=m_{\star,f^{-1}\circ g}
    \end{aligned}}{\FormulaRefAuto{B37TableCodeIsomorphismTransport}[thm]{1,2,3,7,26}}
  \proofstepwidestar[1,2,3,24]{w\in W_A}{\FormulaRefAuto{FiniteSemigroupCanonicalTableCodeDef}[def]{1,26},\rIE{13,31,\rAEb{32},\rAEb{24}}}
  \proofstepwidestar[1,2,3,22]{w\in W_A}{\rEE{23,24,33}}
  \proofstepwidestar[1,2,3]{\forall w\,(w\in W_B\to w\in W_A)}{\rUI{\rRI{22,34}}}
  \proofstepwidestar[1,2,3]{W_B\subseteq W_A}{\FormulaRefAuto{A\subseteq B\coloneq\forall x\,(x\in A\to x\in B)}[def]{35}}
  \proofstepwidestar[1,2,3]{W_A=W_B}{\FormulaRefAuto{A\subseteq B\dsep B\subseteq A\vdash A=B}[thm]{21,36}}
  \proofstepwidestar[1,2,3]{\operatorname{can}(A,\star)=\operatorname{can}(B,\diamond)}{\FormulaRefAuto{FiniteSemigroupCanonicalTableCodeDef}[def]{1,2},\rIE{13,37}}
\closeproofpartwideindDelta
\proofpartwideindDelta[Gleiche Codekomponenten erzwingen Isomorphie]{%
  H_A\dsep H_B\dsep\operatorname{can}(A,\star)=\operatorname{can}(B,\diamond)
  \vdash\exists f\,\FiniteSemigroupIso{f}{A,\star}{B,\diamond}
}{%
  \DeltaRow{Mengen}{A\dsep B}
  \DeltaRow{Funktionen}{f\dsep\star\dsep\diamond}
}[][B37TableCodeEqualCodesIsomorphism]
  \proofstepwidestar[1]{H_A}{\rA}
  \proofstepwidestar[2]{H_B}{\rA}
  \proofstepwidestar[3]{\operatorname{can}(A,\star)=\operatorname{can}(B,\diamond)}{\rA}
  \proofstepwidestar[1,2,3]{n_A=n_B\land W_A=W_B}{\FormulaRefAuto{FiniteSemigroupCanonicalTableCodeDef}[def]{1,2},\FormulaRefAuto{(a,b)=(c,d)\eqvdash a=c\land b=d}[thm]{3}}
  \proofstepwidestar[1]{\Finite A}{\FormulaRefAuto{FiniteSemigroupFinitenessAxiom}[ax]{1}}
  \proofstepwidestar[1]{n_A\in\mathbb N}{\FormulaRefAuto{FiniteCardNatural}[thm]{5}}
  \proofstepwidestar[1]{\NatSeg{n_A}\EqCard A}{\FormulaRefAuto{FiniteCardCounts}[thm]{5}}
  \proofstepwidestar[1]{\exists e_A\,(e_A\colon\NatSeg{n_A}\bij A)}{\FormulaRefAuto{Gleichmächtigkeit}[def]{7}}
  \proofstepwidestar[9]{e_A\colon\NatSeg{n_A}\bij A}{\rA}
  \proofstepwidestar[1,9]{w_{n_A}(m_{\star,e_A})\in W_A}{\FormulaRefAuto{FiniteSemigroupCanonicalTableCodeDef}[def]{1,9}}
  \proofstepwidestar[1,2,3,9]{w_{n_A}(m_{\star,e_A})\in W_B}{\rIE{\rAEb{4},10}}
  \proofstepwidestarbreak[1,2,3,9]{\exists e_B\,(e_B\colon\NatSeg{n_A}\bij B\land w_{n_A}(m_{\star,e_A})=w_{n_A}(m_{\diamond,e_B}))}{\FormulaRefAuto{FiniteSemigroupCanonicalTableCodeDef}[def]{2,11},\rIE{\rAEa{4}}}
  \proofstepwidestarbreak[13]{e_B\colon\NatSeg{n_A}\bij B\land w_{n_A}(m_{\star,e_A})=w_{n_A}(m_{\diamond,e_B})}{\rA}
  \proofstepwidestar[1,2,13,9]{m_{\star,e_A}=m_{\diamond,e_B}}{\FormulaRefAuto{FiniteSemigroupTableTupleInjective}[thm]{6,\FormulaRefAuto{FiniteSemigroupMultiplicationTableDef}[def]{1,6,9},\FormulaRefAuto{FiniteSemigroupMultiplicationTableDef}[def]{2,6,\rAEa{13}},\rAEb{13}}}
  \proofstepwidestar[1,2,13,9]{\FiniteSemigroupIso{e_B\circ e_A^{-1}}{A,\star}{B,\diamond}}{\FormulaRefAuto{B37TableCodeCommonTableIsomorphism}[thm]{1,2,6,9,\rAEa{13},14}}
  \proofstepwidestar[1,2,13,9]{\exists f\,\FiniteSemigroupIso{f}{A,\star}{B,\diamond}}{\rEI{15}}
  \proofstepwidestar[1,2,3,9]{\exists f\,\FiniteSemigroupIso{f}{A,\star}{B,\diamond}}{\rEE{12,13,16}}
  \proofstepwidestar[1,2,3]{\exists f\,\FiniteSemigroupIso{f}{A,\star}{B,\diamond}}{\rEE{8,9,17}}
\closeproofpartwideindDelta
\proofpartwide{Die Äquivalenz}
  \proofstepwidestar[1]{H_A}{\rA}
  \proofstepwidestar[2]{H_B}{\rA}
  \proofstepwidestar[3]{\exists f\,\FiniteSemigroupIso{f}{A,\star}{B,\diamond}}{\rA}
  \proofstepwidestar[4]{\FiniteSemigroupIso{f}{A,\star}{B,\diamond}}{\rA}
  \proofstepwidestar[1,2,4]{\operatorname{can}(A,\star)=\operatorname{can}(B,\diamond)}{\FormulaRefAuto{B37TableCodeIsomorphismInvariant}[thm]{1,2,4}}
  \proofstepwidestar[1,2,3]{\operatorname{can}(A,\star)=\operatorname{can}(B,\diamond)}{\rEE{3,4,5}}
  \proofstepwidestarbreak[1,2]{(\exists f\,\FiniteSemigroupIso{f}{A,\star}{B,\diamond})\to\operatorname{can}(A,\star)=\operatorname{can}(B,\diamond)}{\rRI{3,6}}
  \proofstepwidestar[8]{\operatorname{can}(A,\star)=\operatorname{can}(B,\diamond)}{\rA}
  \proofstepwidestar[1,2,8]{\exists f\,\FiniteSemigroupIso{f}{A,\star}{B,\diamond}}{\FormulaRefAuto{B37TableCodeEqualCodesIsomorphism}[thm]{1,2,8}}
  \proofstepwidestarbreak[1,2]{\operatorname{can}(A,\star)=\operatorname{can}(B,\diamond)\to\exists f\,\FiniteSemigroupIso{f}{A,\star}{B,\diamond}}{\rRI{8,9}}
  \proofstepwidestarbreak[1,2]{(\exists f\,\FiniteSemigroupIso{f}{A,\star}{B,\diamond})\leftrightarrow\operatorname{can}(A,\star)=\operatorname{can}(B,\diamond)}{\rLRI{7,10}}
\closeproofpartwide
\end{tabproofsplit}

\endgroup

\section{Beliebige Träger und wechselnde Nummerierungen}
\hypertarget{table-isomorphisms.proof.arbitrary-carriers}{}
\subsection{Isomorphismen mit und ohne Endlichkeitszusatz}
\FormulaThmDeltaK[Der Endlichkeitszusatz verändert den Isomorphismusbegriff nicht]{%
 \begin{aligned}[t]
 &\FiniteSemigroupAlg{A}{\star}\dsep\FiniteSemigroupAlg{B}{\diamond}\vdash{}\\[-2pt]
 &\FiniteSemigroupIso{f}{A,\star}{B,\diamond}
 \leftrightarrow\SemigroupIso{f}{A,\star}{B,\diamond}
 \end{aligned}
}{B37TIsoBridge}{\DeltaRow{Mengen}{A\dsep B}\DeltaRow{Funktionen}{f\dsep\star\dsep\diamond}}
\begin{tabproofwide}
 \proofstepwidestar[1]{\FiniteSemigroupAlg{A}{\star}}{\rA}
 \proofstepwidestar[2]{\FiniteSemigroupAlg{B}{\diamond}}{\rA}
 \proofstepwidestar[3]{\FiniteSemigroupIso{f}{A,\star}{B,\diamond}}{\rA}
 \proofstepwidestar[3]{\FiniteSemigroupHom{f}{A,\star}{B,\diamond}}{\FormulaRefAuto{FiniteSemigroupIsoHomAxiom}[ax]{3}}
 \proofstepwidestar[3]{\SemigroupHom{f}{A,\star}{B,\diamond}}{\FormulaRefAuto{FiniteSemigroupHomBaseHomAxiom}[ax]{4}}
 \proofstepwidestar[3]{f\colon A\bij B}{\FormulaRefAuto{FiniteSemigroupIsoBijectiveAxiom}[ax]{3}}
 \proofstepwidestar[3]{\SemigroupIso{f}{A,\star}{B,\diamond}}{\FormulaRefAuto{SemigroupIsoDef}[def]{5,6}}
 \proofstepwidestar[]{\FiniteSemigroupIso{f}{A,\star}{B,\diamond}\to\SemigroupIso{f}{A,\star}{B,\diamond}}{\rRI{3,7}}
 \proofstepwidestar[9]{\SemigroupIso{f}{A,\star}{B,\diamond}}{\rA}
 \proofstepwidestar[9]{\SemigroupHom{f}{A,\star}{B,\diamond}}{\FormulaRefAuto{SemigroupIsoHomAxiom}[ax]{9}}
 \proofstepwidestar[9]{f\colon A\bij B}{\FormulaRefAuto{SemigroupIsoBijectiveAxiom}[ax]{9}}
 \proofstepwidestar[1,2,9]{\FiniteSemigroupHom{f}{A,\star}{B,\diamond}}{\FormulaRefAuto{FiniteSemigroupHomDef}[def]{1,2,10}}
 \proofstepwidestar[1,2,9]{\FiniteSemigroupIso{f}{A,\star}{B,\diamond}}{\FormulaRefAuto{FiniteSemigroupIsoDef}[def]{12,11}}
 \proofstepwidestar[1,2]{\SemigroupIso{f}{A,\star}{B,\diamond}\to\FiniteSemigroupIso{f}{A,\star}{B,\diamond}}{\rRI{9,13}}
 \proofstepwidestar[1,2]{\FiniteSemigroupIso{f}{A,\star}{B,\diamond}\leftrightarrow\SemigroupIso{f}{A,\star}{B,\diamond}}{\rLRI{8,14}}
\end{tabproofwide}
\subsection{Das vollständige Kriterium für beliebige endliche Träger}
\FormulaThmDeltaK[Isomorphie ist gleichzeitige Umbenennung der beiden Tafeln]{%
 \begin{aligned}[t]
 &\FiniteSemigroupAlg{A}{\star}\dsep\FiniteSemigroupAlg{B}{\diamond}\dsep n\in\mathbb N\\[-2pt]
 &\dsep e_A\colon\NatSeg n\bij A\dsep e_B\colon\NatSeg n\bij B\vdash{}\\[-2pt]
 &\bigl(\exists f\,\FiniteSemigroupIso{f}{A,\star}{B,\diamond}\bigr)\leftrightarrow{}\\[-2pt]
 &\exists\pi\,\Bigl(\pi\colon\NatSeg n\bij\NatSeg n\land
 \forall i,j\in\NatSeg n\,
 \bigl(\pi(m_{\star,e_A}(i,j))=m_{\diamond,e_B}(\pi(i),\pi(j))\bigr)\Bigr)
 \end{aligned}
}{B37TArbitraryCarrierCriterion}{\DeltaRow{Mengen}{A\dsep B\dsep n\dsep i\dsep j}\DeltaRow{Funktionen}{e_A\dsep e_B\dsep f\dsep\pi\dsep\star\dsep\diamond}}
Wir schreiben \(S=\NatSeg n\), \(m=m_{\star,e_A}\),
\(d=m_{\diamond,e_B}\), \(H_A=\FiniteSemigroupAlg{A}{\star}\),
\(H_B=\FiniteSemigroupAlg{B}{\diamond}\) sowie
\[
 Q(\pi)\equiv\forall i,j\in S\,
   \bigl(\pi(m(i,j))=d(\pi(i),\pi(j))\bigr).
\]
Im Vorwärtszweig bedeutet \(h=e_B^{-1}\circ(f\circ e_A)\), im
Rückwärtszweig \(g=e_B\circ(\pi\circ e_A^{-1})\).
\begin{tabproofwide}
 \proofstepwidestar[1]{H_A}{\rA}
 \proofstepwidestar[2]{H_B}{\rA}
 \proofstepwidestar[3]{n\in\mathbb N}{\rA}
 \proofstepwidestar[4]{e_A\colon S\bij A}{\rA}
 \proofstepwidestar[5]{e_B\colon S\bij B}{\rA}
 \proofstepwidestar[1,3,4]{\FiniteSemigroupAlg{S}{m}\land\FiniteSemigroupIso{e_A}{S,m}{A,\star}}{\FormulaRefAuto{FiniteSemigroupMultiplicationTableTransport}[thm]{1,3,4}}
 \proofstepwidestar[2,3,5]{\FiniteSemigroupAlg{S}{d}\land\FiniteSemigroupIso{e_B}{S,d}{B,\diamond}}{\FormulaRefAuto{FiniteSemigroupMultiplicationTableTransport}[thm]{2,3,5}}
 \proofstepwidestar[1,3,4]{\SemigroupIso{e_A}{S,m}{A,\star}}{\rRE{\rLREa{\FormulaRefAuto{B37TIsoBridge}[thm]{\rAEa{6},1}},\rAEb{6}}}
 \proofstepwidestar[2,3,5]{\SemigroupIso{e_B}{S,d}{B,\diamond}}{\rRE{\rLREa{\FormulaRefAuto{B37TIsoBridge}[thm]{\rAEa{7},2}},\rAEb{7}}}
 \proofstepwidestar[2,3,5]{\SemigroupIso{e_B^{-1}}{B,\diamond}{S,d}}{\FormulaRefAuto{SemigroupIsoInverse}[thm]{9}}
 \proofstepwidestar[1,3,4]{\SemigroupIso{e_A^{-1}}{A,\star}{S,m}}{\FormulaRefAuto{SemigroupIsoInverse}[thm]{8}}
 \proofstepwidestar[12]{\exists f\,\FiniteSemigroupIso{f}{A,\star}{B,\diamond}}{\rA}
 \proofstepwidestar[13]{\FiniteSemigroupIso{f}{A,\star}{B,\diamond}}{\rA}
 \proofstepwidestar[1,2,13]{\SemigroupIso{f}{A,\star}{B,\diamond}}{\rRE{\rLREa{\FormulaRefAuto{B37TIsoBridge}[thm]{1,2}},13}}
 \proofstepwidestar[1,2,3,4,13]{\SemigroupIso{f\circ e_A}{S,m}{B,\diamond}}{\FormulaRefAuto{SemigroupIsoComposition}[thm]{8,14}}
 \proofstepwidestar[1--5,13]{\SemigroupIso{h}{S,m}{S,d}}{\FormulaRefAuto{SemigroupIsoComposition}[thm]{15,10}}
 \proofstepwidestar[1--5,13]{\FiniteSemigroupIso{h}{S,m}{S,d}}{\rRE{\rLREb{\FormulaRefAuto{B37TIsoBridge}[thm]{\rAEa{6},\rAEa{7}}},16}}
 \proofstepwidestar[1--5,13]{h\colon S\bij S}{\FormulaRefAuto{FiniteSemigroupIsoBijectiveAxiom}[ax]{17}}
 \proofstepwidestar[1--5,13]{Q(h)}{\rRE{\rLREa{\FormulaRefAuto{FiniteSemigroupTableRelabelCriterion}[thm]{3,\rAEa{6},\rAEa{7},18}},17}}
 \proofstepwidestar[1--5,13]{\exists\pi\,(\pi\colon S\bij S\land Q(\pi))}{\rEI{\rAI{18,19}}}
 \proofstepwidestar[1--5,12]{\exists\pi\,(\pi\colon S\bij S\land Q(\pi))}{\rEE{12,13,20}}
 \proofstepwidestar[1--5]{\begin{aligned}[t]&(\exists f\,\FiniteSemigroupIso{f}{A,\star}{B,\diamond})\\[-2pt]&\quad\to\exists\pi\,(\pi\colon S\bij S\land Q(\pi))\end{aligned}}{\rRI{12,21}}
 \proofstepwidestar[23]{\exists\pi\,(\pi\colon S\bij S\land Q(\pi))}{\rA}
 \proofstepwidestar[24]{\pi\colon S\bij S\land Q(\pi)}{\rA}
 \proofstepwidestar[1--5,24]{\FiniteSemigroupIso{\pi}{S,m}{S,d}}{\rRE{\rLREb{\FormulaRefAuto{FiniteSemigroupTableRelabelCriterion}[thm]{3,\rAEa{6},\rAEa{7},\rAEa{24}}},\rAEb{24}}}
 \proofstepwidestar[1--5,24]{\SemigroupIso{\pi}{S,m}{S,d}}{\rRE{\rLREa{\FormulaRefAuto{B37TIsoBridge}[thm]{\rAEa{6},\rAEa{7}}},25}}
 \proofstepwidestar[1--5,24]{\SemigroupIso{\pi\circ e_A^{-1}}{A,\star}{S,d}}{\FormulaRefAuto{SemigroupIsoComposition}[thm]{11,26}}
 \proofstepwidestar[1--5,24]{\SemigroupIso{g}{A,\star}{B,\diamond}}{\FormulaRefAuto{SemigroupIsoComposition}[thm]{27,9}}
 \proofstepwidestar[1--5,24]{\FiniteSemigroupIso{g}{A,\star}{B,\diamond}}{\rRE{\rLREb{\FormulaRefAuto{B37TIsoBridge}[thm]{1,2}},28}}
 \proofstepwidestar[1--5,24]{\exists f\,\FiniteSemigroupIso{f}{A,\star}{B,\diamond}}{\rEI{29}}
 \proofstepwidestar[1--5,23]{\exists f\,\FiniteSemigroupIso{f}{A,\star}{B,\diamond}}{\rEE{23,24,30}}
 \proofstepwidestar[1--5]{\begin{aligned}[t]&\exists\pi\,(\pi\colon S\bij S\land Q(\pi))\\[-2pt]&\quad\to\exists f\,\FiniteSemigroupIso{f}{A,\star}{B,\diamond}\end{aligned}}{\rRI{23,31}}
 \proofstepwidestar[1--5]{\begin{aligned}[t]&(\exists f\,\FiniteSemigroupIso{f}{A,\star}{B,\diamond})\\[-2pt]&\quad\leftrightarrow\exists\pi\,(\pi\colon S\bij S\land Q(\pi))\end{aligned}}{\rLRI{22,32}}
\end{tabproofwide}
Die im Hauptband umgekehrt geschriebene Produktgleichung ist wegen der
Symmetrie der Gleichheit dieselbe Bedingung. Keine Tabellenzeile setzt
voraus, dass \(S\), \(A\) oder \(B\) ein Element enthält.

\subsection{Nummerierungswechsel und die Kandidatenmenge}
\hypertarget{table-isomorphisms.proof.renumbering}{}
\FormulaThmDeltaK[Doppelte Umkehrung einer Bijektion]{%
 f\colon A\bij B\vdash(f^{-1})^{-1}=f
}{B37TDoubleInverse}{\DeltaRow{Mengen}{A\dsep B\dsep x}\DeltaRow{Funktionen}{f}}
\begin{tabproofwide}
 \proofstepwidestar[1]{f\colon A\bij B}{\rA}
 \proofstepwidestar[1]{f^{-1}\colon B\bij A}{\FormulaRefAuto{InverseFunctionBijection}[thm]{1}}
 \proofstepwidestar[1]{f\colon A\to B}{\FormulaRefAuto{Bijektive Funktion}[def]{1}}
 \proofstepwidestar[4]{x\in A}{\rA}
 \proofstepwidestar[1,4]{f^{-1}(f(x))=x}{\FormulaRefAuto{InverseFunctionLeftInverse}[thm]{1,4}}
 \proofstepwidestar[1]{\forall x\in A\,(f^{-1}(f(x))=x)}{\rUI{\rRI{4,5}}}
 \proofstepwidestar[1]{f=(f^{-1})^{-1}}{\FormulaRefAuto{InverseFunctionRightInverseUniqueness}[thm]{2,3,6}}
 \proofstepwidestar[1]{(f^{-1})^{-1}=f}{\FormulaRefAuto{a=b\vdash b=a}[thm]{7}}
\end{tabproofwide}

\Needspace{14\baselineskip}
\FormulaThmDeltaK[Jede weitere Nummerierung entsteht durch eine Permutation]{%
 \begin{aligned}[t]
 &e,g\colon S\bij A\vdash{}\\[-2pt]
 &\exists\rho\in\operatorname{Sym}(S)\,(g=e\circ\rho)
 \end{aligned}
}{B37TNumberingFactorization}{\DeltaRow{Mengen}{S\dsep A\dsep x}\DeltaRow{Funktionen}{e\dsep g\dsep\rho}}
In der Tabelle steht \(\rho=e^{-1}\circ g\).
\begin{tabproofwide}
 \proofstepwidestar[1]{e\colon S\bij A}{\rA}
 \proofstepwidestar[2]{g\colon S\bij A}{\rA}
 \proofstepwidestar[1]{e^{-1}\colon A\bij S}{\FormulaRefAuto{InverseFunctionBijection}[thm]{1}}
 \proofstepwidestar[1,2]{\rho\colon S\bij S}{\FormulaRefAuto{BijectiveFunctionComposition}[thm]{2,3}}
 \proofstepwidestar[1,2]{\rho\in\operatorname{Sym}(S)}{\FormulaRefAuto{FiniteSemigroupSymmetricSetDef}[def]{4}}
 \proofstepwidestar[1,2]{e\circ\rho\colon S\bij A}{\FormulaRefAuto{BijectiveFunctionComposition}[thm]{4,1}}
 \proofstepwidestar[7]{x\in S}{\rA}
 \proofstepwidestar[2,7]{g(x)\in A}{\FormulaRefAuto{F\colon A\bij B\dsep x\in A\vdash F(x)\in B}[thm]{2,7}}
 \proofstepwidestar[1,2,7]{\rho(x)=e^{-1}(g(x))}{\FormulaRefAuto{CompositionDef}[def]{\FormulaRefAuto{Bijektive Funktion}[def]{2},\FormulaRefAuto{Bijektive Funktion}[def]{3},7}}
 \proofstepwidestar[1,2,7]{(e\circ\rho)(x)=e(\rho(x))}{\FormulaRefAuto{CompositionDef}[def]{\FormulaRefAuto{Bijektive Funktion}[def]{4},\FormulaRefAuto{Bijektive Funktion}[def]{1},7}}
 \proofstepwidestar[1,2,7]{e(e^{-1}(g(x)))=g(x)}{\FormulaRefAuto{InverseFunctionRightInverse}[thm]{1,8}}
 \proofstepwidestar[1,2,7]{(e\circ\rho)(x)=g(x)}{\rChain{\rIE{9,10},11}}
 \proofstepwidestar[1,2]{\forall x\in S\,((e\circ\rho)(x)=g(x))}{\rUI{\rRI{7,12}}}
 \proofstepwidestar[1,2]{e\circ\rho=g}{\FormulaRefAuto{FunctionExtensionality}[thm]{\FormulaRefAuto{Bijektive Funktion}[def]{6},\FormulaRefAuto{Bijektive Funktion}[def]{2},13}}
 \proofstepwidestar[1,2]{g=e\circ\rho}{\FormulaRefAuto{a=b\vdash b=a}[thm]{14}}
 \proofstepwidestar[1,2]{\exists\rho\in\operatorname{Sym}(S)\,(g=e\circ\rho)}{\rEI{\rAI{5,15}}}
\end{tabproofwide}

\Needspace{14\baselineskip}
\FormulaThmDeltaK[Nummerierungswechsel ist Umbenennung mit der inversen Permutation]{%
 \begin{aligned}[t]
 &\FiniteSemigroupAlg{A}{\star}\dsep n\in\mathbb N\dsep e\colon\NatSeg n\bij A
 \dsep\rho\in\operatorname{Sym}(\NatSeg n)\vdash{}\\[-2pt]
 &e\circ\rho\colon\NatSeg n\bij A\land
   m_{\star,e\circ\rho}=m_{\star,e}^{\rho^{-1}}
 \end{aligned}
}{B37TRenumbering}{\DeltaRow{Mengen}{A\dsep n\dsep i\dsep j}\DeltaRow{Funktionen}{e\dsep\rho\dsep\star}}
Wir schreiben \(S=\NatSeg n\), \(h=e\circ\rho\),
\(m=m_{\star,e}\) und \(d=m_{\star,h}\).
\begin{tabproofwide}
 \proofstepwidestar[1]{\FiniteSemigroupAlg{A}{\star}}{\rA}
 \proofstepwidestar[2]{n\in\mathbb N}{\rA}
 \proofstepwidestar[3]{e\colon S\bij A}{\rA}
 \proofstepwidestar[4]{\rho\in\operatorname{Sym}(S)}{\rA}
 \proofstepwidestar[4]{\rho\colon S\bij S}{\FormulaRefAuto{FiniteSemigroupSymmetricSetBijection}[thm]{4}}
 \proofstepwidestar[3,4]{h\colon S\bij A}{\FormulaRefAuto{BijectiveFunctionComposition}[thm]{5,3}}
 \proofstepwidestar[4]{\rho^{-1}\colon S\bij S}{\FormulaRefAuto{InverseFunctionBijection}[thm]{5}}
 \proofstepwidestar[4]{\rho^{-1}\in\operatorname{Sym}(S)}{\FormulaRefAuto{FiniteSemigroupSymmetricSetDef}[def]{7}}
 \proofstepwidestar[4]{(\rho^{-1})^{-1}=\rho}{\FormulaRefAuto{B37TDoubleInverse}[thm]{5}}
 \proofstepwidestar[1,2,3]{m\colon S\times S\to S}{\FormulaRefAuto{FiniteSemigroupMultiplicationTableDef}[def]{1,2,3}}
 \proofstepwidestar[1,2,3]{\SemigroupAlg{S}{m}}{\FormulaRefAuto{B37TableTransportSemigroup}[thm]{1,2,3}}
 \proofstepwidestar[1,2,3]{m\in\operatorname{Assoc}(S)}{\FormulaRefAuto{FiniteSemigroupAssociativeOperationsDef}[def]{10,11}}
 \proofstepwidestar[1,2,3,4]{d\colon S\times S\to S}{\FormulaRefAuto{FiniteSemigroupMultiplicationTableDef}[def]{1,2,6}}
 \proofstepwidestar[1,2,3,4]{m^{\rho^{-1}}\colon S\times S\to S}{\FormulaRefAuto{FiniteSemigroupRelabelledOperationDef}[def]{12,8}}
 \proofstepwidestar[3,4]{\rho^{-1}\circ e^{-1}=h^{-1}}{\FormulaRefAuto{InverseFunctionComposition}[thm]{5,3}}
 \proofstepwidestar[16]{i\in S}{\rA}
 \proofstepwidestar[17]{j\in S}{\rA}
 \proofstepwidestar[4,16]{\rho(i)\in S}{\FormulaRefAuto{F\colon A\bij B\dsep x\in A\vdash F(x)\in B}[thm]{5,16}}
 \proofstepwidestar[4,17]{\rho(j)\in S}{\FormulaRefAuto{F\colon A\bij B\dsep x\in A\vdash F(x)\in B}[thm]{5,17}}
 \proofstepwidestar[3,4,16]{h(i)=e(\rho(i))}{\FormulaRefAuto{CompositionDef}[def]{\FormulaRefAuto{Bijektive Funktion}[def]{5},\FormulaRefAuto{Bijektive Funktion}[def]{3},16}}
 \proofstepwidestar[3,4,17]{h(j)=e(\rho(j))}{\FormulaRefAuto{CompositionDef}[def]{\FormulaRefAuto{Bijektive Funktion}[def]{5},\FormulaRefAuto{Bijektive Funktion}[def]{3},17}}
 \proofstepwidestar[1]{\SemigroupAlg{A}{\star}}{\FormulaRefAuto{FiniteSemigroupBaseAxiom}[ax]{1}}
 \proofstepwidestar[1,3,4,16,17]{e(\rho(i))\star e(\rho(j))\in A}{\FormulaRefAuto{SemigroupClosureAxiom}[ax]{22,\FormulaRefAuto{F\colon A\bij B\dsep x\in A\vdash F(x)\in B}[thm]{3,18},\FormulaRefAuto{F\colon A\bij B\dsep x\in A\vdash F(x)\in B}[thm]{3,19}}}
 \proofstepwidestar[1,2,3,4,16,17]{d(i,j)=h^{-1}(e(\rho(i))\star e(\rho(j)))}{\rIE{20,21,\FormulaRefAuto{FiniteSemigroupMultiplicationTableDef}[def]{1,2,6,16,17}}}
 \proofstepwidestar[1,2,3,4,16,17]{\begin{aligned}[t]&d(i,j)\\[-2pt]&\quad=\rho^{-1}(e^{-1}(e(\rho(i))\star e(\rho(j))))\end{aligned}}{\rChain{\rIE{\FormulaRefAuto{a=b\vdash b=a}[thm]{15},24},\FormulaRefAuto{CompositionDef}[def]{\FormulaRefAuto{InverseFunctionType}[thm]{3},\FormulaRefAuto{Bijektive Funktion}[def]{7},23}}}
 \proofstepwidestar[1,2,3,4,16,17]{m(\rho(i),\rho(j))=e^{-1}(e(\rho(i))\star e(\rho(j)))}{\FormulaRefAuto{FiniteSemigroupMultiplicationTableDef}[def]{1,2,3,18,19}}
 \proofstepwidestar[1,2,3,4,16,17]{d(i,j)=\rho^{-1}(m(\rho(i),\rho(j)))}{\rIE{\FormulaRefAuto{a=b\vdash b=a}[thm]{26},25}}
 \proofstepwidestar[1,2,3,4,16,17]{m^{\rho^{-1}}(i,j)=\rho^{-1}(m(\rho(i),\rho(j)))}{\rIE{9,\FormulaRefAuto{FiniteSemigroupRelabelledOperationDef}[def]{12,8,16,17}}}
 \proofstepwidestar[1,2,3,4,16,17]{d(i,j)=m^{\rho^{-1}}(i,j)}{\rChain{27,\FormulaRefAuto{a=b\vdash b=a}[thm]{28}}}
 \proofstepwidestar[1,2,3,4]{\forall i,j\in S\,(d(i,j)=m^{\rho^{-1}}(i,j))}{\rUI{\rRI{16,\rUI{\rRI{17,29}}}}}
 \proofstepwidestar[1,2,3,4]{d=m^{\rho^{-1}}}{\FormulaRefAuto{FunctionExtensionality}[thm]{13,14,30}}
 \proofstepwidestar[1,2,3,4]{h\colon S\bij A\land d=m^{\rho^{-1}}}{\rAI{6,31}}
\end{tabproofwide}
\Needspace{14\baselineskip}
\FormulaDefDeltaK[Tafeltupel aller Umbenennungen]{%
 \begin{aligned}[t]
 &n\in\mathbb N\dsep m\in\operatorname{Assoc}(\NatSeg n)\vdash{}\\[-2pt]
 &\mathcal W_n(m)\coloneqq
   \{w_n(m^\pi)\mid\pi\in\operatorname{Sym}(\NatSeg n)\}
 \end{aligned}
}{B37TOrbitTupleCandidatesDef}{\DeltaRow{Mengen}{n}\DeltaRow{Funktionen}{m\dsep\pi}\DeltaRow{\textbf{Neues Symbol}}{}\DeltaPrem{Orbit-Tafeltupel}{\mathcal W_n(m)}}

\Needspace{14\baselineskip}
\FormulaThmDeltaK[Eine Ausgangstafel liefert alle Codekandidaten]{%
 \begin{aligned}[t]
 &\FiniteSemigroupAlg{A}{\star}\dsep
 e\colon\NatSeg{\FiniteCard(A)}\bij A\vdash{}\\[-2pt]
 &\{w_{\FiniteCard(A)}(m_{\star,g})\mid
       g\colon\NatSeg{\FiniteCard(A)}\bij A\}
   =\mathcal W_{\FiniteCard(A)}(m_{\star,e})
 \end{aligned}
}{B37TCodeOrbitCandidates}{\DeltaRow{Mengen}{A\dsep n\dsep w}\DeltaRow{Funktionen}{e\dsep g\dsep\rho\dsep\pi\dsep\star}}
Für diese Tabelle schreiben wir
\[
 \begin{aligned}
 n&=\FiniteCard(A),\quad S=\NatSeg n,\quad m=m_{\star,e},\\
 W&=\{w_n(m_{\star,g})\mid g\colon S\bij A\},
 \qquad V=\mathcal W_n(m).
 \end{aligned}
\]
\begin{tabproofwide}
 \proofstepwidestar[1]{\FiniteSemigroupAlg{A}{\star}}{\rA}
 \proofstepwidestar[2]{e\colon S\bij A}{\rA}
 \proofstepwidestar[1]{\Finite A}{\FormulaRefAuto{FiniteSemigroupFinitenessAxiom}[ax]{1}}
 \proofstepwidestar[1]{n\in\mathbb N}{\FormulaRefAuto{FiniteCardNatural}[thm]{3}}
 \proofstepwidestar[1,2]{m\colon S\times S\to S}{\FormulaRefAuto{FiniteSemigroupMultiplicationTableDef}[def]{1,4,2}}
 \proofstepwidestar[1,2]{\SemigroupAlg{S}{m}}{\FormulaRefAuto{B37TableTransportSemigroup}[thm]{1,4,2}}
 \proofstepwidestar[1,2]{m\in\operatorname{Assoc}(S)}{\FormulaRefAuto{FiniteSemigroupAssociativeOperationsDef}[def]{5,6}}
 \proofstepwidestar[8]{w\in W}{\rA}
 \proofstepwidestar[1,8]{\exists g\,(g\colon S\bij A\land w=w_n(m_{\star,g}))}{\FormulaRefAuto{FiniteSemigroupCanonicalTableCodeDef}[def]{1,8}}
 \proofstepwidestar[10]{g\colon S\bij A\land w=w_n(m_{\star,g})}{\rA}
 \proofstepwidestar[2,10]{\exists\rho\in\operatorname{Sym}(S)\,(g=e\circ\rho)}{\FormulaRefAuto{B37TNumberingFactorization}[thm]{\rAI{2,\rAEa{10}}}}
 \proofstepwidestar[12]{\rho\in\operatorname{Sym}(S)\land g=e\circ\rho}{\rA}
 \proofstepwidestar[1,2,12]{e\circ\rho\colon S\bij A\land m_{\star,e\circ\rho}=m^{\rho^{-1}}}{\FormulaRefAuto{B37TRenumbering}[thm]{1,4,2,\rAEa{12}}}
 \proofstepwidestar[12]{\rho\colon S\bij S}{\FormulaRefAuto{FiniteSemigroupSymmetricSetBijection}[thm]{\rAEa{12}}}
 \proofstepwidestar[12]{\rho^{-1}\colon S\bij S}{\FormulaRefAuto{InverseFunctionBijection}[thm]{14}}
 \proofstepwidestar[12]{\rho^{-1}\in\operatorname{Sym}(S)}{\FormulaRefAuto{FiniteSemigroupSymmetricSetDef}[def]{15}}
 \proofstepwidestar[1,2,12]{m_{\star,g}=m^{\rho^{-1}}}{\rIE{\FormulaRefAuto{a=b\vdash b=a}[thm]{\rAEb{12}},\rAEb{13}}}
 \proofstepwidestar[1,2,10,12]{w=w_n(m^{\rho^{-1}})}{\rIE{17,\rAEb{10}}}
 \proofstepwidestar[1,2,10,12]{w\in V}{\FormulaRefAuto{B37TOrbitTupleCandidatesDef}[def]{4,7,\rEI{\rAI{16,18}}}}
 \proofstepwidestar[1,2,10]{w\in V}{\rEE{11,12,19}}
 \proofstepwidestar[1,2,8]{w\in V}{\rEE{9,10,20}}
 \proofstepwidestar[1,2]{W\subseteq V}{\FormulaRefAuto{A\subseteq B\coloneq\forall x\,(x\in A\to x\in B)}[def]{\rUI{\rRI{8,21}}}}
 \proofstepwidestar[23]{w\in V}{\rA}
 \proofstepwidestar[1,2,23]{\exists\pi\in\operatorname{Sym}(S)\,(w=w_n(m^\pi))}{\FormulaRefAuto{B37TOrbitTupleCandidatesDef}[def]{4,7,23}}
 \proofstepwidestar[25]{\pi\in\operatorname{Sym}(S)\land w=w_n(m^\pi)}{\rA}
 \proofstepwidestar[25]{\pi\colon S\bij S}{\FormulaRefAuto{FiniteSemigroupSymmetricSetBijection}[thm]{\rAEa{25}}}
 \proofstepwidestar[25]{\pi^{-1}\colon S\bij S}{\FormulaRefAuto{InverseFunctionBijection}[thm]{26}}
 \proofstepwidestar[25]{\pi^{-1}\in\operatorname{Sym}(S)}{\FormulaRefAuto{FiniteSemigroupSymmetricSetDef}[def]{27}}
 \proofstepwidestar[1,2,25]{e\circ\pi^{-1}\colon S\bij A\land m_{\star,e\circ\pi^{-1}}=m^{(\pi^{-1})^{-1}}}{\FormulaRefAuto{B37TRenumbering}[thm]{1,4,2,28}}
 \proofstepwidestar[25]{(\pi^{-1})^{-1}=\pi}{\FormulaRefAuto{B37TDoubleInverse}[thm]{26}}
 \proofstepwidestar[1,2,25]{m_{\star,e\circ\pi^{-1}}=m^\pi}{\rIE{30,\rAEb{29}}}
 \proofstepwidestar[1,2,25]{w=w_n(m_{\star,e\circ\pi^{-1}})}{\rIE{\FormulaRefAuto{a=b\vdash b=a}[thm]{31},\rAEb{25}}}
 \proofstepwidestar[1,2,25]{w\in W}{\FormulaRefAuto{FiniteSemigroupCanonicalTableCodeDef}[def]{1,\rEI{\rAI{\rAEa{29},32}}}}
 \proofstepwidestar[1,2,23]{w\in W}{\rEE{24,25,33}}
 \proofstepwidestar[1,2]{V\subseteq W}{\FormulaRefAuto{A\subseteq B\coloneq\forall x\,(x\in A\to x\in B)}[def]{\rUI{\rRI{23,34}}}}
 \proofstepwidestar[1,2]{W=V}{\FormulaRefAuto{A\subseteq B\dsep B\subseteq A\vdash A=B}[thm]{22,35}}
\end{tabproofwide}
Für \(n=0\) durchlaufen die beiden Mengen genau die leere Nummerierung
beziehungsweise die leere Permutation. Beide enthalten daher das eine
leere Tafeltupel; die Existenz eines Grundträgerelements wurde nicht benutzt.

\section{Endliche Kandidatenmenge und minimaler Vertreter}
\hypertarget{table-isomorphisms.proof.finite-candidates}{}
% Ergänzende Endlichkeitsbeweise; öffentliche Tafeldefinitionen bleiben in B37.
\FormulaThmDeltaK[Endliche kartesische Produkte]{%
  \Finite U\dsep\Finite V\vdash\Finite{U\times V}
}{B37TFiniteProduct}{\DeltaRow{Mengen}{U\dsep V\dsep z}}
\begin{tabproofwide}
  \proofstepwidestar[1]{\Finite U}{\rA}
  \proofstepwidestar[2]{\Finite V}{\rA}
  \proofstepwidestar[1,2]{\Finite{U\cup V}}{\FormulaRefAuto{FiniteUnion}[thm]{1,2}}
  \proofstepwidestar[1,2]{\Finite{\powerset(U\cup V)}}{\FormulaRefAuto{FinitePowerSet}[thm]{3}}
  \proofstepwidestar[1,2]{\Finite{\powerset(\powerset(U\cup V))}}{\FormulaRefAuto{FinitePowerSet}[thm]{4}}
  \proofstepwidestar[6]{z\in U\times V}{\rA}
  \proofstepwidestar[6]{\exists u\in U\,\exists v\in V\,z=(u,v)}{\FormulaRefAuto{x\in A\times B\eqvdash\exists a\in A\,\exists b\in B\,x=(a,b)}[thm]{6}}
  \proofstepwidestar[6]{z\in\powerset(\powerset(U\cup V))}{\FormulaRefAuto{\exists a\in A\,\exists b\in B\,(x=(a,b))\vdash x\in\powerset(\powerset(A\cup B))}[thm]{7}}
  \proofstepwidestar[]{U\times V\subseteq\powerset(\powerset(U\cup V))}{\FormulaRefAuto{A\subseteq B\coloneq\forall x\,(x\in A\to x\in B)}[def]{\rUI{\rRI{6,8}}}}
  \proofstepwidestar[1,2]{\Finite{U\times V}}{\FormulaRefAuto{FiniteSubset}[thm]{9,5}}
\end{tabproofwide}

Für einen endlichen Träger \(A\) setzen wir \(n=\FiniteCard(A)\) und
\(S=\NatSeg n\). Die Nummerierungen werden als Funktionsgraphen aufgefasst.
\Needspace{14\baselineskip}
\FormulaDefDeltaK[Menge der Nummerierungen]{%
 E_A\coloneqq\{e\in\powerset(S\times A)\mid e\colon S\bij A\}
}{B37TNumberingsDef}{\DeltaRow{Mengen}{A\dsep n\dsep S}\DeltaRow{Funktionen}{e}}
\Needspace{14\baselineskip}
\FormulaThmDeltaK[Die Nummerierungsmenge ist endlich und nichtleer]{%
 \Finite A\vdash\Finite{E_A}\land E_A\neq\varnothing
}{B37TFiniteNumberings}{\DeltaRow{Mengen}{A\dsep n\dsep S}\DeltaRow{Funktionen}{e}}
\begin{tabproofwide}
 \proofstepwidestar[1]{\Finite A}{\rA}
 \proofstepwidestar[1]{n\in\mathbb N}{\FormulaRefAuto{FiniteCardNatural}[thm]{1}}
 \proofstepwidestar[1]{\Finite S}{\FormulaRefAuto{PeanoNatSegFinite}[thm]{2}}
 \proofstepwidestar[1]{\Finite{S\times A}}{\FormulaRefAuto{B37TFiniteProduct}[thm]{3,1}}
 \proofstepwidestar[1]{\Finite{\powerset(S\times A)}}{\FormulaRefAuto{FinitePowerSet}[thm]{4}}
 \proofstepwidestar[]{E_A\subseteq\powerset(S\times A)}{\FormulaRefAuto{B37TNumberingsDef}[def]}
 \proofstepwidestar[1]{\Finite{E_A}}{\FormulaRefAuto{FiniteSubset}[thm]{6,5}}
 \proofstepwidestar[1]{S\EqCard A}{\FormulaRefAuto{FiniteCardCounts}[thm]{1}}
 \proofstepwidestar[1]{\exists e\,(e\colon S\bij A)}{\FormulaRefAuto{Gleichmächtigkeit}[def]{8}}
 \proofstepwidestar[10]{e\colon S\bij A}{\rA}
 \proofstepwidestar[10]{e\colon S\to A}{\FormulaRefAuto{Bijektive Funktion}[def]{10}}
 \proofstepwidestar[10]{e\subseteq S\times A}{\FormulaRefAuto{FunctionGraphSubsetProduct}[thm]{11}}
 \proofstepwidestar[10]{e\in\powerset(S\times A)}{\FormulaRefAuto{x\in\powerset(A)\eqvdash x\subseteq A}[thm]{12}}
 \proofstepwidestar[10]{e\in E_A}{\FormulaRefAuto{B37TNumberingsDef}[def]{13,10}}
 \proofstepwidestar[10]{\exists e\,(e\in E_A)}{\rEI{14}}
 \proofstepwidestar[1]{\exists e\,(e\in E_A)}{\rEE{9,10,15}}
 \proofstepwidestar[1]{E_A\neq\varnothing}{\FormulaRefAuto{A\neq\varnothing\eqvdash\exists x\,(x\in A)}[thm]{16}}
 \proofstepwidestar[1]{\Finite{E_A}\land E_A\neq\varnothing}{\rAI{7,17}}
\end{tabproofwide}

\Needspace{14\baselineskip}
\FormulaDefDeltaK[Kandidatenmenge des kanonischen Codes]{%
 \FiniteSemigroupAlg{A}{\star}\vdash
 W_A\coloneqq\{w_n(m_{\star,e})\mid e\in E_A\}
}{B37TCodeCandidatesDef}{\DeltaRow{Mengen}{A\dsep E_A\dsep W_A}\DeltaRow{Funktionen}{\star\dsep e}}
Dies ist dieselbe Menge wie in
\FormulaRefAuto{FiniteSemigroupCanonicalTableCodeDef}[def], denn jeder
bijektive Funktionsgraph gehört nach
\FormulaRefAuto{FunctionGraphSubsetProduct}[thm] zu \(\powerset(S\times A)\).
\Needspace{14\baselineskip}
\FormulaDefDeltaK[Abbildung von Nummerierungen auf Tafeltupel]{%
 \begin{aligned}[t]
 &\FiniteSemigroupAlg{A}{\star}\vdash
   \tau_A\colon E_A\to W_A,\\[-2pt]
 &\tau_A(e)\coloneqq w_n(m_{\star,e})\qquad(e\in E_A)
 \end{aligned}
}{B37TTableTupleMapDef}{\DeltaRow{Mengen}{A\dsep E_A\dsep W_A}\DeltaRow{Funktionen}{\star\dsep e\dsep\tau_A}}
\noindent\textbf{Zu zeigen: Existenz und Eindeutigkeit der Abbildung.}
Hier gelten die festgelegten Abkürzungen \(n=\FiniteCard(A)\) und
\(S=\NatSeg n\); \(E_A\) und \(W_A\) sind die oben definierten Mengen.
\[
\begin{aligned}
 &\FiniteSemigroupAlg{A}{\star}\vdash{}\\
 &\exists!F\,\Bigl(F\colon E_A\to W_A\land
 \forall e\in E_A\,F(e)=w_n(m_{\star,e})\Bigr).
\end{aligned}
\]
\begin{tabproofwide}
 \proofstepwidestar[1]{\FiniteSemigroupAlg{A}{\star}}{\rA}
 \proofstepwidestar[2]{e\in E_A}{\rA}
 \proofstepwidestar[1]{\Finite A}{\FormulaRefAuto{FiniteSemigroupFinitenessAxiom}[ax]{1}}
 \proofstepwidestar[1]{n\in\mathbb N}{\FormulaRefAuto{FiniteCardNatural}[thm]{3}}
 \proofstepwidestar[2]{e\colon S\bij A}{\FormulaRefAuto{B37TNumberingsDef}[def]{2}}
 \proofstepwidestar[1,2]{m_{\star,e}\colon S\times S\to S}{\FormulaRefAuto{FiniteSemigroupMultiplicationTableDef}[def]{1,4,5}}
 \proofstepwidestar[1,2]{w_n(m_{\star,e})\in W_A}{\FormulaRefAuto{B37TCodeCandidatesDef}[def]{1,2}}
 \proofstepwidestar[1]{\forall e\in E_A\,(w_n(m_{\star,e})\in W_A)}{\rUI{\rRI{2,7}}}
 \proofstepwidestarbreak[1]{\begin{aligned}[t]
 &\exists!F\,\bigl(F\colon E_A\to W_A\\[-2pt]
 &\qquad{}\land\forall e\in E_A\,F(e)=w_n(m_{\star,e})\bigr)
 \end{aligned}}{\FormulaRefAuto{TypedFunctionDefinitionPrinciple}[thm]{8}}
\end{tabproofwide}

\Needspace{14\baselineskip}
\FormulaThmDeltaK[Die Kandidatenmenge ist das Bild der Nummerierungen]{%
 \FiniteSemigroupAlg{A}{\star}\vdash W_A=\tau_A[E_A]
}{B37TCodeCandidateImage}{\DeltaRow{Mengen}{A\dsep E_A\dsep W_A\dsep z}\DeltaRow{Funktionen}{\star\dsep\tau_A}}
\begin{tabproofwide}
 \proofstepwidestar[1]{\FiniteSemigroupAlg{A}{\star}}{\rA}
 \proofstepwidestar[1]{\tau_A\colon E_A\to W_A}{\FormulaRefAuto{B37TTableTupleMapDef}[def]{1}}
 \proofstepwidestar[3]{z\in W_A}{\rA}
 \proofstepwidestar[1,3]{\exists e\in E_A\,(z=w_n(m_{\star,e}))}{\FormulaRefAuto{B37TCodeCandidatesDef}[def]{1,3}}
 \proofstepwidestar[5]{e\in E_A\land z=w_n(m_{\star,e})}{\rA}
 \proofstepwidestar[5]{e\in E_A}{\rAEa{5}}
 \proofstepwidestar[5]{z=w_n(m_{\star,e})}{\rAEb{5}}
 \proofstepwidestar[1,5]{\tau_A(e)=w_n(m_{\star,e})}{\FormulaRefAuto{B37TTableTupleMapDef}[def]{1,6}}
 \proofstepwidestar[1,5]{z=\tau_A(e)}{\rIE{8,7}}
 \proofstepwidestar[1,5]{\exists e\in E_A\,(z=\tau_A(e))}{\rEI{\rAI{6,9}}}
 \proofstepwidestar[1,3]{\exists e\in E_A\,(z=\tau_A(e))}{\rEE{4,5,10}}
 \proofstepwidestar[1,3]{z\in\tau_A[E_A]}{\FormulaRefAuto{F\colon A\to B\dsep\exists x\in A\,y=F(x)\vdash y\in F[A]}[thm]{2,11}}
 \proofstepwidestar[1]{W_A\subseteq\tau_A[E_A]}{\FormulaRefAuto{A\subseteq B\coloneq\forall x\,(x\in A\to x\in B)}[def]{\rUI{\rRI{3,12}}}}
 \proofstepwidestar[14]{z\in\tau_A[E_A]}{\rA}
 \proofstepwidestar[1,14]{\exists e\in E_A\,(z=\tau_A(e))}{\FormulaRefAuto{F\colon A\to B\dsep y\in F[A]\vdash\exists x\in A\,y=F(x)}[thm]{2,14}}
 \proofstepwidestar[16]{e\in E_A\land z=\tau_A(e)}{\rA}
 \proofstepwidestar[16]{e\in E_A}{\rAEa{16}}
 \proofstepwidestar[16]{z=\tau_A(e)}{\rAEb{16}}
 \proofstepwidestar[1,16]{\tau_A(e)=w_n(m_{\star,e})}{\FormulaRefAuto{B37TTableTupleMapDef}[def]{1,17}}
 \proofstepwidestar[1,16]{z=w_n(m_{\star,e})}{\rChain{18,19}}
 \proofstepwidestar[1,16]{\exists e\in E_A\,(z=w_n(m_{\star,e}))}{\rEI{\rAI{17,20}}}
 \proofstepwidestar[1,16]{z\in W_A}{\FormulaRefAuto{B37TCodeCandidatesDef}[def]{1,21}}
 \proofstepwidestar[1,14]{z\in W_A}{\rEE{15,16,22}}
 \proofstepwidestar[1]{\tau_A[E_A]\subseteq W_A}{\FormulaRefAuto{A\subseteq B\coloneq\forall x\,(x\in A\to x\in B)}[def]{\rUI{\rRI{14,23}}}}
 \proofstepwidestar[1]{W_A=\tau_A[E_A]}{\FormulaRefAuto{A\subseteq B\dsep B\subseteq A\vdash A=B}[thm]{13,24}}
\end{tabproofwide}

\Needspace{14\baselineskip}
\FormulaThmDeltaK[Die Kandidatenmenge des Codes ist endlich und nichtleer]{%
 \FiniteSemigroupAlg{A}{\star}\vdash\Finite{W_A}\land W_A\neq\varnothing
}{B37TFiniteCodeCandidates}{\DeltaRow{Mengen}{A\dsep E_A\dsep W_A}\DeltaRow{Funktionen}{\star\dsep\tau_A}}
\begin{tabproofwide}
 \proofstepwidestar[1]{\FiniteSemigroupAlg{A}{\star}}{\rA}
 \proofstepwidestar[1]{\Finite A}{\FormulaRefAuto{FiniteSemigroupFinitenessAxiom}[ax]{1}}
 \proofstepwidestar[1]{\Finite{E_A}\land E_A\neq\varnothing}{\FormulaRefAuto{B37TFiniteNumberings}[thm]{2}}
 \proofstepwidestar[1]{\tau_A\colon E_A\to W_A}{\FormulaRefAuto{B37TTableTupleMapDef}[def]{1}}
 \proofstepwidestar[]{E_A\subseteq E_A}{\FormulaRefAuto{A\subseteq A}[thm]}
 \proofstepwidestar[1]{\Finite{\tau_A[E_A]}}{\FormulaRefAuto{FiniteImage}[thm]{5,4,\rAEa{3}}}
 \proofstepwidestar[1]{W_A=\tau_A[E_A]}{\FormulaRefAuto{B37TCodeCandidateImage}[thm]{1}}
 \proofstepwidestar[1]{\Finite{W_A}}{\rIE{7,6}}
 \proofstepwidestar[1]{\exists e\,(e\in E_A)}{\FormulaRefAuto{A\neq\varnothing\eqvdash\exists x\,(x\in A)}[thm]{\rAEb{3}}}
 \proofstepwidestar[10]{e\in E_A}{\rA}
 \proofstepwidestar[1,10]{\tau_A(e)\in W_A}{\FormulaRefAuto{F\colon A\to B\dsep x\in A\vdash F(x)\in B}[thm]{4,10}}
 \proofstepwidestar[1,10]{\exists w\,(w\in W_A)}{\rEI{11}}
 \proofstepwidestar[1]{\exists w\,(w\in W_A)}{\rEE{9,10,12}}
 \proofstepwidestar[1]{W_A\neq\varnothing}{\FormulaRefAuto{A\neq\varnothing\eqvdash\exists x\,(x\in A)}[thm]{13}}
 \proofstepwidestar[1]{\Finite{W_A}\land W_A\neq\varnothing}{\rAI{8,14}}
\end{tabproofwide}

Für \(A=\varnothing\) ist \(n=0\), die einzige Nummerierung die leere
Abbildung und das einzige Tafeltupel das leere Tupel. Insbesondere ist
\(W_A\) auch in diesem Randfall eine Einermenge, keine leere Menge.

\subsection{Der kleinste Tafelcode}
\hypertarget{table-isomorphisms.proof.minimal-code}{}
% Zusätzliche Hilfssätze; öffentliche Tafelcodes bleiben in Band 37.
\paragraph{Die lexikographische Ordnung endlicher Tupel.}
Wir fassen ein Tupel der Länge \(r\) mit Werten in \(C\) als Funktion
auf \(I_r=\NatSeg r\) auf und schreiben \(x_i=x(i)\). Als Basis genügt
eine strikt total geordnete Wertemenge. Die folgenden Voraussetzungen
werden in den Hilfssätzen ausdrücklich durch \(\mathcal L\) mitgeführt.

\Needspace{14\baselineskip}
\FormulaDefDeltaK[Strikte Ordnungsannahmen für den Werteträger]{%
 \begin{aligned}[t]
 K(C,\triangleleft)\coloneqq{}&\bigl(\forall a\in C\,\neg(a\triangleleft a)\bigr)\\
 &{}\land\Bigl(\bigl(\forall a,b,c\in C\,
 ((a\triangleleft b\land b\triangleleft c)\to a\triangleleft c)\bigr)\\
 &\qquad{}\land\bigl(\forall a,b\in C\,
 (a\triangleleft b\lor(b\triangleleft a\lor a=b))\bigr)\Bigr)
 \end{aligned}
}{B37TLexBaseOrderDef}{\DeltaRow{Mengen}{C\dsep a\dsep b\dsep c}\DeltaRow{Relationsoperatoren}{\triangleleft}}
Wir schreiben in den folgenden generischen Nachweisen
\(\mathcal L\coloneqq r\in\mathbb N\land K(C,\triangleleft)\).
Der Werteträger und sein strikter Relationsoperator bleiben dabei fest;
bei der Anwendung werden beide ausdrücklich angegeben.

\Needspace{14\baselineskip}
\FormulaThmDeltaK[Irreflexivität auf dem Werteträger]{%
 K(C,\triangleleft)\dsep a\in C\vdash\neg(a\triangleleft a)
}{B37TLexBaseIrreflexive}{\DeltaRow{Mengen}{C\dsep a}}
\begin{tabproofwide}
 \proofstepwidestar[1]{K(C,\triangleleft)}{\rA}
 \proofstepwidestar[2]{a\in C}{\rA}
 \proofstepwidestar[1]{\forall a\in C\,\neg(a\triangleleft a)}{\rAEa{\FormulaRefAuto{B37TLexBaseOrderDef}[def]{1}}}
 \proofstepwidestar[1,2]{\neg(a\triangleleft a)}{\rRE{\rUE{3},2}}
\end{tabproofwide}
\Needspace{14\baselineskip}
\FormulaThmDeltaK[Transitivität auf dem Werteträger]{%
 K(C,\triangleleft)\dsep a,b,c\in C\dsep a\triangleleft b\dsep b\triangleleft c\vdash a\triangleleft c
}{B37TLexBaseTransitive}{\DeltaRow{Mengen}{C\dsep a\dsep b\dsep c}}
\begin{tabproofwide}
 \proofstepwidestar[1]{K(C,\triangleleft)}{\rA}
 \proofstepwidestar[2]{a\in C}{\rA}
 \proofstepwidestar[3]{b\in C}{\rA}
 \proofstepwidestar[4]{c\in C}{\rA}
 \proofstepwidestar[5]{a\triangleleft b}{\rA}
 \proofstepwidestar[6]{b\triangleleft c}{\rA}
 \proofstepwidestar[1]{\forall a,b,c\in C\,((a\triangleleft b\land b\triangleleft c)\to a\triangleleft c)}{\rAEa{\rAEb{\FormulaRefAuto{B37TLexBaseOrderDef}[def]{1}}}}
 \proofstepwidestar[1,2,3,4]{(a\triangleleft b\land b\triangleleft c)\to a\triangleleft c}{\rRE{\rUE{\rRE{\rUE{\rRE{\rUE{7},2}},3}},4}}
 \proofstepwidestar[1,2,3,4,5,6]{a\triangleleft c}{\rRE{8,\rAI{5,6}}}
\end{tabproofwide}
\Needspace{14\baselineskip}
\FormulaThmDeltaK[Vergleichbarkeit auf dem Werteträger]{%
 K(C,\triangleleft)\dsep a,b\in C\vdash a\triangleleft b\lor(b\triangleleft a\lor a=b)
}{B37TLexBaseTrichotomy}{\DeltaRow{Mengen}{C\dsep a\dsep b}}
\begin{tabproofwide}
 \proofstepwidestar[1]{K(C,\triangleleft)}{\rA}
 \proofstepwidestar[2]{a\in C}{\rA}
 \proofstepwidestar[3]{b\in C}{\rA}
 \proofstepwidestar[1]{\forall a,b\in C\,(a\triangleleft b\lor(b\triangleleft a\lor a=b))}{\rAEb{\rAEb{\FormulaRefAuto{B37TLexBaseOrderDef}[def]{1}}}}
 \proofstepwidestar[1,2,3]{a\triangleleft b\lor(b\triangleleft a\lor a=b)}{\rRE{\rUE{\rRE{\rUE{4},2}},3}}
\end{tabproofwide}

\Needspace{14\baselineskip}
\FormulaDefDeltaK[Tupel über einem Werteträger und gemeinsame Präfixe]{%
 \begin{aligned}[t]
 T_r(C)&\coloneqq\{x\in\powerset(I_r\times C)\mid x\colon I_r\to C\},\\
 P_r(x,y;k)&\coloneqq\forall i\in I_r\,(i<k\to x_i=y_i),\\
 D_r(x,y)&\coloneqq\{i\in I_r\mid x_i\ne y_i\}
 \end{aligned}
}{B37TLexTupleDef}{\DeltaRow{Mengen}{r\dsep I_r\dsep T_r(C)\dsep x\dsep y\dsep i\dsep k}}
\Needspace{14\baselineskip}
\FormulaDefDeltaK[Strikte und nichtstrikte lexikographische Ordnung]{%
 \begin{aligned}[t]
 &x,y\in T_r(C)\vdash{}\\[-2pt]
 &x\prec_{r,C} y\coloneqq\exists k\in I_r\,(P_r(x,y;k)\land x_k\triangleleft y_k),\\
 &x\preceq_{r,C} y\coloneqq x=y\lor x\prec_{r,C} y
 \end{aligned}
}{B37TLexOrderDef}{\DeltaRow{Mengen}{r\dsep x\dsep y\dsep k}\DeltaRow{Relationsoperatoren}{\prec_{r,C}\dsep\preceq_{r,C}}}

\Needspace{14\baselineskip}
\FormulaThmDeltaK[Typen von Indizes und Koordinaten]{%
 \begin{aligned}[t]
 &(r\in\mathbb N\land K(C,\triangleleft))\dsep x\in T_r(C)\dsep k\in I_r\vdash{}\\
 &k\in\mathbb N\land x_k\in C
 \end{aligned}
}{B37TLexCoordinate}{\DeltaRow{Mengen}{r\dsep x\dsep k}}
\begin{tabproofwide}
 \proofstepwidestar[1]{\mathcal L}{\rA}
 \proofstepwidestar[2]{x\in T_r(C)}{\rA}
 \proofstepwidestar[3]{k\in I_r}{\rA}
 \proofstepwidestar[2]{x\colon I_r\to C}{\FormulaRefAuto{B37TLexTupleDef}[def]{2}}
 \proofstepwidestar[1,3]{k\in\mathbb N}{\FormulaRefAuto{PeanoNatSegElementNatural}[thm]{\rAEa{1},3}}
 \proofstepwidestar[2,3]{x_k\in C}{\FormulaRefAuto{F\colon A\to B\dsep x\in A\vdash F(x)\in B}[thm]{4,3}}
 \proofstepwidestar[1,2,3]{k\in\mathbb N\land x_k\in C}{\rAI{5,6}}
\end{tabproofwide}

\Needspace{14\baselineskip}
\FormulaThmDeltaK[Die erste abweichende Stelle existiert]{%
 \begin{aligned}[t]
 &(r\in\mathbb N\land K(C,\triangleleft))\dsep x,y\in T_r(C)\dsep x\ne y\vdash{}\\
 &\exists k\in I_r\,(P_r(x,y;k)\land x_k\ne y_k)
 \end{aligned}
}{B37TLexFirstDifference}{\DeltaRow{Mengen}{r\dsep x\dsep y\dsep D=D_r(x,y)\dsep i\dsep k}}
\begin{tabproofwide}
 \proofstepwidestar[1]{\mathcal L}{\rA}
 \proofstepwidestar[2]{x,y\in T_r(C)}{\rA}
 \proofstepwidestar[3]{x\ne y}{\rA}
 \proofstepwidestar[]{D\subseteq I_r}{\FormulaRefAuto{B37TLexTupleDef}[def]}
 \proofstepwidestar[1]{I_r\subseteq\mathbb N}{\FormulaRefAuto{PeanoNatSegSubsetN}[thm]{\rAEa{1}}}
 \proofstepwidestar[1]{D\subseteq\mathbb N}{\FormulaRefAuto{A\subseteq B\dsep B\subseteq C\vdash A\subseteq C}[thm]{4,5}}
 \proofstepwidestar[7]{D=\varnothing}{\rA}
 \proofstepwidestar[8]{i\in I_r}{\rA}
 \proofstepwidestar[9]{x_i\ne y_i}{\rA}
 \proofstepwidestar[8,9]{i\in D}{\FormulaRefAuto{B37TLexTupleDef}[def]{8,9}}
 \proofstepwidestar[7,8,9]{i\in\varnothing}{\rIE{7,10}}
 \proofstepwidestar[]{i\not\in\varnothing}{\FormulaRefAuto{x\not\in\varnothing}[thm]}
 \proofstepwidestar[7,8,9]{\bot}{\rBI{11,12}}
 \proofstepwidestar[7,8]{x_i=y_i}{\rCE{9,13}}
 \proofstepwidestar[7]{\forall i\in I_r\,(x_i=y_i)}{\rUI{\rRI{8,14}}}
 \proofstepwidestar[2]{x\colon I_r\to C\land y\colon I_r\to C}{\FormulaRefAuto{B37TLexTupleDef}[def]{2}}
 \proofstepwidestar[2,7]{x=y}{\FormulaRefAuto{FunctionExtensionality}[thm]{\rAEa{16},\rAEb{16},15}}
 \proofstepwidestar[2,3,7]{\bot}{\rBI{3,17}}
 \proofstepwidestar[2,3]{D\ne\varnothing}{\rCI{7,18}}
 \proofstepwidestar[1,2,3]{\exists k\in D\,\forall i\in D\,(k\leq i)}{\FormulaRefAuto{NaturalWellOrderingPrinciple}[thm]{6,19}}
 \proofstepwidestar[21]{k\in D\land\forall i\in D\,(k\leq i)}{\rA}
 \proofstepwidestar[21]{k\in I_r\land x_k\ne y_k}{\FormulaRefAuto{B37TLexTupleDef}[def]{\rAEa{21}}}
 \proofstepwidestar[1,21]{k\in\mathbb N}{\FormulaRefAuto{PeanoNatSegElementNatural}[thm]{\rAEa{1},\rAEa{22}}}
 \proofstepwidestar[24]{i\in I_r}{\rA}
 \proofstepwidestar[25]{i<k}{\rA}
 \proofstepwidestar[26]{x_i\ne y_i}{\rA}
 \proofstepwidestar[24,26]{i\in D}{\FormulaRefAuto{B37TLexTupleDef}[def]{24,26}}
 \proofstepwidestar[21,24,26]{k\leq i}{\rRE{\rUE{\rAEb{21}},27}}
 \proofstepwidestar[1,24]{i\in\mathbb N}{\FormulaRefAuto{PeanoNatSegElementNatural}[thm]{\rAEa{1},24}}
 \proofstepwidestar[1,21,24,25,26]{\bot}{\FormulaRefAuto{PeanoOrderNoCycle}[thm]{29,23,25,28}}
 \proofstepwidestar[1,21,24,25]{x_i=y_i}{\rCE{26,30}}
 \proofstepwidestar[1,21]{P_r(x,y;k)}{\FormulaRefAuto{B37TLexTupleDef}[def]{\rUI{\rRI{24,\rRI{25,31}}}}}
 \proofstepwidestar[1,21]{\exists k\in I_r\,(P_r(x,y;k)\land x_k\ne y_k)}{\rEI{\rAI{\rAEa{22},\rAI{32,\rAEb{22}}}}}
 \proofstepwidestar[1,2,3]{\exists k\in I_r\,(P_r(x,y;k)\land x_k\ne y_k)}{\rEE{20,21,33}}
\end{tabproofwide}

\Needspace{14\baselineskip}
\FormulaThmDeltaK[Irreflexivität der strikten Tupelordnung]{%
 \begin{aligned}[t]
 &(r\in\mathbb N\land K(C,\triangleleft))\dsep x\in T_r(C)\vdash{}\\
 &\neg(x\prec_{r,C} x)
 \end{aligned}
}{B37TLexIrreflexive}{\DeltaRow{Mengen}{r\dsep x\dsep k}}
\begin{tabproofwide}
 \proofstepwidestar[1]{\mathcal L}{\rA}
 \proofstepwidestar[2]{x\in T_r(C)}{\rA}
 \proofstepwidestar[3]{x\prec_{r,C} x}{\rA}
 \proofstepwidestar[2,3]{\exists k\in I_r\,(P_r(x,x;k)\land x_k\triangleleft x_k)}{\FormulaRefAuto{B37TLexOrderDef}[def]{2,2,3}}
 \proofstepwidestar[5]{k\in I_r\land(P_r(x,x;k)\land x_k\triangleleft x_k)}{\rA}
 \proofstepwidestar[1,2,5]{x_k\in C}{\rAEb{\FormulaRefAuto{B37TLexCoordinate}[thm]{1,2,\rAEa{5}}}}
 \proofstepwidestar[1,2,5]{\neg(x_k\triangleleft x_k)}{\FormulaRefAuto{B37TLexBaseIrreflexive}[thm]{\rAEb{1},6}}
 \proofstepwidestar[1,2,5]{\bot}{\rBI{\rAEb{\rAEb{5}},7}}
 \proofstepwidestar[1,2,3]{\bot}{\rEE{4,5,8}}
 \proofstepwidestar[1,2]{\neg(x\prec_{r,C} x)}{\rCI{3,9}}
\end{tabproofwide}


\paragraph{Beschränkte Quantoren und gebündelte Voraussetzungen.}
Die folgenden beiden logischen Hilfssätze erlauben, die gemeinsam
notierten Trägerannahmen in den Ordnungsaxiomen ohne stillschweigende
Quantorenumformung einzusetzen.
\Needspace{14\baselineskip}
\FormulaThmDeltaK[Zwei beschränkte Quantoren aus einer gemeinsamen Trägerbedingung]{%
 \begin{aligned}[t]
 &\forall x\,\forall y\,((x\in A\land y\in A)\to R(x,y))\vdash{}\\
 &\forall x\in A\,\forall y\in A\,R(x,y)
 \end{aligned}
}{B37TLexBoundedPair}{\DeltaRow{Mengen}{A\dsep x\dsep y}\DeltaRow{Prädikat}{R}}
\begin{tabproofwide}
 \proofstepwidestar[1]{\forall x\,\forall y\,((x\in A\land y\in A)\to R(x,y))}{\rA}
 \proofstepwidestar[2]{x\in A}{\rA}
 \proofstepwidestar[3]{y\in A}{\rA}
 \proofstepwidestar[1]{(x\in A\land y\in A)\to R(x,y)}{\rUE{\rUE{1}}}
 \proofstepwidestar[1,2,3]{R(x,y)}{\rRE{4,\rAI{2,3}}}
 \proofstepwidestar[1,2]{\forall y\in A\,R(x,y)}{\rUI{\rRI{3,5}}}
 \proofstepwidestar[1]{\forall x\in A\,\forall y\in A\,R(x,y)}{\rUI{\rRI{2,6}}}
\end{tabproofwide}
\Needspace{14\baselineskip}
\FormulaThmDeltaK[Drei beschränkte Quantoren aus einer gemeinsamen Trägerbedingung]{%
 \begin{aligned}[t]
 &\forall x\,\forall y\,\forall z\,
 ((x\in A\land(y\in A\land z\in A))\to R(x,y,z))\vdash{}\\
 &\forall x\in A\,\forall y\in A\,\forall z\in A\,R(x,y,z)
 \end{aligned}
}{B37TLexBoundedTriple}{\DeltaRow{Mengen}{A\dsep x\dsep y\dsep z}\DeltaRow{Prädikat}{R}}
\begin{tabproofwide}
 \proofstepwidestar[1]{\forall x\,\forall y\,\forall z\,((x\in A\land(y\in A\land z\in A))\to R(x,y,z))}{\rA}
 \proofstepwidestar[2]{x\in A}{\rA}
 \proofstepwidestar[3]{y\in A}{\rA}
 \proofstepwidestar[4]{z\in A}{\rA}
 \proofstepwidestar[1]{(x\in A\land(y\in A\land z\in A))\to R(x,y,z)}{\rUE{\rUE{\rUE{1}}}}
 \proofstepwidestar[1,2,3,4]{R(x,y,z)}{\rRE{5,\rAI{2,\rAI{3,4}}}}
 \proofstepwidestar[1,2,3]{\forall z\in A\,R(x,y,z)}{\rUI{\rRI{4,6}}}
 \proofstepwidestar[1,2]{\forall y\in A\,\forall z\in A\,R(x,y,z)}{\rUI{\rRI{3,7}}}
 \proofstepwidestar[1]{\forall x\in A\,\forall y\in A\,\forall z\in A\,R(x,y,z)}{\rUI{\rRI{2,8}}}
\end{tabproofwide}

\Needspace{14\baselineskip}
\FormulaThmDeltaK[Vergleichbarkeit durch die erste abweichende Stelle]{%
 \begin{aligned}[t]
 &(r\in\mathbb N\land K(C,\triangleleft))\dsep x,y\in T_r(C)\vdash{}\\
 &x=y\lor(x\prec_{r,C} y\lor y\prec_{r,C} x)
 \end{aligned}
}{B37TLexTrichotomy}{\DeltaRow{Mengen}{r\dsep x\dsep y\dsep k}}
\begin{tabproofwide}
 \proofstepwidestar[1]{\mathcal L}{\rA}
 \proofstepwidestar[2]{x,y\in T_r(C)}{\rA}
 \proofstepwidestar[]{x=y\lor x\ne y}{\FormulaRefAuto{P\lor\neg P}[thm]}
 \proofstepwidestar[4]{x=y}{\rA}
 \proofstepwidestar[4]{x=y\lor(x\prec_{r,C} y\lor y\prec_{r,C} x)}{\rOIa{4}}
 \proofstepwidestar[6]{x\ne y}{\rA}
 \proofstepwidestar[1,2,6]{\exists k\in I_r\,(P_r(x,y;k)\land x_k\ne y_k)}{\FormulaRefAuto{B37TLexFirstDifference}[thm]{1,2,6}}
 \proofstepwidestar[8]{k\in I_r\land(P_r(x,y;k)\land x_k\ne y_k)}{\rA}
 \proofstepwidestar[1,2,8]{x_k\in C\land y_k\in C}{\rAI{\rAEb{\FormulaRefAuto{B37TLexCoordinate}[thm]{1,\rAEa{2},\rAEa{8}}},\rAEb{\FormulaRefAuto{B37TLexCoordinate}[thm]{1,\rAEb{2},\rAEa{8}}}}}
 \proofstepwidestar[1,2,8]{x_k\triangleleft y_k\lor(y_k\triangleleft x_k\lor x_k=y_k)}{\FormulaRefAuto{B37TLexBaseTrichotomy}[thm]{\rAEb{1},\rAEa{9},\rAEb{9}}}
 \proofstepwidestar[11]{x_k\triangleleft y_k}{\rA}
 \proofstepwidestar[2,8,11]{x\prec_{r,C} y}{\FormulaRefAuto{B37TLexOrderDef}[def]{2,\rEI{\rAI{\rAEa{8},\rAI{\rAEa{\rAEb{8}},11}}}}}
 \proofstepwidestar[2,8,11]{x=y\lor(x\prec_{r,C} y\lor y\prec_{r,C} x)}{\rOIb{\rOIa{12}}}
 \proofstepwidestar[14]{y_k\triangleleft x_k\lor x_k=y_k}{\rA}
 \proofstepwidestar[15]{y_k\triangleleft x_k}{\rA}
 \proofstepwidestar[16]{i\in I_r\land i<k}{\rA}
 \proofstepwidestar[8,16]{x_i=y_i}{\FormulaRefAuto{B37TLexTupleDef}[def]{\rAEa{\rAEb{8}},\rAEa{16},\rAEb{16}}}
 \proofstepwidestar[8,16]{y_i=x_i}{\FormulaRefAuto{a=b\vdash b=a}[thm]{17}}
 \proofstepwidestar[8]{P_r(y,x;k)}{\FormulaRefAuto{B37TLexTupleDef}[def]{\rUI{\FormulaRefAuto{(P \land Q) \rightarrow R \eqvdash P \rightarrow (Q \rightarrow R)}[thm]{\rRI{16,18}}}}}
 \proofstepwidestar[2,8,15]{y\prec_{r,C} x}{\FormulaRefAuto{B37TLexOrderDef}[def]{2,\rEI{\rAI{\rAEa{8},\rAI{19,15}}}}}
 \proofstepwidestar[2,8,15]{x=y\lor(x\prec_{r,C} y\lor y\prec_{r,C} x)}{\rOIb{\rOIb{20}}}
 \proofstepwidestar[22]{x_k=y_k}{\rA}
 \proofstepwidestar[8,22]{\bot}{\rBI{\rAEb{\rAEb{8}},22}}
 \proofstepwidestar[8,22]{x=y\lor(x\prec_{r,C} y\lor y\prec_{r,C} x)}{\FormulaRefAuto{\bot\vdash P}[thm]{23}}
 \proofstepwidestar[2,8,14]{x=y\lor(x\prec_{r,C} y\lor y\prec_{r,C} x)}{\rOE{14,15,21,22,24}}
 \proofstepwidestar[1,2,8]{x=y\lor(x\prec_{r,C} y\lor y\prec_{r,C} x)}{\rOE{10,11,13,14,25}}
 \proofstepwidestar[1,2,6]{x=y\lor(x\prec_{r,C} y\lor y\prec_{r,C} x)}{\rEE{7,8,26}}
 \proofstepwidestar[1,2]{x=y\lor(x\prec_{r,C} y\lor y\prec_{r,C} x)}{\rOE{3,4,5,6,27}}
\end{tabproofwide}

\paragraph{Transitivität: die drei Lagen der ersten Unterschiede.}
Für die nächsten drei Hilfssätze verwenden wir die reine Abkürzung
\[
\begin{aligned}
 H\coloneqq{}&\bigl(\mathcal L\land(x,y,z\in T_r(C)\land k,l\in I_r)\bigr)\\
 &{}\land\bigl(P_r(x,y;k)\land(P_r(y,z;l)\land(x_k\triangleleft y_k\land y_l\triangleleft z_l))\bigr).
\end{aligned}
\]
\Needspace{14\baselineskip}
\FormulaThmDeltaK[Der erste Unterschied liegt links früher]{%
 \begin{aligned}[t]
 &r\in\mathbb N\dsep K(C,\triangleleft)\dsep x,y,z\in T_r(C)
 \dsep k,l\in I_r\\
 &\dsep P_r(x,y;k)\dsep P_r(y,z;l)\dsep
 x_k\triangleleft y_k\dsep y_l\triangleleft z_l\\
 &\dsep k<l\vdash x\prec_{r,C} z
 \end{aligned}
}{B37TLexEarlyWitness}{\DeltaRow{Mengen}{r\dsep x\dsep y\dsep z\dsep i\dsep k\dsep l}}
\begin{tabproofwide}
 \proofstepwidestar[1]{H}{\rA}
 \proofstepwidestar[2]{k<l}{\rA}
 \proofstepwidestar[1]{\mathcal L\land(x,y,z\in T_r(C)\land k,l\in I_r)}{\rAEa{1}}
 \proofstepwidestar[1]{P_r(x,y;k)\land(P_r(y,z;l)\land(x_k\triangleleft y_k\land y_l\triangleleft z_l))}{\rAEb{1}}
 \proofstepwidestar[1]{k,l\in\mathbb N}{\rAI{\FormulaRefAuto{PeanoNatSegElementNatural}[thm]{\rAEa{\rAEa{3}},\rAEa{\rAEb{\rAEb{3}}}},\FormulaRefAuto{PeanoNatSegElementNatural}[thm]{\rAEa{\rAEa{3}},\rAEb{\rAEb{\rAEb{3}}}}}}
 \proofstepwidestar[6]{i\in I_r}{\rA}
 \proofstepwidestar[7]{i<k}{\rA}
 \proofstepwidestar[1,6]{i\in\mathbb N}{\FormulaRefAuto{PeanoNatSegElementNatural}[thm]{\rAEa{\rAEa{3}},6}}
 \proofstepwidestar[1,2,6,7]{i<l}{\FormulaRefAuto{PeanoLtTransitive}[thm]{8,\rAEa{5},\rAEb{5},7,2}}
 \proofstepwidestar[1,6,7]{x_i=y_i}{\FormulaRefAuto{B37TLexTupleDef}[def]{\rAEa{4},6,7}}
 \proofstepwidestar[1,2,6,7]{y_i=z_i}{\FormulaRefAuto{B37TLexTupleDef}[def]{\rAEa{\rAEb{4}},6,9}}
 \proofstepwidestar[1,2,6,7]{x_i=z_i}{\rIE{11,10}}
 \proofstepwidestar[1,2]{P_r(x,z;k)}{\FormulaRefAuto{B37TLexTupleDef}[def]{\rUI{\rRI{6,\rRI{7,12}}}}}
 \proofstepwidestar[1,2]{y_k=z_k}{\FormulaRefAuto{B37TLexTupleDef}[def]{\rAEa{\rAEb{4}},\rAEa{\rAEb{\rAEb{3}}},2}}
 \proofstepwidestar[1,2]{x_k\triangleleft z_k}{\rIE{14,\rAEa{\rAEb{\rAEb{4}}}}}
 \proofstepwidestar[1,2]{x\prec_{r,C} z}{\FormulaRefAuto{B37TLexOrderDef}[def]{\rAEa{\rAEb{3}},\rEI{\rAI{\rAEa{\rAEb{\rAEb{3}}},\rAI{13,15}}}}}
\end{tabproofwide}
\Needspace{14\baselineskip}
\FormulaThmDeltaK[Die ersten Unterschiede liegen an derselben Stelle]{%
 \begin{aligned}[t]
 &r\in\mathbb N\dsep K(C,\triangleleft)\dsep x,y,z\in T_r(C)
 \dsep k,l\in I_r\\
 &\dsep P_r(x,y;k)\dsep P_r(y,z;l)\dsep
 x_k\triangleleft y_k\dsep y_l\triangleleft z_l\\
 &\dsep k=l\vdash x\prec_{r,C} z
 \end{aligned}
}{B37TLexEqualWitness}{\DeltaRow{Mengen}{r\dsep x\dsep y\dsep z\dsep i\dsep k\dsep l}}
\begin{tabproofwide}
 \proofstepwidestar[1]{H}{\rA}
 \proofstepwidestar[2]{k=l}{\rA}
 \proofstepwidestar[1]{\mathcal L\land(x,y,z\in T_r(C)\land k,l\in I_r)}{\rAEa{1}}
 \proofstepwidestar[1,2]{P_r(x,y;k)\land(P_r(y,z;k)\land(x_k\triangleleft y_k\land y_k\triangleleft z_k))}{\rIE{\FormulaRefAuto{a=b\vdash b=a}[thm]{2},\rAEb{1}}}
 \proofstepwidestar[5]{i\in I_r}{\rA}
 \proofstepwidestar[6]{i<k}{\rA}
 \proofstepwidestar[1,2,5,6]{x_i=y_i\land y_i=z_i}{\rAI{\FormulaRefAuto{B37TLexTupleDef}[def]{\rAEa{4},5,6},\FormulaRefAuto{B37TLexTupleDef}[def]{\rAEa{\rAEb{4}},5,6}}}
 \proofstepwidestar[1,2,5,6]{x_i=z_i}{\rIE{\rAEb{7},\rAEa{7}}}
 \proofstepwidestar[1,2]{P_r(x,z;k)}{\FormulaRefAuto{B37TLexTupleDef}[def]{\rUI{\rRI{5,\rRI{6,8}}}}}
 \proofstepwidestar[1]{x_k,y_k,z_k\in C}{\rAI{\rAEb{\FormulaRefAuto{B37TLexCoordinate}[thm]{\rAEa{3},\rAEa{\rAEa{\rAEb{3}}},\rAEa{\rAEb{\rAEb{3}}}}},\rAI{\rAEb{\FormulaRefAuto{B37TLexCoordinate}[thm]{\rAEa{3},\rAEa{\rAEb{\rAEa{\rAEb{3}}}},\rAEa{\rAEb{\rAEb{3}}}}},\rAEb{\FormulaRefAuto{B37TLexCoordinate}[thm]{\rAEa{3},\rAEb{\rAEb{\rAEa{\rAEb{3}}}},\rAEa{\rAEb{\rAEb{3}}}}}}}}
 \proofstepwidestar[1,2]{x_k\triangleleft z_k}{\FormulaRefAuto{B37TLexBaseTransitive}[thm]{\rAEb{\rAEa{3}},\rAEa{10},\rAEa{\rAEb{10}},\rAEb{\rAEb{10}},\rAEa{\rAEb{\rAEb{4}}},\rAEb{\rAEb{\rAEb{4}}}}}
 \proofstepwidestar[1,2]{x\prec_{r,C} z}{\FormulaRefAuto{B37TLexOrderDef}[def]{\rAEa{\rAEb{3}},\rEI{\rAI{\rAEa{\rAEb{\rAEb{3}}},\rAI{9,11}}}}}
\end{tabproofwide}
\Needspace{14\baselineskip}
\FormulaThmDeltaK[Der erste Unterschied liegt rechts früher]{%
 \begin{aligned}[t]
 &r\in\mathbb N\dsep K(C,\triangleleft)\dsep x,y,z\in T_r(C)
 \dsep k,l\in I_r\\
 &\dsep P_r(x,y;k)\dsep P_r(y,z;l)\dsep
 x_k\triangleleft y_k\dsep y_l\triangleleft z_l\\
 &\dsep l<k\vdash x\prec_{r,C} z
 \end{aligned}
}{B37TLexLateWitness}{\DeltaRow{Mengen}{r\dsep x\dsep y\dsep z\dsep i\dsep k\dsep l}}
\begin{tabproofwide}
 \proofstepwidestar[1]{H}{\rA}
 \proofstepwidestar[2]{l<k}{\rA}
 \proofstepwidestar[1]{\mathcal L\land(x,y,z\in T_r(C)\land k,l\in I_r)}{\rAEa{1}}
 \proofstepwidestar[1]{P_r(x,y;k)\land(P_r(y,z;l)\land(x_k\triangleleft y_k\land y_l\triangleleft z_l))}{\rAEb{1}}
 \proofstepwidestar[1]{k,l\in\mathbb N}{\rAI{\FormulaRefAuto{PeanoNatSegElementNatural}[thm]{\rAEa{\rAEa{3}},\rAEa{\rAEb{\rAEb{3}}}},\FormulaRefAuto{PeanoNatSegElementNatural}[thm]{\rAEa{\rAEa{3}},\rAEb{\rAEb{\rAEb{3}}}}}}
 \proofstepwidestar[6]{i\in I_r}{\rA}
 \proofstepwidestar[7]{i<l}{\rA}
 \proofstepwidestar[1,6]{i\in\mathbb N}{\FormulaRefAuto{PeanoNatSegElementNatural}[thm]{\rAEa{\rAEa{3}},6}}
 \proofstepwidestar[1,2,6,7]{i<k}{\FormulaRefAuto{PeanoLtTransitive}[thm]{8,\rAEb{5},\rAEa{5},7,2}}
 \proofstepwidestar[1,2,6,7]{x_i=y_i}{\FormulaRefAuto{B37TLexTupleDef}[def]{\rAEa{4},6,9}}
 \proofstepwidestar[1,6,7]{y_i=z_i}{\FormulaRefAuto{B37TLexTupleDef}[def]{\rAEa{\rAEb{4}},6,7}}
 \proofstepwidestar[1,2,6,7]{x_i=z_i}{\rIE{11,10}}
 \proofstepwidestar[1,2]{P_r(x,z;l)}{\FormulaRefAuto{B37TLexTupleDef}[def]{\rUI{\rRI{6,\rRI{7,12}}}}}
 \proofstepwidestar[1,2]{x_l=y_l}{\FormulaRefAuto{B37TLexTupleDef}[def]{\rAEa{4},\rAEb{\rAEb{\rAEb{3}}},2}}
 \proofstepwidestar[1,2]{x_l\triangleleft z_l}{\rIE{\FormulaRefAuto{a=b\vdash b=a}[thm]{14},\rAEb{\rAEb{\rAEb{4}}}}}
 \proofstepwidestar[1,2]{x\prec_{r,C} z}{\FormulaRefAuto{B37TLexOrderDef}[def]{\rAEa{\rAEb{3}},\rEI{\rAI{\rAEb{\rAEb{\rAEb{3}}},\rAI{13,15}}}}}
\end{tabproofwide}

\Needspace{14\baselineskip}
\FormulaThmDeltaK[Transitivität der strikten Tupelordnung]{%
 \begin{aligned}[t]
 &(r\in\mathbb N\land K(C,\triangleleft))\dsep x,y,z\in T_r(C)\\ &\dsep x\prec_{r,C} y\dsep y\prec_{r,C} z\vdash{}\\
 &x\prec_{r,C} z
 \end{aligned}
}{B37TLexTransitive}{\DeltaRow{Mengen}{r\dsep x\dsep y\dsep z\dsep k\dsep l}}
\begin{tabproofwide}
 \proofstepwidestar[1]{\mathcal L}{\rA}
 \proofstepwidestar[2]{x,y,z\in T_r(C)}{\rA}
 \proofstepwidestar[3]{x\prec_{r,C} y}{\rA}
 \proofstepwidestar[4]{y\prec_{r,C} z}{\rA}
 \proofstepwidestar[2,3]{\exists k\in I_r\,(P_r(x,y;k)\land x_k\triangleleft y_k)}{\FormulaRefAuto{B37TLexOrderDef}[def]{2,3}}
 \proofstepwidestar[6]{k\in I_r\land(P_r(x,y;k)\land x_k\triangleleft y_k)}{\rA}
 \proofstepwidestar[2,4]{\exists l\in I_r\,(P_r(y,z;l)\land y_l\triangleleft z_l)}{\FormulaRefAuto{B37TLexOrderDef}[def]{2,4}}
 \proofstepwidestar[8]{l\in I_r\land(P_r(y,z;l)\land y_l\triangleleft z_l)}{\rA}
 \proofstepwidestar[1,2,6,8]{H}{\rAI{\rAI{1,\rAI{2,\rAI{\rAEa{6},\rAEa{8}}}},\rAI{\rAEa{\rAEb{6}},\rAI{\rAEa{\rAEb{8}},\rAI{\rAEb{\rAEb{6}},\rAEb{\rAEb{8}}}}}}}
 \proofstepwidestar[1,6,8]{k,l\in\mathbb N}{\rAI{\FormulaRefAuto{PeanoNatSegElementNatural}[thm]{\rAEa{1},\rAEa{6}},\FormulaRefAuto{PeanoNatSegElementNatural}[thm]{\rAEa{1},\rAEa{8}}}}
 \proofstepwidestar[1,6,8]{k<l\lor(l<k\lor k=l)}{\FormulaRefAuto{PeanoOrderTrichotomy}[thm]{\rAEa{10},\rAEb{10}}}
 \proofstepwidestar[12]{k<l}{\rA}
 \proofstepwidestar[1,2,6,8,12]{x\prec_{r,C} z}{\FormulaRefAuto{B37TLexEarlyWitness}[thm]{9,12}}
 \proofstepwidestar[14]{l<k\lor k=l}{\rA}
 \proofstepwidestar[15]{l<k}{\rA}
 \proofstepwidestar[1,2,6,8,15]{x\prec_{r,C} z}{\FormulaRefAuto{B37TLexLateWitness}[thm]{9,15}}
 \proofstepwidestar[17]{k=l}{\rA}
 \proofstepwidestar[1,2,6,8,17]{x\prec_{r,C} z}{\FormulaRefAuto{B37TLexEqualWitness}[thm]{9,17}}
 \proofstepwidestar[1,2,6,8,14]{x\prec_{r,C} z}{\rOE{14,15,16,17,18}}
 \proofstepwidestar[1,2,6,8]{x\prec_{r,C} z}{\rOE{11,12,13,14,19}}
 \proofstepwidestar[1,2,4,6]{x\prec_{r,C} z}{\rEE{7,8,20}}
 \proofstepwidestar[1,2,3,4]{x\prec_{r,C} z}{\rEE{5,6,21}}
\end{tabproofwide}

\Needspace{14\baselineskip}
\FormulaThmDeltaK[Asymmetrie der strikten Tupelordnung]{%
 \begin{aligned}[t]
 &(r\in\mathbb N\land K(C,\triangleleft))\dsep x,y\in T_r(C)\\ &\dsep x\prec_{r,C} y\vdash{}\\
 &\neg(y\prec_{r,C} x)
 \end{aligned}
}{B37TLexAsymmetric}{\DeltaRow{Mengen}{r\dsep x\dsep y}}
\begin{tabproofwide}
 \proofstepwidestar[1]{\mathcal L}{\rA}
 \proofstepwidestar[2]{x,y\in T_r(C)}{\rA}
 \proofstepwidestar[3]{x\prec_{r,C} y}{\rA}
 \proofstepwidestar[4]{y\prec_{r,C} x}{\rA}
 \proofstepwidestar[1,2,3,4]{x\prec_{r,C} x}{\FormulaRefAuto{B37TLexTransitive}[thm]{1,\rAI{\rAEa{2},\rAI{\rAEb{2},\rAEa{2}}},3,4}}
 \proofstepwidestar[1,2]{\neg(x\prec_{r,C} x)}{\FormulaRefAuto{B37TLexIrreflexive}[thm]{1,\rAEa{2}}}
 \proofstepwidestar[1,2,3,4]{\bot}{\rBI{5,6}}
 \proofstepwidestar[1,2,3]{\neg(y\prec_{r,C} x)}{\rCI{4,7}}
\end{tabproofwide}

\Needspace{14\baselineskip}
\FormulaThmDeltaK[Reflexivität der nichtstrikten Tupelordnung]{%
 \begin{aligned}[t]
 &(r\in\mathbb N\land K(C,\triangleleft))\dsep x\in T_r(C)\vdash{}\\
 &x\preceq_{r,C} x
 \end{aligned}
}{B37TLexWeakReflexive}{\DeltaRow{Mengen}{r\dsep x}}
\begin{tabproofwide}
 \proofstepwidestar[1]{\mathcal L}{\rA}
 \proofstepwidestar[2]{x\in T_r(C)}{\rA}
 \proofstepwidestar[]{x=x}{\rII}
 \proofstepwidestar[]{x=x\lor x\prec_{r,C} x}{\rOIa{3}}
 \proofstepwidestar[2]{x\preceq_{r,C} x}{\FormulaRefAuto{B37TLexOrderDef}[def]{2,2,4}}
\end{tabproofwide}
\Needspace{14\baselineskip}
\FormulaThmDeltaK[Antisymmetrie der nichtstrikten Tupelordnung]{%
 \begin{aligned}[t]
 &(r\in\mathbb N\land K(C,\triangleleft))\dsep x,y\in T_r(C)\\ &\dsep x\preceq_{r,C} y\dsep y\preceq_{r,C} x\vdash{}\\
 &x=y
 \end{aligned}
}{B37TLexWeakAntisymmetric}{\DeltaRow{Mengen}{r\dsep x\dsep y}}
\begin{tabproofwide}
 \proofstepwidestar[1]{\mathcal L}{\rA}
 \proofstepwidestar[2]{x,y\in T_r(C)}{\rA}
 \proofstepwidestar[3]{x\preceq_{r,C} y}{\rA}
 \proofstepwidestar[4]{y\preceq_{r,C} x}{\rA}
 \proofstepwidestar[2,3]{x=y\lor x\prec_{r,C} y}{\FormulaRefAuto{B37TLexOrderDef}[def]{2,3}}
 \proofstepwidestar[6]{x=y}{\rA}
 \proofstepwidestar[7]{x\prec_{r,C} y}{\rA}
 \proofstepwidestar[2,4]{y=x\lor y\prec_{r,C} x}{\FormulaRefAuto{B37TLexOrderDef}[def]{2,4}}
 \proofstepwidestar[9]{y=x}{\rA}
 \proofstepwidestar[9]{x=y}{\FormulaRefAuto{a=b\vdash b=a}[thm]{9}}
 \proofstepwidestar[11]{y\prec_{r,C} x}{\rA}
 \proofstepwidestar[1,2,7]{\neg(y\prec_{r,C} x)}{\FormulaRefAuto{B37TLexAsymmetric}[thm]{1,2,7}}
 \proofstepwidestar[1,2,7,11]{\bot}{\rBI{11,12}}
 \proofstepwidestar[1,2,7,11]{x=y}{\FormulaRefAuto{\bot\vdash P}[thm]{13}}
 \proofstepwidestar[1,2,4,7]{x=y}{\rOE{8,9,10,11,14}}
 \proofstepwidestar[1,2,3,4]{x=y}{\rOE{5,6,6,7,15}}
\end{tabproofwide}
\Needspace{14\baselineskip}
\FormulaThmDeltaK[Transitivität der nichtstrikten Tupelordnung]{%
 \begin{aligned}[t]
 &(r\in\mathbb N\land K(C,\triangleleft))\dsep x,y,z\in T_r(C)\\ &\dsep x\preceq_{r,C} y\dsep y\preceq_{r,C} z\vdash{}\\
 &x\preceq_{r,C} z
 \end{aligned}
}{B37TLexWeakTransitive}{\DeltaRow{Mengen}{r\dsep x\dsep y\dsep z}}
\begin{tabproofwide}
 \proofstepwidestar[1]{\mathcal L}{\rA}
 \proofstepwidestar[2]{x,y,z\in T_r(C)}{\rA}
 \proofstepwidestar[3]{x\preceq_{r,C} y}{\rA}
 \proofstepwidestar[4]{y\preceq_{r,C} z}{\rA}
 \proofstepwidestar[2,3]{x=y\lor x\prec_{r,C} y}{\FormulaRefAuto{B37TLexOrderDef}[def]{2,3}}
 \proofstepwidestar[6]{x=y}{\rA}
 \proofstepwidestar[4,6]{x\preceq_{r,C} z}{\rIE{\FormulaRefAuto{a=b\vdash b=a}[thm]{6},4}}
 \proofstepwidestar[8]{x\prec_{r,C} y}{\rA}
 \proofstepwidestar[2,4]{y=z\lor y\prec_{r,C} z}{\FormulaRefAuto{B37TLexOrderDef}[def]{2,4}}
 \proofstepwidestar[10]{y=z}{\rA}
 \proofstepwidestar[8,10]{x\prec_{r,C} z}{\rIE{10,8}}
 \proofstepwidestar[2,8,10]{x\preceq_{r,C} z}{\FormulaRefAuto{B37TLexOrderDef}[def]{2,\rOIb{11}}}
 \proofstepwidestar[13]{y\prec_{r,C} z}{\rA}
 \proofstepwidestar[1,2,8,13]{x\prec_{r,C} z}{\FormulaRefAuto{B37TLexTransitive}[thm]{1,2,8,13}}
 \proofstepwidestar[1,2,8,13]{x\preceq_{r,C} z}{\FormulaRefAuto{B37TLexOrderDef}[def]{2,\rOIb{14}}}
 \proofstepwidestar[1,2,4,8]{x\preceq_{r,C} z}{\rOE{9,10,12,13,15}}
 \proofstepwidestar[1,2,3,4]{x\preceq_{r,C} z}{\rOE{5,6,7,8,16}}
\end{tabproofwide}

\Needspace{14\baselineskip}
\FormulaThmDeltaK[Die lexikographische Tupelordnung ist total]{%
 (r\in\mathbb N\land K(C,\triangleleft))\vdash\TotOrd{T_r(C),\preceq_{r,C}}
}{B37TLexTotalOrder}{\DeltaRow{Mengen}{r\dsep x\dsep y\dsep z}\DeltaRow{Relationsoperatoren}{\preceq_{r,C}}}
\begin{tabproofwide}
 \proofstepwidestar[1]{\mathcal L}{\rA}
 \proofstepwidestar[2]{x\in T_r(C)}{\rA}
 \proofstepwidestar[1,2]{x\preceq_{r,C} x}{\FormulaRefAuto{B37TLexWeakReflexive}[thm]{1,2}}
 \proofstepwidestar[1]{\forall x\in T_r(C)\,(x\preceq_{r,C} x)}{\rUI{\rRI{2,3}}}
 \proofstepwidestar[5]{x,y,z\in T_r(C)\land(x\preceq_{r,C} y\land y\preceq_{r,C} z)}{\rA}
 \proofstepwidestar[1,5]{x\preceq_{r,C} z}{\FormulaRefAuto{B37TLexWeakTransitive}[thm]{1,\rAEa{5},\rAEa{\rAEb{5}},\rAEb{\rAEb{5}}}}
 \proofstepwidestar[1]{\forall x,y,z\in T_r(C)\,((x\preceq_{r,C} y\land y\preceq_{r,C} z)\to x\preceq_{r,C} z)}{\FormulaRefAuto{B37TLexBoundedTriple}[thm]{\rUI{\rUI{\rUI{\FormulaRefAuto{(P \land Q) \rightarrow R \eqvdash P \rightarrow (Q \rightarrow R)}[thm]{\rRI{5,6}}}}}}}
 \proofstepwidestar[8]{x,y\in T_r(C)\land(x\preceq_{r,C} y\land y\preceq_{r,C} x)}{\rA}
 \proofstepwidestar[1,8]{x=y}{\FormulaRefAuto{B37TLexWeakAntisymmetric}[thm]{1,\rAEa{8},\rAEa{\rAEb{8}},\rAEb{\rAEb{8}}}}
 \proofstepwidestar[1]{\forall x,y\in T_r(C)\,((x\preceq_{r,C} y\land y\preceq_{r,C} x)\to x=y)}{\FormulaRefAuto{B37TLexBoundedPair}[thm]{\rUI{\rUI{\FormulaRefAuto{(P \land Q) \rightarrow R \eqvdash P \rightarrow (Q \rightarrow R)}[thm]{\rRI{8,9}}}}}}
 \proofstepwidestar[1]{\PartOrd{T_r(C),\preceq_{r,C}}}{\FormulaRefAuto{Partielle Ordnung}[def]{4,7,10}}
 \proofstepwidestar[12]{x,y\in T_r(C)}{\rA}
 \proofstepwidestar[1,12]{x=y\lor(x\prec_{r,C} y\lor y\prec_{r,C} x)}{\FormulaRefAuto{B37TLexTrichotomy}[thm]{1,12}}
 \proofstepwidestar[14]{x=y}{\rA}
 \proofstepwidestar[12,14]{x\preceq_{r,C} y\lor y\preceq_{r,C} x}{\rOIa{\FormulaRefAuto{B37TLexOrderDef}[def]{12,\rOIa{14}}}}
 \proofstepwidestar[16]{x\prec_{r,C} y\lor y\prec_{r,C} x}{\rA}
 \proofstepwidestar[17]{x\prec_{r,C} y}{\rA}
 \proofstepwidestar[12,17]{x\preceq_{r,C} y\lor y\preceq_{r,C} x}{\rOIa{\FormulaRefAuto{B37TLexOrderDef}[def]{12,\rOIb{17}}}}
 \proofstepwidestar[19]{y\prec_{r,C} x}{\rA}
 \proofstepwidestar[12,19]{x\preceq_{r,C} y\lor y\preceq_{r,C} x}{\rOIb{\FormulaRefAuto{B37TLexOrderDef}[def]{12,\rOIb{19}}}}
 \proofstepwidestar[12,16]{x\preceq_{r,C} y\lor y\preceq_{r,C} x}{\rOE{16,17,18,19,20}}
 \proofstepwidestar[1,12]{x\preceq_{r,C} y\lor y\preceq_{r,C} x}{\rOE{13,14,15,16,21}}
 \proofstepwidestar[1]{\forall x,y\in T_r(C)\,(x\preceq_{r,C} y\lor y\preceq_{r,C} x)}{\FormulaRefAuto{B37TLexBoundedPair}[thm]{\rUI{\rUI{\rRI{12,22}}}}}
 \proofstepwidestar[1]{\TotOrd{T_r(C),\preceq_{r,C}}}{\FormulaRefAuto{Totale Ordnung}[def]{11,23}}
\end{tabproofwide}

\paragraph{Existenz und Eindeutigkeit des kleinsten Kandidaten.}
\Needspace{14\baselineskip}
\FormulaThmDeltaK[Endliche nichtleere Tupelmengen besitzen genau ein Minimum]{%
 \begin{aligned}[t]
 &(r\in\mathbb N\land K(C,\triangleleft))\dsep W\subseteq T_r(C)\dsep\Finite W\dsep W\ne\varnothing\vdash{}\\
 &\exists!q\in W\,\forall w\in W\,(q\preceq_{r,C} w)
 \end{aligned}
}{B37TLexFiniteMinimum}{\DeltaRow{Mengen}{r\dsep W\dsep q\dsep w}}
\begin{tabproofwide}
 \proofstepwidestar[1]{\mathcal L}{\rA}
 \proofstepwidestar[2]{W\subseteq T_r(C)}{\rA}
 \proofstepwidestar[3]{\Finite W}{\rA}
 \proofstepwidestar[4]{W\ne\varnothing}{\rA}
 \proofstepwidestar[1]{\TotOrd{T_r(C),\preceq_{r,C}}}{\FormulaRefAuto{B37TLexTotalOrder}[thm]{1}}
 \proofstepwidestar[1]{\PartOrd{T_r(C),\preceq_{r,C}}}{\FormulaRefAuto{Totale Ordnung}[def]{5}}
 \proofstepwidestar[1,2,3,4]{\exists q\in W\,\forall w\in W\,(w\preceq_{r,C} q\to w=q)}{\FormulaRefAuto{FinitePosetMinimalElement}[thm]{6,3,4,2}}
 \proofstepwidestar[8]{q\in W\land\forall w\in W\,(w\preceq_{r,C} q\to w=q)}{\rA}
 \proofstepwidestar[2,8]{\operatorname{MinEl}_{T_r(C),\preceq_{r,C}}(q,W)}{\FormulaRefAuto{OrderMinimalElement}[def]{2,8}}
 \proofstepwidestar[1,2,8]{q=\min_{T_r(C),\preceq_{r,C}}(W)}{\FormulaRefAuto{TotalOrderMinimalElementIsMinimum}[thm]{5,2,9}}
 \proofstepwidestar[2,8]{q\in T_r(C)}{\FormulaRefAuto{A\subseteq B,\,x\in A\vdash x\in B}[thm]{2,\rAEa{8}}}
 \proofstepwidestar[12]{w\in W}{\rA}
 \proofstepwidestar[2,12]{w\in T_r(C)}{\FormulaRefAuto{A\subseteq B,\,x\in A\vdash x\in B}[thm]{2,12}}
 \proofstepwidestar[1,2,8,12]{q\preceq_{r,C} w\lor w\preceq_{r,C} q}{\FormulaRefAuto{OrderComparability}[ax]{5,11,13}}
 \proofstepwidestar[15]{q\preceq_{r,C} w}{\rA}
 \proofstepwidestar[16]{w\preceq_{r,C} q}{\rA}
 \proofstepwidestar[8,12,16]{w=q}{\rRE{\rRE{\rUE{\rAEb{8}},12},16}}
 \proofstepwidestar[1,2,8]{q\preceq_{r,C} q}{\FormulaRefAuto{B37TLexWeakReflexive}[thm]{1,11}}
 \proofstepwidestar[1,2,8,12,16]{q\preceq_{r,C} w}{\rIE{\FormulaRefAuto{a=b\vdash b=a}[thm]{17},18}}
 \proofstepwidestar[1,2,8,12]{q\preceq_{r,C} w}{\rOE{14,15,15,16,19}}
 \proofstepwidestar[1,2,8]{\forall w\in W\,(q\preceq_{r,C} w)}{\rUI{\rRI{12,20}}}
 \proofstepwidestar[1,2,8]{\exists q\in W\,\forall w\in W\,(q\preceq_{r,C} w)}{\rEI{\rAI{\rAEa{8},21}}}
 \proofstepwidestar[1,2,3,4]{\exists q\in W\,\forall w\in W\,(q\preceq_{r,C} w)}{\rEE{7,8,22}}
 \proofstepwidestar[1,2,3,4]{\exists!q\in W\,\forall w\in W\,(q\preceq_{r,C} w)}{\FormulaRefAuto{MinimumUnique}[thm]{6,2,23}}
\end{tabproofwide}

\Needspace{14\baselineskip}
\FormulaThmDeltaK[Der Randfall der Tupellänge null]{%
 x,y\in T_0(C)\vdash x=y
}{B37TLexEmptyLength}{\DeltaRow{Mengen}{x\dsep y\dsep i}}
\begin{tabproofwide}
 \proofstepwidestar[1]{x,y\in T_0(C)}{\rA}
 \proofstepwidestar[1]{x\colon I_0\to C\land y\colon I_0\to C}{\FormulaRefAuto{B37TLexTupleDef}[def]{1}}
 \proofstepwidestar[]{I_0=\varnothing}{\FormulaRefAuto{PeanoNatSegZero}[thm]}
 \proofstepwidestar[4]{i\in I_0}{\rA}
 \proofstepwidestar[4]{i\in\varnothing}{\rIE{3,4}}
 \proofstepwidestar[]{i\not\in\varnothing}{\FormulaRefAuto{x\not\in\varnothing}[thm]}
 \proofstepwidestar[4]{\bot}{\rBI{5,6}}
 \proofstepwidestar[4]{x_i=y_i}{\FormulaRefAuto{\bot\vdash P}[thm]{7}}
 \proofstepwidestar[]{\forall i\in I_0\,(x_i=y_i)}{\rUI{\rRI{4,8}}}
 \proofstepwidestar[1]{x=y}{\FormulaRefAuto{FunctionExtensionality}[thm]{\rAEa{2},\rAEb{2},9}}
\end{tabproofwide}
\Needspace{5\baselineskip}
Insbesondere bedarf die Minimumskonstruktion keiner Voraussetzung
\mbox{\(r>0\)}. In einer nichtleeren Kandidatenmenge der Länge null sind
alle Kandidaten gleich dem leeren Tupel.


\Needspace{14\baselineskip}
\FormulaThmDeltaK[Die natürlichen Zahlen erfüllen die strikten Ordnungsannahmen]{%
 K(\mathbb N,<)
}{B37TLexNaturalBase}{\DeltaRow{Mengen}{a\dsep b\dsep c}}
\begin{tabproofwide}
 \proofstepwidestar[1]{a\in\mathbb N}{\rA}
 \proofstepwidestar[1]{a\leq a}{\FormulaRefAuto{PeanoLeqReflexive}[thm]{1}}
 \proofstepwidestar[3]{a<a}{\rA}
 \proofstepwidestar[1,3]{\bot}{\FormulaRefAuto{PeanoOrderNoCycle}[thm]{1,1,3,2}}
 \proofstepwidestar[1]{\neg(a<a)}{\rCI{3,4}}
 \proofstepwidestar[]{\forall a\in\mathbb N\,\neg(a<a)}{\rUI{\rRI{1,5}}}
 \proofstepwidestar[7]{a,b,c\in\mathbb N\land(a<b\land b<c)}{\rA}
 \proofstepwidestar[7]{a<c}{\FormulaRefAuto{PeanoLtTransitive}[thm]{\rAEa{\rAEa{7}},\rAEa{\rAEb{\rAEa{7}}},\rAEb{\rAEb{\rAEa{7}}},\rAEa{\rAEb{7}},\rAEb{\rAEb{7}}}}
 \proofstepwidestar[]{\forall a,b,c\in\mathbb N\,((a<b\land b<c)\to a<c)}{\FormulaRefAuto{B37TLexBoundedTriple}[thm]{\rUI{\rUI{\rUI{\FormulaRefAuto{(P \land Q) \rightarrow R \eqvdash P \rightarrow (Q \rightarrow R)}[thm]{\rRI{7,8}}}}}}}
 \proofstepwidestar[10]{a,b\in\mathbb N}{\rA}
 \proofstepwidestar[10]{a<b\lor(b<a\lor a=b)}{\FormulaRefAuto{PeanoOrderTrichotomy}[thm]{\rAEa{10},\rAEb{10}}}
 \proofstepwidestar[]{\forall a,b\in\mathbb N\,(a<b\lor(b<a\lor a=b))}{\FormulaRefAuto{B37TLexBoundedPair}[thm]{\rUI{\rUI{\rRI{10,11}}}}}
 \proofstepwidestar[]{K(\mathbb N,<)}{\FormulaRefAuto{B37TLexBaseOrderDef}[def]{\rAI{6,\rAI{9,12}}}}
\end{tabproofwide}

\Needspace{14\baselineskip}
\FormulaThmDeltaK[Die strikten Ordnungsannahmen übertragen sich auf Tupel]{%
 (r\in\mathbb N\land K(C,\triangleleft))\vdash K(T_r(C),\prec_{r,C})
}{B37TLexIterationBase}{\DeltaRow{Mengen}{r\dsep C\dsep x\dsep y\dsep z}}
\begin{tabproofwide}
 \proofstepwidestar[1]{\mathcal L}{\rA}
 \proofstepwidestar[2]{x\in T_r(C)}{\rA}
 \proofstepwidestar[1,2]{\neg(x\prec_{r,C}x)}{\FormulaRefAuto{B37TLexIrreflexive}[thm]{1,2}}
 \proofstepwidestar[1]{\forall x\in T_r(C)\,\neg(x\prec_{r,C}x)}{\rUI{\rRI{2,3}}}
 \proofstepwidestar[5]{x,y,z\in T_r(C)\land(x\prec_{r,C}y\land y\prec_{r,C}z)}{\rA}
 \proofstepwidestar[1,5]{x\prec_{r,C}z}{\FormulaRefAuto{B37TLexTransitive}[thm]{1,\rAEa{5},\rAEa{\rAEb{5}},\rAEb{\rAEb{5}}}}
 \proofstepwidestar[1]{\forall x,y,z\in T_r(C)\,((x\prec_{r,C}y\land y\prec_{r,C}z)\to x\prec_{r,C}z)}{\FormulaRefAuto{B37TLexBoundedTriple}[thm]{\rUI{\rUI{\rUI{\FormulaRefAuto{(P \land Q) \rightarrow R \eqvdash P \rightarrow (Q \rightarrow R)}[thm]{\rRI{5,6}}}}}}}
 \proofstepwidestar[8]{x,y\in T_r(C)}{\rA}
 \proofstepwidestar[1,8]{x=y\lor(x\prec_{r,C}y\lor y\prec_{r,C}x)}{\FormulaRefAuto{B37TLexTrichotomy}[thm]{1,8}}
 \proofstepwidestar[10]{x=y}{\rA}
 \proofstepwidestar[10]{x\prec_{r,C}y\lor(y\prec_{r,C}x\lor x=y)}{\rOIb{\rOIb{10}}}
 \proofstepwidestar[12]{x\prec_{r,C}y\lor y\prec_{r,C}x}{\rA}
 \proofstepwidestar[13]{x\prec_{r,C}y}{\rA}
 \proofstepwidestar[13]{x\prec_{r,C}y\lor(y\prec_{r,C}x\lor x=y)}{\rOIa{13}}
 \proofstepwidestar[15]{y\prec_{r,C}x}{\rA}
 \proofstepwidestar[15]{x\prec_{r,C}y\lor(y\prec_{r,C}x\lor x=y)}{\rOIb{\rOIa{15}}}
 \proofstepwidestar[12]{x\prec_{r,C}y\lor(y\prec_{r,C}x\lor x=y)}{\rOE{12,13,14,15,16}}
 \proofstepwidestar[1,8]{x\prec_{r,C}y\lor(y\prec_{r,C}x\lor x=y)}{\rOE{9,10,11,12,17}}
 \proofstepwidestar[1]{\forall x,y\in T_r(C)\,(x\prec_{r,C}y\lor(y\prec_{r,C}x\lor x=y))}{\FormulaRefAuto{B37TLexBoundedPair}[thm]{\rUI{\rUI{\rRI{8,18}}}}}
 \proofstepwidestar[1]{K(T_r(C),\prec_{r,C})}{\FormulaRefAuto{B37TLexBaseOrderDef}[def]{\rAI{4,\rAI{7,19}}}}
\end{tabproofwide}

\paragraph{Anwendung auf die Tafelkandidaten.}
Wir ordnen die verschachtelten Tafeltupel direkt. Setze
\(U_n=T_n(\mathbb N)\) und \(V_n=T_n(U_n)\). Auf \(U_n\) wird die
lexikographische Ordnung der natürlichen Koordinaten verwendet, auf
\(V_n\) anschließend die lexikographische Ordnung dieser Zeilen.
\(W_A\) bleibt die zweite Komponente des öffentlichen Codes;
\(n_A=\FiniteCard(A)\).
% Typisierung ohne Identifikation verschachtelter und flacher Tupel.
\Needspace{14\baselineskip}
\FormulaThmDeltaK[Tafelkandidaten sind verschachtelte natürliche Tupel]{%
 \FiniteSemigroupAlg{A}{\star}\vdash
 n_A\in\mathbb N\land W_A\subseteq V_{n_A}
}{B37TLexCodeCandidatesTyped}{%
 \DeltaRow{Mengen}{A\dsep n_A\dsep W_A\dsep w\dsep i\dsep j}
 \DeltaRow{Funktionen}{e\dsep\star\dsep m\dsep\ell_i}
}
Hier sind \(n_A=\FiniteCard(A)\), \(I=I_{n_A}\),
\(U=T_{n_A}(\mathbb N)\) und \(V_{n_A}=T_{n_A}(U)\).
Im Zweig mit der Nummerierung \(e\) kürzen wir
\(m=m_{\star,e}\) und \(\ell_i=(m(i,j))_{j\in I}\) ab.
Die innere Folge und danach die äußere Folge werden als die durch ihre
Vorschriften eindeutig bestimmten Funktionen gelesen. In beiden Fällen
wird zuerst jeder Wert typisiert; anschließend rechtfertigt
\FormulaRefAuto{TypedFunctionDefinitionPrinciple}[thm] die in
\FormulaRefAuto{FiniteSemigroupTableTupleDef}[def] verwendete Funktion.
Der dort definierte verschachtelte Kandidat bleibt dabei unverändert.
\begin{tabproofwide}
 \proofstepwidestar[1]{\FiniteSemigroupAlg{A}{\star}}{\rA}
 \proofstepwidestar[1]{\Finite A}{\FormulaRefAuto{FiniteSemigroupFinitenessAxiom}[ax]{1}}
 \proofstepwidestar[1]{n_A\in\mathbb N}{\FormulaRefAuto{FiniteCardNatural}[thm]{2}}
 \proofstepwidestar[1]{I\subseteq\mathbb N}{\FormulaRefAuto{PeanoNatSegSubsetN}[thm]{3}}
 \proofstepwidestar[5]{w\in W_A}{\rA}
 \proofstepwidestar[1,5]{\exists e\,(e\colon I\bij A\land w=w_{n_A}(m_{\star,e}))}{\FormulaRefAuto{FiniteSemigroupCanonicalTableCodeDef}[def]{1,5}}
 \proofstepwidestar[7]{e\colon I\bij A\land w=w_{n_A}(m_{\star,e})}{\rA}
 \proofstepwidestar[1,7]{m\colon I\times I\to I}{\FormulaRefAuto{FiniteSemigroupMultiplicationTableDef}[def]{1,3,\rAEa{7}}}
 \proofstepwidestar[9]{i\in I}{\rA}
 \proofstepwidestar[10]{j\in I}{\rA}
 \proofstepwidestar[9,10]{(i,j)\in I\times I}{\FormulaRefAuto{a\in A\dsep b\in B\vdash(a,b)\in A\times B}[thm]{9,10}}
 \proofstepwidestar[1,7,9,10]{m(i,j)\in I}{\FormulaRefAuto{F\colon A\to B\dsep x\in A\vdash F(x)\in B}[thm]{8,11}}
 \proofstepwidestar[1,7,9,10]{m(i,j)\in\mathbb N}{\FormulaRefAuto{A\subseteq B\dsep x\in A\vdash x\in B}[thm]{4,12}}
 \proofstepwidestar[1,7,9]{\forall j\in I\,(m(i,j)\in\mathbb N)}{\rUI{\rRI{10,13}}}
 \proofstepwidestar[1,7,9]{\exists!F\,(F\colon I\to\mathbb N\land\forall j\in I\,(F(j)=m(i,j)))}{\FormulaRefAuto{TypedFunctionDefinitionPrinciple}[thm]{14}}
 \proofstepwidestar[1,7,9]{\ell_i\colon I\to\mathbb N}{\FormulaRefAuto{FiniteSemigroupTableTupleDef}[def]{3,8,9,15}}
 \proofstepwidestar[1,7,9]{\ell_i\subseteq I\times\mathbb N}{\FormulaRefAuto{FunctionGraphSubsetProduct}[thm]{16}}
 \proofstepwidestar[1,7,9]{\ell_i\in\powerset(I\times\mathbb N)}{\FormulaRefAuto{x\in\powerset(A)\eqvdash x\subseteq A}[thm]{17}}
 \proofstepwidestar[1,7,9]{\ell_i\in U}{\FormulaRefAuto{B37TLexTupleDef}[def]{18,16}}
 \proofstepwidestar[1,7]{\forall i\in I\,(\ell_i\in U)}{\rUI{\rRI{9,19}}}
 \proofstepwidestar[1,7]{\exists!G\,(G\colon I\to U\land\forall i\in I\,(G(i)=\ell_i))}{\FormulaRefAuto{TypedFunctionDefinitionPrinciple}[thm]{20}}
 \proofstepwidestar[1,7]{w_{n_A}(m)\colon I\to U}{\FormulaRefAuto{FiniteSemigroupTableTupleDef}[def]{3,8,21}}
 \proofstepwidestar[1,7]{w\colon I\to U}{\rIE{\FormulaRefAuto{a=b\vdash b=a}[thm]{\rAEb{7}},22}}
 \proofstepwidestar[1,7]{w\subseteq I\times U}{\FormulaRefAuto{FunctionGraphSubsetProduct}[thm]{23}}
 \proofstepwidestar[1,7]{w\in\powerset(I\times U)}{\FormulaRefAuto{x\in\powerset(A)\eqvdash x\subseteq A}[thm]{24}}
 \proofstepwidestar[1,7]{w\in V_{n_A}}{\FormulaRefAuto{B37TLexTupleDef}[def]{25,23}}
 \proofstepwidestar[1,5]{w\in V_{n_A}}{\rEE{6,7,26}}
 \proofstepwidestar[1]{W_A\subseteq V_{n_A}}{\FormulaRefAuto{A\subseteq B\coloneq\forall x\,(x\in A\to x\in B)}[def]{\rUI{\rRI{5,27}}}}
 \proofstepwidestar[1]{n_A\in\mathbb N\land W_A\subseteq V_{n_A}}{\rAI{3,28}}
\end{tabproofwide}
Bei \(n_A=0\) sind beide Allquantoren über \(I=\varnothing\) leer.
Das Funktionsdefinitionsprinzip liefert dann die leere äußere Funktion;
keine positive Tupellänge wurde vorausgesetzt.

\Needspace{14\baselineskip}
\FormulaThmDeltaK[Genau ein kleinster Tafelkandidat]{%
 \FiniteSemigroupAlg{A}{\star}\vdash\exists!q\in W_A\,\forall w\in W_A\,(q\preceq_{n_A,U_{n_A}}w)
}{B37TLexCodeMinimum}{\DeltaRow{Mengen}{A\dsep n_A\dsep W_A\dsep q\dsep w}\DeltaRow{Funktionen}{\star}}
\begin{tabproofwide}
 \proofstepwidestar[1]{\FiniteSemigroupAlg{A}{\star}}{\rA}
 \proofstepwidestar[1]{n_A\in\mathbb N\land W_A\subseteq V_{n_A}}{\FormulaRefAuto{B37TLexCodeCandidatesTyped}[thm]{1}}
 \proofstepwidestar[]{K(\mathbb N,<)}{\FormulaRefAuto{B37TLexNaturalBase}[thm]}
 \proofstepwidestar[1]{K(U_{n_A},\prec_{n_A,\mathbb N})}{\FormulaRefAuto{B37TLexIterationBase}[thm]{\rAI{\rAEa{2},3}}}
 \proofstepwidestar[1]{n_A\in\mathbb N\land K(U_{n_A},\prec_{n_A,\mathbb N})}{\rAI{\rAEa{2},4}}
 \proofstepwidestar[1]{\Finite{W_A}\land W_A\ne\varnothing}{\FormulaRefAuto{B37TFiniteCodeCandidates}[thm]{1}}
 \proofstepwidestar[1]{\exists!q\in W_A\,\forall w\in W_A\,(q\preceq_{n_A,U_{n_A}}w)}{\FormulaRefAuto{B37TLexFiniteMinimum}[thm]{5,\rAEb{2},\rAEa{6},\rAEb{6}}}
\end{tabproofwide}
\Needspace{14\baselineskip}
\FormulaDefDeltaK[Minimaler Tafelcode]{%
 \begin{aligned}[t]
 &\FiniteSemigroupAlg{A}{\star}\vdash{}\\
 &\mu_A\coloneqq\min_{V_{n_A},\preceq_{n_A,U_{n_A}}}(W_A),\qquad
 \operatorname{mcan}(A,\star)\coloneqq(n_A,\mu_A)
 \end{aligned}
}{B37TMinimalCodeDef}{\DeltaRow{Mengen}{A\dsep n_A\dsep W_A\dsep\mu_A}\DeltaRow{Funktionen}{\star}}
\Needspace{14\baselineskip}
\FormulaThmDeltaK[Der gewählte minimale Code ist ein Kandidat]{%
 \FiniteSemigroupAlg{A}{\star}\vdash\mu_A\in W_A\land\forall w\in W_A\,(\mu_A\preceq_{n_A,U_{n_A}}w)
}{B37TLexMinimumMembership}{\DeltaRow{Mengen}{A\dsep\mu_A\dsep w}\DeltaRow{Funktionen}{\star}}
\begin{tabproofwide}
 \proofstepwidestar[1]{\FiniteSemigroupAlg{A}{\star}}{\rA}
 \proofstepwidestar[1]{\exists!q\in W_A\,\forall w\in W_A\,(q\preceq_{n_A,U_{n_A}}w)}{\FormulaRefAuto{B37TLexCodeMinimum}[thm]{1}}
 \proofstepwidestar[1]{\mu_A\in W_A\land\forall w\in W_A\,(\mu_A\preceq_{n_A,U_{n_A}}w)}{\FormulaRefAuto{B37TMinimalCodeDef}[def]{1,\FormulaRefAuto{Minimum}[def]{2,\FormulaRefAuto{IotaSymbolDef}[def]{2}}}}
\end{tabproofwide}

\Needspace{14\baselineskip}
\FormulaThmDeltaK[Isomorphie erhält den minimalen Tafelcode]{%
 \begin{aligned}[t]
 &\FiniteSemigroupAlg{A}{\star}\dsep\FiniteSemigroupAlg{B}{\diamond}\dsep
 \exists f\,\FiniteSemigroupIso{f}{A,\star}{B,\diamond}\vdash{}\\
 &\operatorname{mcan}(A,\star)=\operatorname{mcan}(B,\diamond)
 \end{aligned}
}{B37TMinimalCodeForward}{\DeltaRow{Mengen}{A\dsep B\dsep n_A\dsep n_B\dsep W_A\dsep W_B}\DeltaRow{Funktionen}{\star\dsep\diamond\dsep f}}
\begin{tabproofwide}
 \proofstepwidestar[1]{\FiniteSemigroupAlg{A}{\star}}{\rA}
 \proofstepwidestar[2]{\FiniteSemigroupAlg{B}{\diamond}}{\rA}
 \proofstepwidestar[3]{\exists f\,\FiniteSemigroupIso{f}{A,\star}{B,\diamond}}{\rA}
 \proofstepwidestar[1,2,3]{\operatorname{can}(A,\star)=\operatorname{can}(B,\diamond)}{\rRE{\rLREa{\FormulaRefAuto{FiniteSemigroupCanonicalTableCodeComplete}[thm]{1,2}},3}}
 \proofstepwidestar[1,2,3]{n_A=n_B\land W_A=W_B}{\FormulaRefAuto{(a,b)=(c,d)\eqvdash a=c\land b=d}[thm]{\FormulaRefAuto{FiniteSemigroupCanonicalTableCodeDef}[def]{1,2,4}}}
 \proofstepwidestar[1,2,3]{\mu_A=\mu_B}{\FormulaRefAuto{B37TMinimalCodeDef}[def]{1,2,\rAEa{5},\rAEb{5}}}
 \proofstepwidestar[1,2,3]{\operatorname{mcan}(A,\star)=\operatorname{mcan}(B,\diamond)}{\FormulaRefAuto{B37TMinimalCodeDef}[def]{1,2,\rAEa{5},6}}
\end{tabproofwide}

\Needspace{14\baselineskip}
\FormulaThmDeltaK[Gleiche minimale Codes liefern einen Isomorphismus]{%
 \begin{aligned}[t]
 &\FiniteSemigroupAlg{A}{\star}\dsep\FiniteSemigroupAlg{B}{\diamond}\dsep
 \operatorname{mcan}(A,\star)=\operatorname{mcan}(B,\diamond)\vdash{}\\
 &\exists f\,\FiniteSemigroupIso{f}{A,\star}{B,\diamond}
 \end{aligned}
}{B37TMinimalCodeBackward}{\DeltaRow{Mengen}{A\dsep B\dsep n_A\dsep n_B\dsep W_A\dsep W_B}\DeltaRow{Funktionen}{\star\dsep\diamond\dsep e_A\dsep e_B\dsep f}}
\begin{tabproofwide}
 \proofstepwidestar[1]{\FiniteSemigroupAlg{A}{\star}}{\rA}
 \proofstepwidestar[2]{\FiniteSemigroupAlg{B}{\diamond}}{\rA}
 \proofstepwidestar[3]{\operatorname{mcan}(A,\star)=\operatorname{mcan}(B,\diamond)}{\rA}
 \proofstepwidestar[1,2,3]{n_A=n_B\land\mu_A=\mu_B}{\FormulaRefAuto{(a,b)=(c,d)\eqvdash a=c\land b=d}[thm]{\FormulaRefAuto{B37TMinimalCodeDef}[def]{1,2,3}}}
 \proofstepwidestar[1]{\mu_A\in W_A}{\rAEa{\FormulaRefAuto{B37TLexMinimumMembership}[thm]{1}}}
 \proofstepwidestar[2]{\mu_B\in W_B}{\rAEa{\FormulaRefAuto{B37TLexMinimumMembership}[thm]{2}}}
 \proofstepwidestar[1]{n_A\in\mathbb N}{\FormulaRefAuto{FiniteCardNatural}[thm]{\FormulaRefAuto{FiniteSemigroupFinitenessAxiom}[ax]{1}}}
 \proofstepwidestar[1]{\exists e_A\,(e_A\colon I_{n_A}\bij A\land\mu_A=w_{n_A}(m_{\star,e_A}))}{\FormulaRefAuto{FiniteSemigroupCanonicalTableCodeDef}[def]{1,5}}
 \proofstepwidestar[9]{e_A\colon I_{n_A}\bij A\land\mu_A=w_{n_A}(m_{\star,e_A})}{\rA}
 \proofstepwidestar[1,2,3]{\exists e_B\,(e_B\colon I_{n_A}\bij B\land\mu_A=w_{n_A}(m_{\diamond,e_B}))}{\rIE{\FormulaRefAuto{a=b\vdash b=a}[thm]{\rAEa{4}},\FormulaRefAuto{a=b\vdash b=a}[thm]{\rAEb{4}},\FormulaRefAuto{FiniteSemigroupCanonicalTableCodeDef}[def]{2,6}}}
 \proofstepwidestar[11]{e_B\colon I_{n_A}\bij B\land\mu_A=w_{n_A}(m_{\diamond,e_B})}{\rA}
 \proofstepwidestar[9,11]{w_{n_A}(m_{\star,e_A})=w_{n_A}(m_{\diamond,e_B})}{\rIE{\rAEb{9},\rAEb{11}}}
 \proofstepwidestar[1,9]{m_{\star,e_A}\colon I_{n_A}\times I_{n_A}\to I_{n_A}}{\FormulaRefAuto{FiniteSemigroupMultiplicationTableDef}[def]{1,7,\rAEa{9}}}
 \proofstepwidestar[1,2,11]{m_{\diamond,e_B}\colon I_{n_A}\times I_{n_A}\to I_{n_A}}{\FormulaRefAuto{FiniteSemigroupMultiplicationTableDef}[def]{2,7,\rAEa{11}}}
 \proofstepwidestar[1,2,9,11]{m_{\star,e_A}=m_{\diamond,e_B}}{\FormulaRefAuto{FiniteSemigroupTableTupleInjective}[thm]{7,13,14,12}}
 \proofstepwidestar[1,2,9,11]{\FiniteSemigroupIso{e_B\circ e_A^{-1}}{A,\star}{B,\diamond}}{\FormulaRefAuto{B37TableCodeCommonTableIsomorphism}[thm]{1,2,7,\rAEa{9},\rAEa{11},15}}
 \proofstepwidestar[1,2,9,11]{\exists f\,\FiniteSemigroupIso{f}{A,\star}{B,\diamond}}{\rEI{16}}
 \proofstepwidestar[1,2,3,9]{\exists f\,\FiniteSemigroupIso{f}{A,\star}{B,\diamond}}{\rEE{10,11,17}}
 \proofstepwidestar[1,2,3]{\exists f\,\FiniteSemigroupIso{f}{A,\star}{B,\diamond}}{\rEE{8,9,18}}
\end{tabproofwide}

\Needspace{14\baselineskip}
\FormulaThmDeltaK[Der minimale Tafelcode ist vollständig]{%
 \begin{aligned}[t]
 &\FiniteSemigroupAlg{A}{\star}\dsep\FiniteSemigroupAlg{B}{\diamond}\vdash{}\\
 &\bigl(\exists f\,\FiniteSemigroupIso{f}{A,\star}{B,\diamond}\bigr)
 \leftrightarrow\operatorname{mcan}(A,\star)=\operatorname{mcan}(B,\diamond)
 \end{aligned}
}{B37TMinimalCodeCriterion}{\DeltaRow{Mengen}{A\dsep B}\DeltaRow{Funktionen}{\star\dsep\diamond\dsep f}}
\begin{tabproofwide}
 \proofstepwidestar[1]{\FiniteSemigroupAlg{A}{\star}}{\rA}
 \proofstepwidestar[2]{\FiniteSemigroupAlg{B}{\diamond}}{\rA}
 \proofstepwidestar[3]{\exists f\,\FiniteSemigroupIso{f}{A,\star}{B,\diamond}}{\rA}
 \proofstepwidestar[1,2,3]{\operatorname{mcan}(A,\star)=\operatorname{mcan}(B,\diamond)}{\FormulaRefAuto{B37TMinimalCodeForward}[thm]{1,2,3}}
 \proofstepwidestar[1,2]{(\exists f\,\FiniteSemigroupIso{f}{A,\star}{B,\diamond})\to\operatorname{mcan}(A,\star)=\operatorname{mcan}(B,\diamond)}{\rRI{3,4}}
 \proofstepwidestar[6]{\operatorname{mcan}(A,\star)=\operatorname{mcan}(B,\diamond)}{\rA}
 \proofstepwidestar[1,2,6]{\exists f\,\FiniteSemigroupIso{f}{A,\star}{B,\diamond}}{\FormulaRefAuto{B37TMinimalCodeBackward}[thm]{1,2,6}}
 \proofstepwidestar[1,2]{\operatorname{mcan}(A,\star)=\operatorname{mcan}(B,\diamond)\to\exists f\,\FiniteSemigroupIso{f}{A,\star}{B,\diamond}}{\rRI{6,7}}
 \proofstepwidestar[1,2]{(\exists f\,\FiniteSemigroupIso{f}{A,\star}{B,\diamond})\leftrightarrow\operatorname{mcan}(A,\star)=\operatorname{mcan}(B,\diamond)}{\rLRI{5,8}}
\end{tabproofwide}
Die erste Komponente \(n_A\) gehört ausdrücklich zum minimalen Code.
Der Nachweis verwendet keine Nichtleerheit von \(A\) oder \(B\);
im Fall \(n_A=0\) ist der einzige Kandidat das leere Tupel.
